% La santé reproductive — SVTEEHB Tle TI — Cours signé OnBuch+
% Programme officiel MINESEC. Compiler avec : tectonic N07-sante-reproductive.tex
\newcommand{\DOCMATIERE}{SVTEEHB}
\newcommand{\DOCNIVEAU}{Tle TI}
\newcommand{\DOCMODULE}{Module II : L'Éducation à la Santé}
\newcommand{\DOCLECON}{N07}
\newcommand{\DOCTITRE}{La santé reproductive}
% OnBuch+ — préambule « cours signés » ENRICHI (compilé avec tectonic / XeLaTeX)
% Système visuel « Pop » d'OnBuch+ : fond crème (jamais blanc), bordures encre
% épaisses, ombres « dures » décalées, titres Archivo Black en capitales.
% Version MATHÉMATIQUES : graphes (pgfplots/tikz), boîtes pédagogiques
% (définition, propriété, méthode, attention, le savais-tu, exercice résolu,
% exercice, TP, bilan), page de garde et signature OnBuch+ ancrée en bas.
%
% Macros attendues dans main.tex AVANT \input{preamble} :
%   \DOCMATIERE  \DOCNIVEAU  \DOCMODULE  \DOCLECON (numéro)  \DOCTITRE
\documentclass[11pt]{article}

\usepackage{fontspec}
\usepackage[french]{babel}
\usepackage[a4paper,margin=2cm,top=3.4cm,bottom=2.8cm,headheight=1.4cm,footskip=1.3cm]{geometry}
\usepackage{amsmath,amssymb,amsfonts}
\usepackage{mathtools} % \xmapsto, \coloneqq... souvent utilisés par les modèles
\usepackage{cancel} % simplifications barrées (bilans de forces)
% \abs{z} (module, valeur absolue) et \norm{u} (norme), utilisés par les modèles
\providecommand{\abs}{}\let\abs\relax\DeclarePairedDelimiter{\abs}{\lvert}{\rvert}
\providecommand{\norm}{}\let\norm\relax\DeclarePairedDelimiter{\norm}{\lVert}{\rVert}
\usepackage[table]{xcolor} % [table] : fournit aussi \rowcolors (alternance de lignes)
\usepackage{pagecolor}
\usepackage{tikz}
\usetikzlibrary{arrows.meta,positioning,calc,shapes.geometric,shapes.misc,shapes.arrows,decorations.pathmorphing,decorations.pathreplacing,decorations.markings,patterns,shadows,mindmap,trees,fit,backgrounds,matrix,intersections,angles,quotes}
\usepackage{pgfplots}
\usepackage[version=4]{mhchem} % équations nucléaires \ce{^{235}_{92}U}
\usepackage[european,siunitx]{circuitikz} % schémas électriques
\pgfplotsset{compat=1.18}
\usepgfplotslibrary{fillbetween,groupplots} % aires entre courbes (calcul intégral)
\usepackage{enumitem}
\usepackage{fancyhdr}
% [most] charge toutes les bibliothèques tcolorbox (listings, minted, raster,
% documentation...) et gonfle inutilement la mémoire TeX sur un document long
% (~300 boîtes) : on ne charge que ce qui est réellement utilisé.
\usepackage[skins,breakable]{tcolorbox}
\usepackage{graphicx}
\usepackage{colortbl}
\usepackage{booktabs}
\usepackage{tabularx}
\usepackage{multirow}
\usepackage{diagbox} % case d'en-tête diagonale des tableaux à double entrée
% Colonnes extensibles centrée / gauche / droite (C, L, R), souvent utilisées par les modèles
\newcolumntype{C}{>{\centering\arraybackslash}X}
\newcolumntype{L}{>{\raggedright\arraybackslash}X}
\newcolumntype{R}{>{\raggedleft\arraybackslash}X}
\usepackage{array}
\usepackage{multicol}
\usepackage{siunitx}
% parse-numbers=false : accepte aussi \SI{2,5 \times 10^{-3}}{...}
\sisetup{locale=FR,output-decimal-marker={,},group-minimum-digits=4,parse-numbers=false}
\DeclareSIUnit\ohm{\ensuremath{\symup{\Omega}}} % U+2126 absent de la police
\DeclareSIUnit\division{div} % réglages d'oscilloscope (ms/div, V/div)
\DeclareSIUnit\tr{tr} % tours (vitesse de rotation, ex. \SI{1500}{\tr\per\minute})
\DeclareSIUnit\molar{M} % concentration molaire (ex. \SI{3}{\molar}, \SI{1}{\micro\molar})
\DeclareSIUnit\an{an} % année (le modèle confond parfois avec \year, un compteur TeX) — ex. \SI{5}{\kilo\meter\per\an}
\DeclareSIUnit\FCFA{FCFA} % franc CFA (ex. \SI{500}{\FCFA\per\kilo\gram})
\DeclareSIUnit\degreeBrix{\textdegree Brix} % teneur en sucre d'un sirop/jus (réfractométrie)
\DeclareSIUnit\month{mois} % le modèle utilise parfois \month, un compteur TeX, comme unité de durée
\DeclareSIUnit\microliter{\micro\liter} % le modèle écrit parfois \microliter en un seul token
\DeclareSIUnit\million{millions} % dénombrement (ex. \SI{15}{\million\per\mL} spermatozoïdes)
\usepackage{qrcode}
% \degree : commande fréquemment utilisée par les modèles pour un angle ou une
% température en dehors de \SI{}{} (non fournie par les paquets chargés).
\providecommand{\degree}{^{\circ}}
% \qed (fin de démonstration), utilisé par les modèles sans amsthm
\providecommand{\qed}{\hfill$\square$}
% \barre : double barre des valeurs interdites dans un tableau de variations (tkz-tab)
\providecommand{\barre}{\ensuremath{\|}}
% environnement proof (amsthm non chargé) utilisé par les modèles
\ifdefined\proof\else\newenvironment{proof}[1][Démonstration]{\par\noindent\textit{#1.}\ }{\qed\par}\fi
\usepackage{lastpage}
\usepackage{hyperref}
\hypersetup{hidelinks}

% ============================================================
% Palette « Pop » OnBuch+ — mêmes hex que la classe P (plus_ui.dart)
% ============================================================
\definecolor{poporange}{HTML}{F59321}
\definecolor{popdark}{HTML}{E07A0C}
\definecolor{popcream}{HTML}{FFF6E8}
\definecolor{popink}{HTML}{1C1714}
\definecolor{poppurple}{HTML}{7A5AE0}
\definecolor{popgreen}{HTML}{2E9E6A}
\definecolor{poppink}{HTML}{E0457B}
\definecolor{popblue}{HTML}{2F7DE1}
\definecolor{popgold}{HTML}{C98A12}
\definecolor{popmuted}{HTML}{8A7F76}
\definecolor{popline}{HTML}{EADFD2}
\definecolor{popgray}{HTML}{8A7F76}
\colorlet{popgrey}{popgray}
% Alias tolérants (le modèle nomme parfois les couleurs différemment)
\colorlet{popwhite}{white}
\colorlet{popblack}{popink}
\colorlet{popred}{poppink}
\colorlet{popyellow}{popgold}
% Teintes claires (fonds de tableaux/encadrés) — le modèle nomme parfois
% les couleurs avec un suffixe L (ex. popblueL) sans les avoir définies.
\colorlet{poporangeL}{poporange!20}
\colorlet{popdarkL}{popdark!20}
\colorlet{poppurpleL}{poppurple!20}
\colorlet{popgreenL}{popgreen!20}
\colorlet{poppinkL}{poppink!20}
\colorlet{popblueL}{popblue!20}
\colorlet{popgoldL}{popgold!20}
\colorlet{popredL}{popred!20}
\colorlet{popgrayL}{popgray!20}
% Teintes claires pour fonds de tableaux / zones
\colorlet{popblueL}{popblue!10!white}
\colorlet{poporangeL}{poporange!14!white}
\colorlet{popgreenL}{popgreen!12!white}
\colorlet{poppurpleL}{poppurple!12!white}
\colorlet{poppinkL}{poppink!12!white}
\colorlet{popgoldL}{popgold!14!white}

\pagecolor{popcream}

% ============================================================
% Typographie
% ============================================================
\defaultfontfeatures{Ligatures=TeX}
\setmainfont{PlusJakartaSans-Medium.ttf}[Path=../../../fonts/,BoldFont=PlusJakartaSans-Bold.ttf]
% Maths en sans-serif (Fira Math) avec chiffres et lettres droites de Plus
% Jakarta, pour que les formules se fondent dans le texte.
\usepackage[math-style=ISO,bold-style=ISO]{unicode-math}
\setmathfont{FiraMath-Regular.otf}[Path=../../../fonts/]
\setmathfont{PlusJakartaSans-Medium.ttf}[Path=../../../fonts/,range={up/num,up/{Latin,latin},\mathup}]
\usepackage{icomma} % virgule décimale sans espace en mode math
% Caractères mathématiques Unicode absents des polices de texte (Plus Jakarta,
% Archivo Black) : rendus par leur équivalent mathématique, en texte comme en
% formule — sinon ils s'affichent en carré vide (ex. « Arithmétique dans ℤ »).
\usepackage{newunicodechar}
\newunicodechar{ℝ}{\ensuremath{\mathbb{R}}}
\newunicodechar{ℤ}{\ensuremath{\mathbb{Z}}}
\newunicodechar{ℂ}{\ensuremath{\mathbb{C}}}
\newunicodechar{ℕ}{\ensuremath{\mathbb{N}}}
\newunicodechar{ℚ}{\ensuremath{\mathbb{Q}}}
\newunicodechar{𝔻}{\ensuremath{\mathbb{D}}}
\newunicodechar{α}{\ensuremath{\alpha}}
\newunicodechar{β}{\ensuremath{\beta}}
\newunicodechar{γ}{\ensuremath{\gamma}}
\newunicodechar{δ}{\ensuremath{\delta}}
\newunicodechar{ε}{\ensuremath{\varepsilon}}
\newunicodechar{θ}{\ensuremath{\theta}}
\newunicodechar{λ}{\ensuremath{\lambda}}
\newunicodechar{μ}{\ensuremath{\mu}}
\newunicodechar{σ}{\ensuremath{\sigma}}
\newunicodechar{τ}{\ensuremath{\tau}}
\newunicodechar{φ}{\ensuremath{\varphi}}
\newunicodechar{ψ}{\ensuremath{\psi}}
\newunicodechar{ω}{\ensuremath{\omega}}
\newunicodechar{Σ}{\ensuremath{\Sigma}}
% majuscules produites par \MakeUppercase dans les titres (α-aminés -> Α)
\newunicodechar{Α}{\ensuremath{\alpha}}
\newunicodechar{Β}{\ensuremath{\beta}}
\newunicodechar{Γ}{\ensuremath{\Gamma}}
\newunicodechar{Δ}{\ensuremath{\Delta}}
\newunicodechar{Ε}{\ensuremath{\varepsilon}}
\newunicodechar{Θ}{\ensuremath{\Theta}}
\newunicodechar{Λ}{\ensuremath{\Lambda}}
\newunicodechar{Μ}{\ensuremath{\mu}}
\newunicodechar{Τ}{\ensuremath{\tau}}
\newunicodechar{Φ}{\ensuremath{\Phi}}
\newunicodechar{Ψ}{\ensuremath{\Psi}}
\newunicodechar{Ω}{\ensuremath{\Omega}}
\newunicodechar{↦}{\ensuremath{\mapsto}}
\newunicodechar{∈}{\ensuremath{\in}}
\newunicodechar{∉}{\ensuremath{\notin}}
\newunicodechar{∧}{\ensuremath{\wedge}}
\newunicodechar{⇔}{\ensuremath{\Leftrightarrow}}
\newunicodechar{⟺}{\ensuremath{\Longleftrightarrow}}
\newunicodechar{⇒}{\ensuremath{\Rightarrow}}
\newunicodechar{⟹}{\ensuremath{\Longrightarrow}}
\newunicodechar{∘}{\ensuremath{\circ}}
\newunicodechar{∪}{\ensuremath{\cup}}
\newunicodechar{∩}{\ensuremath{\cap}}
\newunicodechar{≡}{\ensuremath{\equiv}}
\newunicodechar{⊂}{\ensuremath{\subset}}
\newunicodechar{⊥}{\ensuremath{\perp}}
\newunicodechar{∃}{\ensuremath{\exists}}
\newunicodechar{∀}{\ensuremath{\forall}}
\newunicodechar{∛}{\ensuremath{\sqrt[3]{\,}}}
\newunicodechar{∜}{\ensuremath{\sqrt[4]{\,}}}
\newunicodechar{⃗}{\ensuremath{{}^{\to}}}
\newunicodechar{ⁿ}{\ifmmode^{n}\else\textsuperscript{\ensuremath{n}}\fi}
\newunicodechar{ˣ}{\ifmmode^{x}\else\textsuperscript{\ensuremath{x}}\fi}
\newunicodechar{ᵇ}{\ifmmode^{b}\else\textsuperscript{\ensuremath{b}}\fi}
\newunicodechar{ʳ}{\ifmmode^{r}\else\textsuperscript{\ensuremath{r}}\fi}
\newunicodechar{ᵘ}{\ifmmode^{u}\else\textsuperscript{\ensuremath{u}}\fi}
\newunicodechar{ʸ}{\ifmmode^{y}\else\textsuperscript{\ensuremath{y}}\fi}
\newunicodechar{ᵃ}{\ifmmode^{a}\else\textsuperscript{\ensuremath{a}}\fi}
\newunicodechar{ᵉ}{\ifmmode^{e}\else\textsuperscript{\ensuremath{e}}\fi}
\newunicodechar{ᵏ}{\ifmmode^{k}\else\textsuperscript{\ensuremath{k}}\fi}
\newunicodechar{ᵖ}{\ifmmode^{p}\else\textsuperscript{\ensuremath{p}}\fi}
\newunicodechar{ᴺ}{\ifmmode^{N}\else\textsuperscript{\ensuremath{N}}\fi}
\newunicodechar{ᶜ}{\ifmmode^{c}\else\textsuperscript{\ensuremath{c}}\fi}
\newunicodechar{ʰ}{\ifmmode^{h}\else\textsuperscript{\ensuremath{h}}\fi}
\newunicodechar{ᵠ}{\ifmmode^{q}\else\textsuperscript{\ensuremath{q}}\fi}
\newunicodechar{ₙ}{\ifmmode_{n}\else\textsubscript{\ensuremath{n}}\fi}
\newunicodechar{ₐ}{\ifmmode_{a}\else\textsubscript{\ensuremath{a}}\fi}
\newunicodechar{ᵢ}{\ifmmode_{i}\else\textsubscript{\ensuremath{i}}\fi}
\newunicodechar{ᵣ}{\ifmmode_{r}\else\textsubscript{\ensuremath{r}}\fi}
\newunicodechar{ₖ}{\ifmmode_{k}\else\textsubscript{\ensuremath{k}}\fi}
\newunicodechar{ᵦ}{\ifmmode_{\beta}\else\textsubscript{\ensuremath{\beta}}\fi}
\newunicodechar{ₓ}{\ifmmode_{x}\else\textsubscript{\ensuremath{x}}\fi}
\newfontfamily\popdisplay{ArchivoBlack-Regular.ttf}[Path=../../../fonts/]
\newfontfamily\poplogo{SpaceGrotesk-Bold.ttf}[Path=../../../fonts/]
\newfontfamily\popsemi{PlusJakartaSans-SemiBold.ttf}[Path=../../../fonts/]
\newfontfamily\popxbold{PlusJakartaSans-ExtraBold.ttf}[Path=../../../fonts/]
\newfontfamily\popxboldls{PlusJakartaSans-ExtraBold.ttf}[Path=../../../fonts/,LetterSpace=3.0]

\setlist[enumerate]{leftmargin=1.7em,itemsep=5pt,topsep=4pt}
\setlist[itemize]{leftmargin=1.7em,itemsep=4pt,topsep=4pt,label={\color{poporange}\textbullet}}
\setlength{\parindent}{0pt}
\setlength{\parskip}{5pt}
\renewcommand{\arraystretch}{1.25}

% pgfplots : style Pop par défaut
\pgfplotsset{
  every axis/.append style={axis line style={popink,line width=1pt},tick style={popink},
    grid style={popline,line width=0.5pt},label style={font=\small},tick label style={font=\footnotesize}},
  popcurve/.style={poporange,line width=1.8pt,no markers,smooth},
}
% Même style disponible dans un \draw TikZ simple (hors pgfplots)
\tikzset{popcurve/.style={poporange,line width=1.8pt}}
\tikzset{popgrid/.style={popline,very thin}}
\colorlet{popcurve}{poporange} % le nom du style est parfois utilisé comme couleur

% ============================================================
% Pill / squiggle
% ============================================================
\newcommand{\poppill}[3]{%
  \tikz[baseline=-0.55ex]\node[draw=popink,line width=1.2pt,rounded corners=2.6pt,%
    fill=#2,inner xsep=6.5pt,inner ysep=3pt]{%
    \popxboldls\fontsize{7}{7}\selectfont\color{#3}\MakeUppercase{#1}};%
}
\newcommand{\popsquiggle}[1]{%
  \tikz[baseline=-0.4ex]{
    \draw[#1,line width=2.6pt,line cap=round]
      plot [smooth] coordinates {(0,0.09) (0.34,0.01) (0.68,0.21) (1.02,0.01) (1.36,0.10)};
  }%
}
% Mot-clé surligné (marqueur orange sous le texte)
\newcommand{\cle}[1]{{\popxbold\color{popdark}#1}}

% ============================================================
% En-tête / pied
% ============================================================
\pagestyle{fancy}
\fancyhf{}
\renewcommand{\headrulewidth}{0pt}
\renewcommand{\footrulewidth}{0pt}
\lhead{\raisebox{-0.15ex}{\poplogo\fontsize{13}{13}\selectfont\color{popink}OnBuch\color{poporange}+}}
\rhead{\poppill{\DOCMATIERE\ \textbullet\ \DOCNIVEAU\ \textbullet\ Leçon \DOCLECON}{white}{popink}}
\lfoot{\popxbold\fontsize{7.6}{7.6}\selectfont\color{popmuted}\MakeUppercase{Cours signé OnBuch+}\ \textbullet\ {\popsemi\DOCTITRE}}
\rfoot{\tikz[baseline=-0.6ex]\node[rounded corners=8.5pt,fill=popink,minimum height=17pt,minimum width=17pt,inner xsep=5pt,inner sep=0]{\popxbold\fontsize{8}{8}\selectfont\color{popcream}\thepage\,/\,\pageref*{LastPage}};}

% ============================================================
% Titres
% ============================================================
\newcommand{\chaptitre}[2]{%
  \begin{tcolorbox}[enhanced,colback=poporange,colframe=popink,boxrule=2.6pt,arc=11pt,
      drop shadow=popink,left=15pt,right=15pt,top=12pt,bottom=11pt,before skip=3pt,after skip=18pt]
    \poppill{#2}{popink}{popcream}\par\vspace{9pt}
    {\raggedright\hyphenpenalty=10000\exhyphenpenalty=10000\popdisplay\fontsize{22}{24}\selectfont\color{popcream}\MakeUppercase{#1}\par}
  \end{tcolorbox}
}
% Section numérotée : pastille orange avec le numéro + titre Archivo
\newcounter{popsec}
\newcommand{\coursec}[1]{%
  \stepcounter{popsec}\setcounter{popsub}{0}%
  \par\vspace{16pt}\noindent
  \begin{minipage}{\linewidth}
  \tikz[baseline=-0.7ex]\node[draw=popink,line width=1.6pt,fill=poporange,rounded corners=4pt,minimum size=22pt,inner sep=0]{\popdisplay\fontsize{12}{12}\selectfont\color{popcream}\thepopsec};%
  \hspace{8pt}{\popdisplay\fontsize{15}{17}\selectfont\color{popink}\MakeUppercase{#1}}\par
  \vspace{4pt}\noindent{\color{poporange}\rule{36pt}{3.6pt}}
  \end{minipage}\par\nopagebreak\vspace{8pt}
}
\newcounter{popsub}
\newcommand{\courssub}[1]{%
  \stepcounter{popsub}%
  \par\vspace{9pt}\noindent{\popxbold\fontsize{11.5}{13}\selectfont\color{popink}{\color{poporange}\thepopsec.\thepopsub}\hspace{6pt}#1}\par\nopagebreak\vspace{4pt}
}

% ============================================================
% Boîtes pédagogiques — même recette : fond blanc, bordure épaisse colorée,
% ombre dure encre, badge plein. Titre optionnel [#1].
% ============================================================
\tcbset{popbase/.style={breakable,enhanced,colback=white,boxrule=2.3pt,arc=9pt,
  shadow={3pt}{-3pt}{0pt}{popink},left=14pt,right=14pt,top=11pt,bottom=11pt,before skip=11pt,after skip=15pt,
  fontupper=\popsemi\fontsize{10.6}{14.5}\selectfont\color{popink},
  fontlower=\popsemi\fontsize{10.6}{14.5}\selectfont\color{popink}}}
% Coche dessinée (le glyphe ✓ n'existe pas dans les polices du document)
\newcommand{\popcheck}[1]{\tikz[baseline=-0.5ex]\draw[#1,line width=1.6pt,line cap=round,line join=round](-0.13,0.0)--(-0.04,-0.09)--(0.13,0.1);}
\providecommand{\coche}{\popcheck{popgreen}}
\providecommand{\resultat}[1]{\par\smallskip\noindent\fcolorbox{popgreen}{popgreenL}{\parbox{\dimexpr\linewidth-2\fboxsep-2\fboxrule\relax}{\textbf{Résultat.} #1}}\par\smallskip}
\newcommand{\popboxhead}[3]{\poppill{#1}{#2}{white}\ifx\relax#3\relax\else\hspace{6pt}{\popxbold\fontsize{10.6}{12}\selectfont\color{popink}#3}\fi\par\vspace{7pt}}

\newtcolorbox{aretenir}[1][]{popbase,colframe=popgreen,before upper={\popboxhead{À retenir}{popgreen}{#1}}}
\newtcolorbox{exemplebox}[1][]{popbase,colframe=poporange,before upper={\popboxhead{Exemple}{poporange}{#1}}}
\newtcolorbox{definition}[1][]{popbase,colframe=popblue,before upper={\popboxhead{Définition}{popblue}{#1}}}
\newtcolorbox{propriete}[1][]{popbase,colframe=poppurple,before upper={\popboxhead{Propriété}{poppurple}{#1}}}
\newtcolorbox{methode}[1][]{popbase,colframe=popgold,before upper={\popboxhead{Méthode}{popgold}{#1}}}
\newtcolorbox{attention}[1][]{popbase,colframe=poppink,before upper={\popboxhead{Attention}{poppink}{#1}}}
\newtcolorbox{experience}[1][]{popbase,colframe=popblue,colback=popblueL,before upper={\popboxhead{Expérience}{popblue}{#1}}}
% « Le savais-tu ? » : carte violette pleine (contexte, culture, Cameroun)
\newtcolorbox{savaistu}[1][]{popbase,colback=poppurple,colframe=popink,
  fontupper=\popsemi\fontsize{10.4}{14.2}\selectfont\color{white},
  before upper={\poppill{Le savais-tu\,?}{white}{poppurple}\ifx\relax#1\relax\else\hspace{6pt}{\popxbold\color{white}#1}\fi\par\vspace{7pt}}}
% Exercice résolu : énoncé en haut, \tcblower puis la solution sur fond crème
\newcounter{popexo}\newcounter{popexr}
\newtcolorbox{exoresolu}[1][]{popbase,colframe=poporange,colbacklower=popcream,
  segmentation style={popink,line width=1.2pt,dashed},
  before upper={\stepcounter{popexr}\popboxhead{Exercice résolu \thepopexr}{poporange}{#1}},
  before lower={\poppill{Solution}{popink}{popcream}\par\vspace{6pt}}}
% Exercice (non résolu en place) — la correction est dans la partie Corrigés
\newtcolorbox{exercice}[1][]{popbase,colframe=popink,
  before upper={\stepcounter{popexo}\popboxhead{Exercice \thepopexo}{popink}{#1}}}
% Corrigé
\newtcolorbox{corrige}[1][]{popbase,colframe=popgreen,colback=popgreenL,
  before upper={\popboxhead{Corrigé}{popgreen}{#1}}}
% Objectifs / prérequis / bilan
\newtcolorbox{objectifs}{popbase,colframe=popink,colback=poporangeL,
  before upper={\popboxhead{Objectifs de la leçon}{popink}{}}}
\newtcolorbox{prerequis}{popbase,colframe=popink,colback=popblueL,
  before upper={\popboxhead{Prérequis}{popblue}{}}}
\newtcolorbox{bilan}{popbase,colframe=popink,colback=white,boxrule=2.8pt,
  before upper={\popboxhead{Fiche bilan}{poporange}{L'essentiel à connaître pour l'examen}}}
% Figure encadrée avec légende
\newtcolorbox{popfigure}[1][]{enhanced,breakable=false,colback=white,colframe=popline,boxrule=1.4pt,arc=8pt,
  left=10pt,right=10pt,top=10pt,bottom=8pt,before skip=10pt,after skip=12pt,halign=center,
  lower separated=false,halign lower=center,
  fontlower=\popsemi\fontsize{9}{11.5}\selectfont\color{popmuted},
  #1}
\newcommand{\legende}[1]{\par\vspace{6pt}{\popsemi\fontsize{9}{11.5}\selectfont\color{popmuted}#1}}

% ============================================================
% Signature OnBuch+
% ============================================================
\newcommand{\poplogoinline}[1]{{\poplogo\fontsize{#1}{#1}\selectfont\color{popink}OnBuch\color{poporange}+}}
\newcommand{\signatureonbuch}{%
  \begin{tcolorbox}[enhanced,colback=popink,colframe=popink,boxrule=2.2pt,arc=10pt,
      drop shadow=poporange,left=14pt,right=14pt,top=10pt,bottom=10pt,before skip=0pt,after skip=0pt]
    \tikz[baseline=-0.7ex]\node[circle,fill=popgreen,draw=popcream,line width=1.3pt,minimum size=22pt,inner sep=0]{\popcheck{white}};%
    \hspace{9pt}\begin{minipage}[c]{0.62\linewidth}
      {\popxbold\fontsize{12}{13}\selectfont\color{popcream}Signé\ \textbullet\ {\poplogo OnBuch\color{poporange}+}}\par\vspace{2pt}
      {\raggedright\popsemi\fontsize{8}{10.5}\selectfont\color{popline}\MakeUppercase{Équipe pédagogique OnBuch+}\\\MakeUppercase{Conforme au programme officiel MINESEC}\par}
    \end{minipage}\hfill
    {\poplogo\fontsize{22}{22}\selectfont\color{popcream}OnBuch\color{poporange}+}
  \end{tcolorbox}%
}

% ============================================================
% Page de garde
% #1 titre · #2 matière · #3 niveau · #4 module · #5 n° leçon · #6 durée ·
% #7 description / objectifs (texte ou itemize)
% ============================================================
\newcommand{\pagegarde}[7]{%
  \thispagestyle{empty}
  \newgeometry{a4paper,margin=1.7cm,top=1.6cm,bottom=1.4cm}%
  \begin{tikzpicture}[remember picture,overlay]
    \fill[poporange!18] ([xshift=-3.2cm,yshift=-3.6cm]current page.north east) circle (4.6cm);
    \fill[poppurple!14] ([xshift=2.4cm,yshift=7.6cm]current page.south west) circle (3.8cm);
  \end{tikzpicture}%
  {\poplogo\fontsize{30}{30}\selectfont\color{popink}OnBuch\color{poporange}+}\hfill
  \raisebox{0.6ex}{\poppill{Cours signé \textbullet\ Premium}{popink}{popcream}}\par
  \vspace{1pt}\popsquiggle{poppurple}\par
  \vspace{8pt}
  \begin{tcolorbox}[enhanced,colback=poporange,colframe=popink,boxrule=2.8pt,arc=13pt,
      drop shadow=popink,left=18pt,right=18pt,top=12pt,bottom=13pt,after skip=10pt]
    \poppill{#2 \textbullet\ #3}{popink}{popcream}\hspace{5pt}\poppill{#4}{white}{popink}\par\vspace{8pt}
    {\popdisplay\fontsize{14}{14}\selectfont\color{popink}LEÇON #5}\par\vspace{4pt}
    {\raggedright\hyphenpenalty=10000\exhyphenpenalty=10000\popdisplay\fontsize{25}{27}\selectfont\color{popcream}\MakeUppercase{#1}\par}
    \vspace{8pt}
    {\popxbold\fontsize{9.5}{9.5}\selectfont\color{popink}\MakeUppercase{Durée conseillée : #6}}
  \end{tcolorbox}
  \begin{tcolorbox}[enhanced,colback=white,colframe=popink,boxrule=2.2pt,arc=10pt,
      drop shadow=popink,left=16pt,right=16pt,top=10pt,bottom=10pt,after skip=8pt]
    \poppill{Dans cette leçon}{poppurple}{white}\par\vspace{5pt}
    \popsemi\fontsize{9.3}{12.6}\selectfont\color{popink}#7
  \end{tcolorbox}
  \noindent\begin{minipage}[t]{0.487\linewidth}
    \begin{tcolorbox}[enhanced,colback=popgreen,colframe=popink,boxrule=2.2pt,arc=10pt,
        drop shadow=popink,left=11pt,right=11pt,top=10pt,bottom=10pt,height=3.1cm,valign=center]
      \tcbox[colback=white,colframe=popink,boxrule=1.3pt,arc=5pt,boxsep=0pt,nobeforeafter,
        left=3pt,right=3pt,top=3pt,bottom=3pt,valign=center]{\qrcode[height=1.9cm]{https://play.google.com/store/apps/details?id=com.luvvix.onbuch&hl=fr}}%
      \hspace{8pt}\begin{minipage}[c]{0.5\linewidth}
        {\popdisplay\fontsize{11}{12.5}\selectfont\color{white}\MakeUppercase{Télécharge OnBuch}}\par\vspace{4pt}
        {\raggedright\popsemi\fontsize{8.4}{11}\selectfont\color{white}Gratuit \textbullet\ Google Play\\Tuteur IA, annales, concours, résultats\par}
      \end{minipage}
    \end{tcolorbox}
  \end{minipage}\hfill
  \begin{minipage}[t]{0.487\linewidth}
    \begin{tcolorbox}[enhanced,colback=poppurple,colframe=popink,boxrule=2.2pt,arc=10pt,
        drop shadow=popink,left=11pt,right=11pt,top=10pt,bottom=10pt,height=3.1cm,valign=center]
      \tikz[baseline=-0.7ex]\node[circle,fill=white,draw=popink,line width=1.3pt,minimum size=1.6cm,inner sep=0]{\popdisplay\fontsize{24}{24}\selectfont\color{poppurple}+};%
      \hspace{8pt}\begin{minipage}[c]{0.6\linewidth}
        {\popdisplay\fontsize{11}{12.5}\selectfont\color{white}\MakeUppercase{Passe à OnBuch+}}\par\vspace{4pt}
        {\raggedright\popsemi\fontsize{8.4}{11}\selectfont\color{white}Tous les cours signés, Tuteur IA illimité, fiches PDF — onglet OnBuch+ de l'app\par}
      \end{minipage}
    \end{tcolorbox}
  \end{minipage}
  \vspace{8pt}
  \popcontenu
  \vfill
  \signatureonbuch
  \restoregeometry
  \newpage
}

% Rangée « Ce cours contient » (page de garde)
\newcommand{\poptile}[3]{% #1 couleur #2 gros texte #3 légende
  \begin{tcolorbox}[enhanced,colback=#1,colframe=popink,boxrule=2pt,arc=9pt,shadow={2.5pt}{-2.5pt}{0pt}{popink},
      left=6pt,right=6pt,top=9pt,bottom=9pt,halign=center,valign=center,height=1.75cm,width=\linewidth,nobeforeafter]
    {\popdisplay\fontsize{12.5}{13}\selectfont\color{white}\MakeUppercase{#2}}\par\vspace{3pt}
    {\centering\popsemi\fontsize{7.8}{9.5}\selectfont\color{white}#3\par}
  \end{tcolorbox}}
\newcommand{\popcontenu}{%
  \par\noindent{\popxboldls\fontsize{7.5}{7.5}\selectfont\color{popmuted}CE COURS SIGNÉ CONTIENT}\par\vspace{7pt}
  \noindent\begin{minipage}[t]{0.235\linewidth}\poptile{popblue}{Cours}{complet, illustré, pas à pas}\end{minipage}\hfill
  \begin{minipage}[t]{0.235\linewidth}\poptile{poporange}{Méthodes}{et exercices résolus}\end{minipage}\hfill
  \begin{minipage}[t]{0.235\linewidth}\poptile{poppink}{Exercices}{type examen + corrigés}\end{minipage}\hfill
  \begin{minipage}[t]{0.235\linewidth}\poptile{popgreen}{Fiche bilan}{l'essentiel à retenir}\end{minipage}\par}

% Sommaire
\newtcolorbox{sommaire}{enhanced,colback=white,colframe=popink,boxrule=2.2pt,arc=10pt,
  drop shadow=popink,left=18pt,right=18pt,top=13pt,bottom=13pt,before skip=6pt,after skip=20pt,
  before upper={\poppill{Sommaire}{poppurple}{white}\par\vspace{9pt}\popsemi\fontsize{10.6}{14.5}\selectfont\color{popink}}}

% Signature de fin de cours — poussée tout en bas de la dernière page
\newcommand{\findecours}{%
  \par\vfill\nopagebreak
  \begin{center}\popsquiggle{poporange}\end{center}\vspace{4pt}
  \signatureonbuch
}
\AtBeginDocument{\def\llbracket{\mathopen{[\mkern-3.5mu[}}\def\rrbracket{\mathclose{]\mkern-3.5mu]}}} % intervalles d'entiers (glyphe ⟦ absent de la police)
% Glyphes absents de Fira Math : redéfinis à partir de glyphes disponibles
\AtBeginDocument{%
  \def\setminus{\mathbin{\backslash}}\let\smallsetminus\setminus
  \def\vdots{\mathord{\vbox{\baselineskip3.2pt\lineskiplimit0pt\kern4pt\hbox{.}\hbox{.}\hbox{.}}}}%
  \def\ddots{\mathinner{\mkern1mu\raise6.5pt\hbox{.}\mkern2mu\raise3.6pt\hbox{.}\mkern2mu\raise0.7pt\hbox{.}\mkern1mu}}%
  \def\triangle{\mathord{\scalebox{0.9}{$\symup{\Delta}$}}}%
  \def\nabla{\mathord{\raisebox{\depth}{\scalebox{1}[-1]{$\symup{\Delta}$}}}}%
  \def\checkmark{\popcheck{popdark}}%
  \def\star{\mathbin{\ast}}\def\bigstar{\mathord{\ast}}%
  \def\diamond{\mathbin{\scalebox{0.6}{$\Diamond$}}}%
  \def\bigcup{\mathop{\mathchoice{\raisebox{-0.3em}{\scalebox{1.7}{$\cup$}}}{\scalebox{1.25}{$\cup$}}{\cup}{\cup}}\displaylimits}%
  \def\bigcap{\mathop{\mathchoice{\raisebox{-0.3em}{\scalebox{1.7}{$\cap$}}}{\scalebox{1.25}{$\cap$}}{\cap}{\cap}}\displaylimits}%
  \def\lesseqqgtr{\mathrel{\lessgtr}}\def\bigoplus{\mathop{\oplus}\displaylimits}%
}
\providecommand{\ssi}{si et seulement si} % abréviation fréquente
\AtBeginDocument{\providecommand{\wideparen}{\overparen}} % arc de cercle

%% FIGFIX-BEGIN v4 — correctifs globaux des figures (docs/onbuch-generation/tools/figfix)
\usepackage{adjustbox}\usepackage{etoolbox}
% toute figure TikZ/pgfplots plus large que la ligne est réduite à la largeur disponible
\BeforeBeginEnvironment{tikzpicture}{\begin{adjustbox}{max width=\linewidth}}
\AfterEndEnvironment{tikzpicture}{\end{adjustbox}}
% légendes pgfplots lisibles par-dessus les courbes
\pgfplotsset{every axis/.append style={legend style={fill=white,fill opacity=0.92,text opacity=1,draw=black!15,inner sep=3pt}}}
% étiquettes d'arêtes (syntaxe edge["..."]) avec fond blanc
\tikzset{every edge quotes/.append style={fill=white,inner sep=1.2pt}}
%% FIGFIX-END
\begin{document}

\pagegarde{La santé reproductive}{SVTEEHB}{Tle TI}{Module II : L'Éducation à la Santé}{N07}{4 h}{Cette leçon étudie la régulation hormonale des fonctions reproductrices chez l'homme et chez la femme, établie à partir d'expériences et de courbes de dosages. Elle aborde ensuite les causes de l'infertilité du couple et les techniques de procréation médicalement assistée (insémination artificielle, FIVETE).
\vspace{4pt}
\begin{multicols}{2}\small
\begin{itemize}
\item À la fin de cette leçon, je sais identifier les hormones sexuelles mâles et femelles et leur organe de sécrétion.
\item Je sais analyser des résultats d'expériences (castration, greffe, injections) pour déterminer le rôle des hormones sexuelles.
\item Je sais interpréter des courbes de dosages hormonaux au cours du cycle ovarien.
\item Je sais décrire le rétrocontrôle exercé par les hormones gonadiques sur l'axe hypothalamo-hypophysaire.
\item Je sais citer les principales causes possibles d'infertilité chez l'homme et chez la femme.
\item Je sais élaborer un outil de sensibilisation (affiche, dépliant, sketch) sur les causes de l'infertilité.
\end{itemize}\end{multicols}}

\begin{sommaire}
\begin{multicols}{2}
\begin{enumerate}
\item I. Rappels et mise en place expérimentale : les hormones sexuelles
\item II. Régulation des hormones sexuelles chez l'homme
\item III. Régulation des hormones sexuelles chez la femme
\item IV. L'infertilité : causes possibles chez l'homme et chez la femme
\item V. La procréation médicalement assistée (PMA)
\item VI. Synthèse et intégration
\item Activité d'intégration
\item Méthodes et exercices résolus
\item Exercices
\item Corrigés des exercices
\item Fiche bilan
\end{enumerate}
\end{multicols}
\end{sommaire}

\begin{objectifs}
\begin{itemize}
\item À la fin de cette leçon, je sais identifier les hormones sexuelles mâles et femelles et leur organe de sécrétion.
\item Je sais analyser des résultats d'expériences (castration, greffe, injections) pour déterminer le rôle des hormones sexuelles.
\item Je sais interpréter des courbes de dosages hormonaux au cours du cycle ovarien.
\item Je sais décrire le rétrocontrôle exercé par les hormones gonadiques sur l'axe hypothalamo-hypophysaire.
\item Je sais citer les principales causes possibles d'infertilité chez l'homme et chez la femme.
\item Je sais élaborer un outil de sensibilisation (affiche, dépliant, sketch) sur les causes de l'infertilité.
\item Je sais décrire les étapes de l'insémination artificielle et de la FIVETE.
\item Je sais comparer les indications de l'insémination artificielle et de la FIVETE à partir de cas cliniques.
\end{itemize}
\end{objectifs}

\begin{prerequis}
\begin{itemize}
\item Anatomie des appareils reproducteurs masculin et féminin (classe de Première)
\item Spermatogenèse et ovogenèse : production des gamètes (Première)
\item Notion d'hormone : substance chimique véhiculée par le sang, organe cible (Seconde/Première)
\item Le cycle menstruel et la fécondation (Première)
\item Lecture et interprétation de courbes expérimentales
\end{itemize}
\end{prerequis}

\newpage

\coursec{I. Rappels et mise en place expérimentale : les hormones sexuelles}

\textbf{Situation d'accroche.} À Yaoundé comme à Bafoussam, de nombreux couples consultent les services de gynécologie après plusieurs années de vie commune sans enfant. Dans l'entourage, on entend souvent : « c'est la femme qui est stérile ». Pourtant, les statistiques hospitalières (CHU de Yaoundé, Hôpital Laquintinie de Douala) montrent que l'homme est en cause dans près de 50 \% des cas d'infertilité de couple. Comment le corps humain régule-t-il la production des gamètes et des hormones sexuelles ? Quelles peuvent être les causes d'une infertilité, et comment la science médicale aide-t-elle aujourd'hui ces couples à concevoir ?

\textbf{Problématique.} Pour répondre à ces questions, il faut d'abord identifier les acteurs chimiques de la reproduction : les \cle{hormones sexuelles}. Qui les produit ? Qui commande leur sécrétion ? C'est ce que nous allons établir dans cette première partie, à partir d'expériences historiques rigoureuses, avant d'étudier leur régulation dynamique (sections II et III), puis l'infertilité et la procréation médicalement assistée (sections IV et V).

\courssub{Rappels d'anatomie : les gonades, organes à double fonction}

\begin{definition}[Gonade]
Une \cle{gonade} est une glande génitale qui assure une \textbf{double fonction} :
\begin{itemize}
\item une fonction \textbf{exocrine} : la production des gamètes (spermatozoïdes chez l'homme, ovules chez la femme) ;
\item une fonction \textbf{endocrine} : la sécrétion d'hormones sexuelles déversées directement dans le sang.
\end{itemize}
Chez l'homme, les gonades sont les \cle{testicules} ; chez la femme, ce sont les \cle{ovaires}.
\end{definition}

\textbf{Le testicule.} Chaque testicule est formé de lobules contenant des \cle{tubes séminifères}, longs tubes repliés sur eux-mêmes (plusieurs centaines de mètres au total par testicule) où se déroule la \textbf{spermatogenèse} (production des spermatozoïdes). Entre les tubes séminifères se trouvent des amas de cellules appelées \cle{cellules de Leydig} (ou cellules interstitielles), baignées par de nombreux capillaires sanguins : ce sont elles qui sécrètent la testostérone. À l'intérieur des tubes séminifères, les \textbf{cellules de Sertoli} assurent la nutrition, la protection et le soutien des cellules germinales en cours de différenciation ; elles sécrètent aussi l'inhibine et l'hormone anti-müllérienne (AMH).

\textbf{L'ovaire.} Chaque ovaire contient, dès la naissance, un stock de \cle{follicules} primordiaux (environ \num{1000000} à la naissance, \num{400000} à la puberté, dont seulement 400 à 500 arriveront à maturation au cours de la vie génitale). Chaque follicule renferme une cellule-œuf (ovocyte) en maturation. La paroi du follicule (cellules de la thèque et de la granulose) sécrète les œstrogènes. Après l'ovulation, le follicule rompu se transforme en une glande éphémère, le \cle{corps jaune} (corpus luteum), qui sécrète principalement de la progestérone et des œstrogènes pendant environ 10 à 14 jours avant de régresser (corps blanc) si aucune fécondation n'a lieu.

\begin{aretenir}[Structures productrices à retenir]
\begin{itemize}
\item \textbf{Testicule} : tubes séminifères (spermatogenèse, cellules de Sertoli) + cellules de Leydig (sécrétion de testostérone).
\item \textbf{Ovaire} : follicules (maturation de l'ovocyte + sécrétion d'œstrogènes par les cellules de la granulose) + corps jaune (sécrétion de progestérone et d'œstrogènes).
\end{itemize}
\end{aretenir}

\courssub{Les hormones sexuelles : définition et identification}

\begin{definition}[Hormone sexuelle]
Une \cle{hormone sexuelle} est une substance chimique sécrétée par une \textbf{glande endocrine} (hypothalamus, hypophyse ou gonade), \textbf{véhiculée par le sang}, et qui agit à distance sur des \textbf{organes cibles} de la reproduction (gonades, voies génitales, glandes annexes, caractères sexuels secondaires, métabolisme osseux, cerveau). Elle est active à très faible concentration (de l'ordre du nanogramme par millilitre de sang, \SI{e-9}{g/mL}).
\end{definition}

Les hormones impliquées dans la fonction de reproduction se répartissent en trois niveaux hiérarchiques :

\begin{itemize}
\item \textbf{Neurohormone hypothalamique} : la \cle{GnRH} (Gonadotropin Releasing Hormone), sécrétée par les neurones de l'\textbf{hypothalamus} (aire préoptique médiane).
\item \textbf{Hormones gonadotropes hypophysaires} (commandent les gonades) : la \cle{FSH} (Follicle Stimulating Hormone) et la \cle{LH} (Luteinizing Hormone), deux glycoprotéines sécrétées par l'\textbf{antéhypophyse} (adénohypophyse, partie antérieure de l'hypophyse).
\item \textbf{Hormones gonadiques} (stéroïdes, dérivées du cholestérol) :
\begin{itemize}
\item chez l'homme : la \cle{testostérone}, sécrétée par les cellules de Leydig ;
\item chez la femme : les \cle{œstrogènes} (principalement l'estradiol, sécrétés par les cellules de la granulose des follicules) et la \cle{progestérone} (sécrétée par les cellules du corps jaune).
\end{itemize}
\end{itemize}

\begin{attention}[Erreur fréquente : ne pas confondre les niveaux]
Les élèves confondent souvent FSH/LH avec testostérone/œstrogènes. Retiens bien :
\begin{itemize}
\item \textbf{FSH et LH} sont des hormones \textbf{hypophysaires} (produites par le cerveau, via l'antéhypophyse) : elles \emph{commandent} les gonades. Ce sont des protéines (glycoprotéines).
\item \textbf{Testostérone, œstrogènes, progestérone} sont des hormones \textbf{gonadiques} (produites par les testicules et les ovaires) : ce sont des stéroïdes dérivés du cholestérol.
\end{itemize}
Dire « la FSH est sécrétée par les ovaires » ou « la testostérone est sécrétée par l'hypophyse » est une erreur grave qui coûte des points au Baccalauréat.
\end{attention}

\courssub{Mise en évidence expérimentale du rôle endocrine des gonades et de l'hypophyse}

Comment a-t-on découvert que les gonades produisent des substances chimiques actives et que l'hypophyse les commande ? Par la démarche classique d'\textbf{ablation--greffe--injection}.

\begin{experience}[Expérience 1 : castration et greffe chez le coq (Berthold, 1849)]
\textbf{Protocole.} On compare trois lots de jeunes coqs de même âge et race :
\begin{itemize}
\item \textbf{Lot témoin} : coqs intacts (non opérés) ;
\item \textbf{Lot A} : coqs \textbf{castrés} (ablation chirurgicale des deux testicules) ;
\item \textbf{Lot B} : coqs castrés, puis recevant une \textbf{greffe de testicule} dans la cavité abdominale (sans connexion nerveuse avec le système nerveux central) ou des injections régulières d'extrait testiculaire.
\end{itemize}
\textbf{Observations (après plusieurs semaines).}
\begin{itemize}
\item Lot témoin : développement normal de la crête et des barbillons (rouges, volumineux), comportement de cour et de combat, chant du coq.
\item Lot A : \textbf{régression des caractères sexuels secondaires} : la crête et les barbillons restent atrophiés et pâles, le comportement sexuel et agressif disparaît, les voies génitales (vésicules séminales) s'atrophient.
\item Lot B : \textbf{restauration} des caractères sexuels : la crête se développe et se colore, le comportement de coq réapparaît, les vésicules séminales regagnent leur volume, alors même que le testicule greffé n'est relié au reste du corps \emph{que par les vaisseaux sanguins} (vascularisation néoformée).
\end{itemize}
\textbf{Interprétation.} La restauration se fait sans connexion nerveuse : le message ne peut donc pas être nerveux. Le testicule libère dans le \textbf{sang} une substance chimique active, responsable des caractères sexuels mâles. Cette substance sera identifiée plus tard (1935) : la \textbf{testostérone}.
\textbf{Conclusion.} Le testicule est une glande endocrine : il produit une hormone (la testostérone) véhiculée par le sang, qui agit sur les organes cibles (crête, cerveau, voies génitales, muscles, os).
\end{experience}

\begin{methode}[Interpréter une expérience d'ablation--greffe--injection (type Bac)]
Pour analyser ce type d'expérience (très fréquent au Bac), procède en quatre étapes rigoureuses :
\begin{enumerate}
\item \textbf{Identifier les lots} : le lot témoin (non opéré) sert de référence physiologique normale ; le lot opéré (ablation) montre l'effet de l'absence de l'organe ; le lot opéré + greffe ou injection d'extrait teste si l'organe suffit à restaurer le phénomène.
\item \textbf{Comparer} les résultats du lot opéré au lot témoin : toute différence significative est due à l'ablation de l'organe.
\item \textbf{Comparer} le lot « opéré + greffe/injection » au lot opéré : si le phénomène est restauré (retour à la normale), l'organe (ou son extrait) contient la substance active.
\item \textbf{Conclure} sur la nature du message : si la restauration se fait sans lien nerveux (greffe ectopique hors de la localisation naturelle, ou extrait injecté par voie sanguine), la substance agit par voie \textbf{sanguine} : c'est une \textbf{hormone}, et l'organe est une \textbf{glande endocrine}.
\end{enumerate}
\end{methode}

\begin{experience}[Expérience 2 : hypophysectomie chez le rat]
\textbf{Protocole.} On réalise l'\textbf{ablation de l'hypophyse} (hypophysectomie) chez de jeunes rats mâles et femelles prépubères, puis on injecte à certains d'entre eux des extraits d'antéhypophyse ou des hormones gonadiques.
\textbf{Observations.}
\begin{itemize}
\item Rats témoins : développement normal des testicules et des ovaires, puberté normale, gamétogenèse active.
\item Rats hypophysectomisés : \textbf{atrophie des gonades} (testicules et ovaires restent infantiles, tubes séminifères réduits, pas de follicules matures), arrêt complet de la gamétogenèse, absence de puberté. L'injection de testostérone (mâle) ou d'œstrogènes (femelle) restaure les caractères sexuels secondaires (poils, voix, musculature, muqueuses) mais \emph{ne restaure pas} la production de gamètes (spermatogenèse, ovogenèse).
\item Rats hypophysectomisés + extraits d'antéhypophyse : \textbf{restauration} du développement des gonades, de la gamétogenèse et de la sécrétion d'hormones gonadiques.
\end{itemize}
\textbf{Interprétation.} L'hypophyse (plus précisément l'antéhypophyse) sécrète des substances indispensables au développement et au fonctionnement des gonades (croissance, gamétogenèse, stéroïdogenèse). Leur fractionnement biochimique a permis d'isoler deux hormones distinctes : la \textbf{FSH} et la \textbf{LH}, dites \cle{gonadotropines} (ou gonadotrophines).
\textbf{Conclusion.} L'antéhypophyse contrôle les gonades par l'intermédiaire de la FSH et de la LH, véhiculées par le sang. La FSH et la LH agissent \emph{directement sur la gonade} (récepteurs membranaires sur cellules de Sertoli/Leydig ou granulose/thèque), et non en imitant les stéroïdes.
\end{experience}

\begin{attention}[Piège d'interprétation majeur]
Dans l'expérience 2, l'injection d'hormones gonadiques (testostérone) à un animal hypophysectomisé restaure les caractères sexuels secondaires mais \textbf{pas} la spermatogenèse (ni l'ovogenèse). Cela prouve formellement que la FSH et la LH agissent \emph{directement sur la gonade} pour stimuler la gamétogenèse et la stéroïdogenèse locale, et non simplement en imitant l'action des stéroïdes sur les organes cibles périphériques. Ne conclus jamais trop vite : précise toujours \emph{quel} effet est restauré par \emph{quelle} injection.
\end{attention}

\courssub{L'axe hypothalamo-hypophyso-gonadique}

Les expériences complémentaires (lésions localisées de l'hypothalamus, section de la tige pituitaire, dosages sanguins porte vs périphérique) ont montré que l'hypophyse elle-même est commandée par l'\textbf{hypothalamus}, une région du cerveau qui sécrète la \cle{GnRH}. Cette neurohormone gagne l'antéhypophyse par un court réseau de vaisseaux sanguins (le \textbf{système porte hypothalamo-hypophysaire}) et y déclenche la libération pulsatile de FSH et de LH.

L'ensemble forme un axe de commande à trois étages, appelé \cle{axe hypothalamo-hypophyso-gonadique} :

\begin{popfigure}
\begin{center}
\begin{tikzpicture}[
  node distance=1.8cm,
  organe/.style={rectangle, rounded corners=6pt, draw=popdark, thick, fill=popblueL, minimum width=4.5cm, minimum height=1.1cm, align=center, font=\small\bfseries},
  hormone/.style={font=\small\itshape, text=popdark},
  fleche/.style={-{Stealth[length=3mm]}, very thick, poporange},
  cible/.style={rectangle, rounded corners=6pt, draw=popdark, thick, fill=popgreenL, minimum width=4.5cm, minimum height=1.1cm, align=center, font=\small\bfseries}
]
\node[organe, align=center] (hypo){Hypothalamus\\ (neurones neurosecréteurs)};
\node[organe, below=of hypo, fill=poppurpleL, align=center] (hypophyse){Antéhypophyse\\ (cellules gonadotropes)};
\node[cible, below left=2.0cm and 0.5cm of hypophyse, align=center] (test){Testicule\\ (cellules de Leydig, tubes séminifères,\\ cellules de Sertoli)};
\node[cible, below right=2.0cm and 0.5cm of hypophyse, align=center] (ovaire){Ovaire\\ (follicules : cellules thèque/granulose,\\ corps jaune)};

\draw[fleche] (hypo) -- node[hormone, right, xshift=3pt] {GnRH} (hypophyse);
\draw[fleche] (hypophyse.south west) -- node[hormone, left, xshift=-3pt, align=right] {FSH + LH} (test.north);
\draw[fleche] (hypophyse.south east) -- node[hormone, right, xshift=3pt, align=left] {FSH + LH} (ovaire.north);

\node[hormone, below=0.3cm of test, text=popblue, align=center]{Testostérone\\(+ Inhibine)};
\node[hormone, below=0.3cm of ovaire, text=popblue, align=center]{Œstrogènes + Progestérone\\(+ Inhibine)};
\end{tikzpicture}
\end{center}
\legende{Schéma fonctionnel de l'axe hypothalamo-hypophyso-gonadique. L'hypothalamus sécrète la GnRH (voie porte), qui stimule l'antéhypophyse ; celle-ci libère la FSH et la LH (voie sanguine générale), qui agissent sur les gonades ; les gonades produisent alors les hormones sexuelles stéroïdes et les gamètes. Les flèches oranges indiquent la voie de stimulation.}
\end{popfigure}

\begin{center}
\begin{tabularx}{\linewidth}{|l|X|X|X|}
\hline
\rowcolor{popblueL} \textbf{Hormone} & \textbf{Organe de sécrétion} & \textbf{Organe cible principal} & \textbf{Rôle principal dans la reproduction} \\
\hline
GnRH & Hypothalamus (neurones) & Antéhypophyse (cellules gonadotropes) & Déclenche la synthèse et la libération pulsatile de FSH et de LH \\
\hline
FSH & Antéhypophyse & Gonades & \textbf{Homme} : stimule les cellules de Sertoli (spermatogenèse, sécrétion d'inhibine/AMH) \\
& & & \textbf{Femme} : stimule la croissance folliculaire et la synthèse d'œstrogènes (cellules de la granulose) \\
\hline
LH & Antéhypophyse & Gonades & \textbf{Homme} : stimule les cellules de Leydig (synthèse de testostérone) \\
& & & \textbf{Femme} : déclenche l'ovulation, luteinisation (formation du corps jaune), sécrétion de progestérone \\
\hline
Testostérone & Cellules de Leydig (testicule) & Voies génitales, muscles, os, peau, cerveau, cellules de Sertoli & Caractères sexuels secondaires masculins, maintien de la spermatogenèse (synergie FSH), rétroaction négative \\
\hline
Œstrogènes (Estradiol) & Follicules ovariens (cell. granulose) & Utérus, seins, os, cerveau, hypophyse & Caractères sexuels féminins, prolifération endomètre, rétroaction négative/positive \\
\hline
Progestérone & Corps jaune (ovaire) & Utérus, seins, hypothalamus/hypophyse & Préparation endomètre à la nidation, maintien grossesse, rétroaction négative, blocage LH \\
\hline
Inhibine & Cellules de Sertoli (H) / Granulose (F) & Antéhypophyse & Inhibition sélective de la sécrétion de FSH (rétroaction négative) \\
\hline
\end{tabularx}
\end{center}

\courssub{Le dosage des hormones : la méthode radio-immunologique (RIA)}

Comment mesure-t-on des hormones présentes dans le sang à des concentrations infimes (parfois moins de \SI{1}{ng/mL}, soit \SI{e-9}{g/mL}) ? On utilise le \cle{dosage radio-immunologique} (RIA, mis au point par Rosalyn Yalow et Solomon Berson, prix Nobel de physiologie/médecine 1977). Aujourd'hui, les dosages immuno-enzymatiques (ELISA) ou chimiluminescents, non radioactifs, ont largement remplacé le RIA en routine, mais le principe de \textbf{compétition} reste identique.

\begin{definition}[Principe du dosage par compétition immuno-chimique (admis)]
Le principe repose sur la \textbf{compétition} entre l'hormone à doser (présente dans l'échantillon de sang/urine du patient, \textbf{non marquée}) et une quantité connue de la même hormone \textbf{marquée} (radioactivement pour RIA, par enzyme pour ELISA, etc.), pour se fixer sur un nombre limité et limitant d'\textbf{anticorps} spécifiques (anticorps monoclonaux anti-hormone).
\begin{itemize}
\item Plus l'échantillon contient d'hormone naturelle, plus celle-ci « chasse » l'hormone marquée des anticorps.
\item On sépare ensuite les complexes anticorps-hormone de l'hormone libre (par précipitation, adsorption sur charbon, ou puits ELISA).
\item On mesure le signal (radioactivité, absorbance, luminescence) lié à la fraction \textbf{liée aux anticorps}.
\item \textbf{Résultat} : le signal mesuré est \textbf{inversement proportionnel} à la concentration de l'hormone dans l'échantillon. Une courbe étalon (connue) permet de convertir le signal en concentration hormonale.
\end{itemize}
\end{definition}

\begin{methode}[Exploiter un résultat de dosage hormonal (démarche diagnostique)]
\begin{enumerate}
\item \textbf{Relever} la valeur mesurée et son unité (ex. testostérone : \SI{5,2}{ng/mL} ; FSH : \SI{8,5}{UI/L} ; LH : \SI{6,2}{UI/L} ; Progestérone : \SI{15}{ng/mL}).
\item \textbf{Comparer} aux \textbf{valeurs de référence} fournies dans l'énoncé (varient selon l'âge, le sexe, le moment du cycle pour la femme). Exemple valeurs de référence homme adulte : Testostérone \SIrange{3}{10}{ng/mL} ; FSH \SIrange{1,5}{12}{UI/L} ; LH \SIrange{1,5}{9}{UI/L}.
\item \textbf{Conclure} pour chaque hormone : valeur normale, effondrée (hypo-sécrétion) ou élevée (hyper-sécrétion).
\item \textbf{Corréler} les profils (FSH/LH vs Hormones gonadiques) pour localiser le niveau de la pathologie :
\begin{itemize}
\item FSH/LH basses + Hormones gonadiques basses $\rightarrow$ \textbf{Atteinte centrale} (hypothalamo-hypophysaire).
\item FSH/LH élevées + Hormones gonadiques basses $\rightarrow$ \textbf{Atteinte périphérique} (insuffisance gonadique primitive : testicule/ovaire ne répondent pas, perte de rétroaction négative).
\item FSH/LH normales/élevées + Hormones gonadiques élevées $\rightarrow$ Tumeur sécrétrice ou pathologie surrénalienne.
\end{itemize}
\item Cette démarche est la base du \textbf{bilan hormonal d'infertilité} (section IV).
\end{enumerate}
\end{methode}

\begin{exoresolu}[Exploitation des expériences d'ablation--greffe--injection]
\textbf{Énoncé.} Chez trois lots de rats mâles adultes, on réalise les opérations suivantes. Lot 1 : rats témoins non opérés. Lot 2 : ablation bilatérale des testicules (castration). Lot 3 : ablation bilatérale des testicules puis injection quotidienne sous-cutanée d'extrait testiculaire (contenant de la testostérone). Après 8 semaines, on mesure la masse des vésicules séminales (glandes annexes dont le développement et le maintien dépendent strictement des androgènes) : Lot 1 : \SI{450}{mg} ; Lot 2 : \SI{60}{mg} ; Lot 3 : \SI{430}{mg}. Interpréter ces résultats et conclure sur la nature du facteur testiculaire.
\tcblower
\textbf{Solution détaillée.}
\begin{enumerate}
\item \textbf{Comparaison Lot 2 / Lot 1 (effet de l'ablation)} : La castration provoque une chute drastique de la masse des vésicules séminales, passant de \SI{450}{mg} (témoin) à \SI{60}{mg} (castré), soit une régression de \SI{87}{\%} environ. L'ablation des testicules supprime donc un facteur indispensable au développement et au maintien des glandes annexes.
\item \textbf{Comparaison Lot 3 / Lot 2 (effet de l'injection)} : L'injection quotidienne d'extrait testiculaire restaure presque totalement la masse des vésicules séminales (\SI{430}{mg}, valeur non significativement différente du témoin \SI{450}{mg}). L'extrait contient donc le facteur actif absent chez les castrés.
\item \textbf{Nature de la voie de communication} : Le facteur actif est injecté par voie sous-cutanée (dans le sang) et agit à distance sur les vésicules séminales, sans qu'aucune connexion nerveuse n'ait été réalisée entre le site d'injection et l'organe cible. Le message est donc \textbf{humoral (sanguin)}.
\item \textbf{Conclusion} : Le testicule est une \textbf{glande endocrine}. Il sécrète dans le sang une hormone (identifiée comme la \textbf{testostérone}, produite par les cellules de Leydig) responsable du développement et du maintien des organes génitaux accessoires (vésicules séminales, prostate) et des caractères sexuels secondaires.
\end{enumerate}
\end{exoresolu}

\begin{savaistu}[Histoire des sciences : de Berthold à la contraception camerounaise]
Le mot « hormone » vient du grec \emph{hormao} (« j'excite, je mets en mouvement »), introduit en 1905 par Ernest Starling. Mais la première démonstration expérimentale d'une sécrétion interne remonte à \textbf{1849} : le médecin allemand \textbf{Arnold Berthold} castra des coqs, observe la régression de leur crête, puis greffe un testicule dans l'abdomen de certains et constate la restauration des caractères mâles — sans connexion nerveuse. Il fondait, sans le nommer, l'endocrinologie.
Aujourd'hui, au Cameroun, les dosages hormonaux (FSH, LH, Testostérone, Progestérone, AMH, Prolactine, TSH) font partie du \textbf{bilan de routine} des couples consultant pour infertilité dans les centres spécialisés (CHU, Hôpital Général, cliniques privées). La compréhension de l'axe hypothalamo-hypophyso-gonadique a aussi permis le développement de la \textbf{contraception hormonale} (pilule œstro-progestative, injectables, implants) largement disponible dans les centres de planification familiale (ex. CAMNAFAW, centres de santé de district), permettant aux femmes de choisir l'espacement des naissances.
\end{savaistu}

\begin{aretenir}[L'essentiel de la section I]
\begin{itemize}
\item Les gonades (testicules, ovaires) ont une double fonction : production des gamètes (exocrine) et sécrétion d'hormones sexuelles stéroïdes (endocrine).
\item \textbf{Testicule} : cellules de Leydig $\rightarrow$ Testostérone ; tubes séminifères (cellules de Sertoli) $\rightarrow$ Spermatogenèse + Inhibine/AMH.
\item \textbf{Ovaire} : Follicules (cellules granulose/thèque) $\rightarrow$ Œstrogènes ; Corps jaune $\rightarrow$ Progestérone + Œstrogènes.
\item La régulation est hiérarchique : \textbf{Hypothalamus (GnRH)} $\rightarrow$ \textbf{Antéhypophyse (FSH, LH)} $\rightarrow$ \textbf{Gonades (Stéroïdes, Inhibine)}. C'est l'\textbf{axe hypothalamo-hypophyso-gonadique}.
\item Les expériences d'\textbf{ablation--greffe--injection} (Berthold, 1849) et d'\textbf{hypophysectomie} prouvent la nature hormonale (sanguine) de ces commandes et distinguent le rôle des gonadotropines (action sur la gonade) de celui des stéroïdes (action sur organes cibles périphériques).
\item Le \textbf{dosage immuno-chimique} (RIA/ELISA) permet de quantifier ces hormones dans le sang (ng/mL ou UI/L) : outil indispensable du diagnostic d'infertilité et du suivi de PMA.
\end{itemize}
\end{aretenir}

\coursec{II. Régulation des hormones sexuelles chez l'homme}

Dans la section précédente, nous avons identifié les principales hormones sexuelles : les hormones hypophysaires (FSH et LH) et les hormones gonadiques (testostérone chez l'homme, \oe{}strogènes et progestérone chez la femme). Nous avons aussi vu que l'hypothalamus, par l'intermédiaire de la GnRH, commande l'activité de l'hypophyse. Une question se pose alors naturellement : comment tous ces acteurs sont-ils coordonnés pour maintenir une production continue de spermatozoïdes chez l'homme adulte ? C'est l'objet de cette section : montrer, à partir de données expérimentales, que la fonction reproductrice masculine est contrôlée par un mécanisme élégant appelé \cle{rétrocontrôle négatif}, complété par une régulation spécifique de la FSH par l'\cle{inhibine}.

\courssub{Un fonctionnement continu : le constat expérimental}

\textbf{Observation.} Contrairement à la femme, dont les taux d'hormones sexuelles varient de façon cyclique au cours du cycle menstruel, l'homme présente des taux sanguins de FSH, de LH et de testostérone remarquablement \textbf{stables} au cours du temps. Il n'existe pas de « cycle » masculin : la spermatogenèse se poursuit de façon continue, jour et nuit, de la puberté jusqu'à un âge avancé.

\begin{experience}[Dosages hormonaux chez l'homme adulte sur 24 h]
\textbf{Protocole.} On prélève du sang toutes les 2 heures pendant 24 heures chez un homme adulte sain, et on dose la LH et la testostérone par des méthodes immunologiques (radio-immunologie ou ELISA).

\textbf{Observations.} Les taux de LH oscillent légèrement autour d'une valeur moyenne (environ \num{5} UI/L), reflétant la sécrétion pulsatile de la GnRH. Le taux de testostérone reste également stable, avec de faibles fluctuations nycthémérales (il est légèrement plus élevé le matin que le soir).

\textbf{Interprétation.} L'absence de variation brutale traduit une sécrétion continue et régulée des hormones hypophysaires et testiculaires. Il existe donc un mécanisme de régulation qui maintient ces taux dans des limites étroites, comme un thermostat maintient la température d'une pièce.
\end{experience}

\begin{popfigure}
\begin{center}
\begin{tikzpicture}
\begin{axis}[
    width=0.92\linewidth, height=6.2cm,
    xlabel={Temps (h)}, ylabel={Taux sanguin},
    xmin=0, xmax=24, ymin=0, ymax=12,
    xtick={0,4,8,12,16,20,24},
    ytick={0,2,4,6,8,10,12},
    legend style={at={(0.02,0.97)}, anchor=north west, font=\small},
    grid=major, grid style={popline!40},
    axis line style={popdark},
]
% LH : pulsatile simulé par une série de pics gaussiens lissés + bruit
\addplot[popcurve, color=popblue, domain=0:24, samples=500]
    {5 + 1.2*exp(-((x-2)/0.5)^2) + 1.0*exp(-((x-5.5)/0.4)^2) + 0.8*exp(-((x-9)/0.6)^2) + 1.1*exp(-((x-12.5)/0.5)^2) + 0.9*exp(-((x-16)/0.4)^2) + 1.3*exp(-((x-19.5)/0.5)^2) + 0.7*exp(-((x-23)/0.6)^2)};
\addlegendentry{LH (UI/L) -- pulsatile}
% Testostérone : lissée, suit vaguement LH avec décalage
\addplot[popcurve, color=poporange, domain=0:24, samples=200]
    {7 + 0.8*sin(deg(x*0.26)) + 0.3*sin(deg(x*1.5))};
\addlegendentry{Testostérone (ng/mL) -- lissée}
\end{axis}
\end{tikzpicture}
\end{center}
\legende{Évolution des taux sanguins de LH (pulsatile, reflet de la GnRH) et de testostérone (lissée) chez l'homme adulte au cours d'une journée. Les fluctuations de la testostérone sont faibles autour d'une valeur moyenne stable : il n'y a pas de cycle hormonal masculin.}
\end{popfigure}

\begin{exemplebox}[Ordres de grandeur à connaître]
Chez l'homme adulte sain, le taux de testostérone sanguine est d'environ \textbf{3 à \num{10} ng/mL} de sang (soit \num{3} à \num{10} µg/L). Le taux de LH est de l'ordre de 2 à 9 UI/L et celui de FSH de 1 à 8 UI/L. Ces valeurs, mesurées dans les laboratoires d'analyses médicales (par exemple au Centre Hospitalier d'Essos à Yaoundé), permettent au médecin de vérifier le bon fonctionnement de l'axe de reproduction.
\end{exemplebox}

\courssub{Les rôles respectifs de la LH et de la FSH}

Comment l'hypophyse contrôle-t-elle le testicule ? L'expérimentation animale (castration, injections d'extraits hypophysaires, hypophysectomie) a permis d'établir précisément le rôle de chaque gonadostimuline.

\begin{experience}[Mise en évidence des rôles de la LH et de la FSH]
\textbf{Protocole.} Chez un rat mâle, on pratique l'ablation de l'hypophyse (hypophysectomie) : les testicules s'atrophient et la spermatogenèse s'arrête. On injecte ensuite séparément, à des lots différents d'animaux hypophysectomisés, soit de la LH pure, soit de la FSH pure.

\textbf{Résultats.}
\begin{itemize}
\item L'injection de \textbf{LH} provoque le développement des \textbf{cellules de Leydig} (cellules interstitielles situées entre les tubes séminifères) et la reprise de la sécrétion de \textbf{testostérone}. En revanche, les tubes séminifères restent peu actifs : la spermatogenèse ne reprend pas complètement.
\item L'injection de \textbf{FSH} restaure l'activité des \textbf{tubes séminifères} en agissant sur les \textbf{cellules de Sertoli} : la spermatogenèse reprend partiellement, mais les caractères sexuels secondaires (dépendant de la testostérone) ne se développent pas.
\item L'injection \textbf{simultanée} de FSH et de LH restaure totalement la fonction testiculaire : spermatogenèse complète et sécrétion normale de testostérone.
\end{itemize}

\textbf{Interprétation.} Les cellules de Sertoli, stimulées par la FSH, sécrètent l'\textbf{inhibine}, une hormone qui freine spécifiquement la sécrétion de FSH par l'antéhypophyse (régulation fine).
\textbf{Conclusion.} La LH et la FSH ont des cibles distinctes et complémentaires au niveau du testicule.
\end{experience}

\begin{aretenir}[Rôles des gonadostimulines et de l'inhibine chez l'homme]
\begin{itemize}
\item La \cle{LH} (hormone lutéinisante) stimule les \textbf{cellules de Leydig} des testicules, qui sécrètent la \textbf{testostérone}.
\item La \cle{FSH} (hormone folliculo-stimulante) agit sur les \textbf{cellules de Sertoli} des tubes séminifères ; elle est indispensable au \textbf{maintien de la spermatogenèse} (les cellules de Sertoli nourrissent et accompagnent les cellules germinales en cours de différenciation).
\item Les \textbf{cellules de Sertoli} sécrètent l'\cle{inhibine}, qui exerce un \textbf{rétrocontrôle négatif spécifique sur la FSH} au niveau de l'antéhypophyse.
\item La \cle{testostérone}, sécrétée en continu, assure : le déroulement de la \textbf{spermatogenèse} (en synergie avec la FSH), l'apparition et le maintien des \textbf{caractères sexuels secondaires} (pilosité, mue de la voix, développement musculaire) et la \textbf{libido}.
\end{itemize}
\end{aretenir}

\courssub{La preuve du rétrocontrôle : l'expérience d'injection de testostérone}

Reste à comprendre \emph{comment} les taux hormonaux restent stables. Si l'hypophyse stimulait le testicule sans aucun frein, la testostérone s'accumulerait indéfiniment. Or ce n'est pas le cas. L'expérience décisive est la suivante.

\begin{experience}[Injection de fortes doses de testostérone chez l'animal]
\textbf{Hypothèse.} La testostérone exerce un contrôle sur les organes qui commandent sa propre sécrétion.

\textbf{Protocole.} On injecte à un rat mâle intact des doses élevées de testostérone pendant plusieurs jours. On dose ensuite la FSH et la LH sanguines et on examine l'histologie des testicules.

\textbf{Observations.}
\begin{itemize}
\item Les taux sanguins de \textbf{FSH et de LH chutent fortement}.
\item Les testicules \textbf{s'atrophient} et la \textbf{spermatogenèse s'arrête} (privée de FSH et de testostérone locale élevée, les cellules germinales ne se développent plus).
\item L'arrêt des injections ramène progressivement les taux de FSH et de LH à la normale, et la spermatogenèse reprend.
\end{itemize}

\textbf{Interprétation.} La testostérone injectée en excès agit sur l'\textbf{hypothalamus} (diminution de la sécrétion de GnRH) et sur l'\textbf{antéhypophyse} (diminution de la sensibilité à la GnRH), ce qui réduit la libération de FSH et de LH. Le testicule, moins stimulé, réduit alors sa propre production de testostérone \emph{et} d'inhibine. C'est un freinage : la testostérone limite sa propre sécrétion.

\textbf{Conclusion.} Il existe un \cle{rétrocontrôle négatif} (ou rétroaction négative) exercé par la testostérone sur l'axe hypothalamo-hypophysaire.
\end{experience}

\begin{definition}[Rétrocontrôle négatif]
Le \cle{rétrocontrôle négatif} est un mécanisme de régulation par lequel une hormone \textbf{freine} (inhibe) la sécrétion des hormones qui commandent sa propre production. Il fonctionne comme un thermostat : dès que le taux de l'hormone dépasse une valeur de consigne, la chaîne de commande est ralentie ; dès que le taux baisse, le frein se relâche et la production repart. Ce mécanisme assure la \textbf{stabilité} des taux hormonaux.
\end{definition}

\begin{propriete}[Sécrétion pulsatile de la GnRH (admise)]
La sécrétion de GnRH par l'hypothalamus est \textbf{pulsatile} : des bouffées de GnRH sont libérées toutes les 60 à 90 minutes environ. Ce rythme conditionne la libération de FSH et de LH par l'antéhypophyse. (Propriété admise ; une perfusion continue de GnRH inhibe paradoxalement l'hypophyse, fait exploité par certains médicaments analogues de la GnRH.)
\end{propriete}

\courssub{Schéma fonctionnel de l'axe hypothalamo-hypophyso-testiculaire}

On peut résumer l'ensemble des interactions par un schéma fonctionnel, à savoir reproduire et légender le jour du Baccalauréat.

\begin{popfigure}
\begin{center}
\begin{tikzpicture}[
    node distance=1.6cm,
    org/.style={rectangle, rounded corners=4pt, draw=popdark, fill=popcream, thick, minimum width=3.4cm, minimum height=0.95cm, align=center, font=\small},
    horm/.style={font=\small\itshape, text=popdark},
    stim/.style={-{Stealth[length=3mm]}, very thick, popgreen!70!black},
    inhib/.style={-{Stealth[length=3mm]}, very thick, poppink!80!black, dashed}
]
\node[org] (hypoth) {Hypothalamus};
\node[org, below=of hypoth] (hypo) {Antéhypophyse};
\node[org, below left=1.7cm and -0.4cm of hypo, align=center] (leydig){Cellules de Leydig\\ (interstitielles)};
\node[org, below right=1.7cm and -0.4cm of hypo, align=center] (sertoli){Cellules de Sertoli\\ (tubes séminifères)};
\node[org, below=1.7cm of leydig, align=center] (cible){Spermatogenèse,\\ caractères sexuels\\ secondaires, libido};

% Stimulations
\draw[stim] (hypoth) -- node[horm, right]{GnRH} (hypo);
\draw[stim] (hypo.south west) -- node[horm, left, pos=0.45]{LH} (leydig.north);
\draw[stim] (hypo.south east) -- node[horm, right, pos=0.45]{FSH} (sertoli.north);
\draw[stim] (leydig) -- node[horm, left]{Testostérone} (cible);
\draw[stim] (sertoli.south) |- node[horm, right, pos=0.75]{soutien de la spermatogenèse} ($(cible.east)+(0,0.15)$) -- (cible.east);

% Rétrocontrôles
\draw[inhib] (leydig.west) -- ++(-2.6,0) |- node[horm, left, pos=0.25]{rétrocontrôle négatif ($-$)} (hypoth.west);
\draw[inhib] ($(leydig.west)+(-2.6,0)$) |- (hypo.west);
% Inhibine
\draw[inhib] (sertoli.east) -- ++(1.8,0) |- node[horm, right, pos=0.25]{Inhibine ($-$ sur FSH)} (hypo.east);

\node[font=\footnotesize, text=popmuted, align=left] at ($(hypoth.north west)+(-2.4,0.9)$) {($-$) : inhibition};
\end{tikzpicture}
\end{center}
\legende{Schéma fonctionnel de la régulation des hormones sexuelles chez l'homme. Flèches vertes pleines : stimulations (GnRH $\rightarrow$ FSH/LH $\rightarrow$ testostérone, inhibine, spermatogenèse). Flèches rouges pointillées : rétrocontrôles négatifs. La testostérone inhibe l'hypothalamus (GnRH) et l'antéhypophyse (FSH, LH). L'inhibine (cellules de Sertoli) inhibe spécifiquement la FSH.}
\end{popfigure}

\begin{methode}[Construire le raisonnement sur un rétrocontrôle (méthode Bac)]
\begin{enumerate}
\item \textbf{Identifier la chaîne de commande} : hypothalamus (GnRH) $\rightarrow$ antéhypophyse (FSH, LH) $\rightarrow$ gonade (testostérone, inhibine).
\item \textbf{Repérer la perturbation} dans l'expérience proposée (injection d'hormone, ablation d'une glande).
\item \textbf{Comparer} les taux mesurés avant/après la perturbation et confronter à l'hypothèse.
\item \textbf{Conclure} en nommant le mécanisme : si l'hormone périphérique fait chuter les hormones qui la commandent, c'est un rétrocontrôle \emph{négatif}. Préciser si c'est global (testostérone sur GnRH/LH/FSH) ou spécifique (inhibine sur FSH).
\end{enumerate}
\end{methode}

\begin{attention}[L'erreur classique à ne jamais commettre]
Beaucoup de candidats écrivent que « la testostérone stimule l'hypophyse ». C'est \textbf{faux} : c'est exactement l'inverse. La testostérone \textbf{inhibe} l'hypothalamus et l'antéhypophyse (rétrocontrôle \emph{négatif}). Autre piège : ne pas confondre les cibles de la LH (cellules de \textbf{Leydig}, sécrétion de testostérone) et de la FSH (cellules de \textbf{Sertoli}, spermatogenèse + sécrétion d'inhibine). Enfin, le rétrocontrôle porte sur l'axe hypothalamo-hypophysaire, jamais directement sur les cellules germinales.
\end{attention}

\courssub{Application : le principe de la contraception masculine hormonale}

\begin{savaistu}[Vers un « contraceptif » masculin ?]
L'expérience précédente a une application directe, parfois évoquée au Baccalauréat en guise de complément : la \textbf{contraception masculine hormonale}. L'administration régulière de testostérone (souvent associée à un progestatif potentialisant le blocage gonadotrope) à un homme entretient un rétrocontrôle négatif permanent : les taux de FSH et de LH s'effondrent, et la \textbf{spermatogenèse s'interrompt} (azoospermie : absence de spermatozoïdes dans le sperme). Les caractères sexuels secondaires et la libido sont préservés grâce à la testostérone apportée de l'extérieur. L'effet est \textbf{réversible} à l'arrêt du traitement. Des essais cliniques sont en cours dans plusieurs pays ; cette méthode n'est pas encore commercialisée, notamment en raison d'effets secondaires à long terme encore mal évalués (prise de poids, acné, modifications lipidiques).
\end{savaistu}

\begin{exoresolu}[Interprétation d'un document expérimental]
\textbf{Énoncé.} Chez un singe adulte, on mesure les taux plasmatiques de LH et de testostérone avant et après injection quotidienne de testostérone pendant 15 jours. Résultats (valeurs moyennes) :

\begin{center}
\begin{tabularx}{\linewidth}{|l|X|X|}
\hline
\rowcolor{popblueL} \textbf{Paramètre} & \textbf{Avant injection} & \textbf{Après 15 jours d'injection} \\
\hline
LH plasmatique (UI/L) & 6 & 1 \\
\hline
Testostérone endogène (ng/mL) & 6 & 0,5 \\
\hline
Spermatozoïdes dans le sperme & nombreux & absents \\
\hline
\end{tabularx}
\end{center}

\begin{enumerate}
\item Analyser les résultats obtenus.
\item Déduire le mécanisme de régulation mis en jeu et le nommer.
\item Expliquer l'absence de spermatozoïdes.
\end{enumerate}

\tcblower
\textbf{Solution détaillée.}

\textbf{1. Analyse.} Après 15 jours d'injection de testostérone, le taux de LH passe de 6 à 1 UI/L, soit une chute d'environ 83 \%. Le taux de testostérone \emph{endogène} (produite par le testicule) s'effondre de 6 à 0,5 ng/mL, et les spermatozoïdes disparaissent du sperme. L'injection de testostérone a donc provoqué à la fois la baisse de la LH et l'arrêt de l'activité testiculaire.

\textbf{2. Mécanisme.} La testostérone exogène (injectée) agit sur l'hypothalamus et l'antéhypophyse en freinant la sécrétion de GnRH, puis de LH et de FSH. Le testicule, moins stimulé par la LH, réduit sa production propre de testostérone. Il s'agit d'un \textbf{rétrocontrôle négatif} (rétroaction négative) de la testostérone sur l'axe hypothalamo-hypophysaire.

\textbf{3. Absence de spermatozoïdes.} La spermatogenèse exige l'action conjointe de la FSH (sur les cellules de Sertoli) et d'une forte concentration locale de testostérone dans les tubes séminifères. La chute de FSH et de LH supprime ces deux conditions : les cellules germinales ne peuvent plus achever leur différenciation, d'où une azoospermie. Ce principe est celui de la contraception masculine hormonale, réversible à l'arrêt des injections.
\end{exoresolu}

\begin{exercice}[Pour t'entraîner]
Chez un homme adulte, on pratique une vasectomie (section des canaux déférents). Les dosages hormonaux réalisés un mois après l'opération montrent des taux de FSH, de LH et de testostérone normaux.
\begin{enumerate}
\item Ce résultat est-il compatible avec le schéma de régulation étudié ? Justifier.
\item Expliquer pourquoi la libido et les caractères sexuels secondaires sont conservés après vasectomie.
\end{enumerate}
\end{exercice}

\begin{corrige}[Exercice : vasectomie et équilibre hormonal]
\textbf{1.} Oui, ce résultat est parfaitement compatible. La vasectomie interrompt uniquement le \emph{transport} des spermatozoïdes (voie excrétrice) ; elle ne touche ni les cellules de Leydig, ni les cellules de Sertoli, ni les vaisseaux sanguins. La sécrétion de testostérone (et d'inhibine) se poursuit normalement, le rétrocontrôle négatif fonctionne toujours : les taux de FSH et de LH restent donc stables.

\textbf{2.} La libido et les caractères sexuels secondaires dépendent de la testostérone, qui est libérée dans le \textbf{sang} (voie endocrine) et non dans les canaux déférents (voie exocrine). Comme la sécrétion testiculaire est intacte, ces effets sont pleinement conservés. La vasectomie rend stérile sans modifier la masculinité : c'est un argument essentiel en éducation à la santé pour lever les craintes infondées.
\end{corrige}

\begin{aretenir}[L'essentiel de la section]
\begin{itemize}
\item Chez l'homme, les taux de FSH, LH et testostérone sont \textbf{constants} : fonctionnement continu, sans cycle (sécrétion pulsatile de GnRH/LH, testostérone lissée).
\item \textbf{LH} $\rightarrow$ cellules de Leydig $\rightarrow$ testostérone ; \textbf{FSH} $\rightarrow$ cellules de Sertoli $\rightarrow$ maintien de la spermatogenèse + sécrétion d'\textbf{inhibine}.
\item La testostérone assure la spermatogenèse, les caractères sexuels secondaires et la libido.
\item La testostérone exerce un \textbf{rétrocontrôle négatif global} sur l'hypothalamus (GnRH) et l'antéhypophyse (LH, FSH).
\item L'inhibine exerce un \textbf{rétrocontrôle négatif spécifique} sur la FSH (antéhypophyse).
\item Application : la contraception masculine hormonale repose sur le maintien artificiel de ce freinage par apport exogène de testostérone.
\end{itemize}
\end{aretenir}

\coursec{III. Régulation des hormones sexuelles chez la femme}

Dans la section précédente, nous avons vu que chez l'homme, la sécrétion de testostérone est \emph{continue} et soumise à un rétrocontrôle négatif simple : l'axe hypothalamo-hypophysaire fonctionne comme un thermostat réglé une fois pour toutes. Chez la femme, l'observation la plus banale — les règles qui reviennent environ tous les 28 jours — suggère un fonctionnement radicalement différent : un fonctionnement \cle{cyclique}. Comment l'organisme féminin orchestre-t-il ce cycle ? Quelles hormones en sont les chefs d'orchestre, et quel dialogue s'établit entre l'ovaire, l'hypophyse et l'utérus ? C'est toute la question de la régulation des hormones sexuelles chez la femme, que nous allons aborder comme des chercheurs : d'abord un constat expérimental (des dosages sanguins), puis une interprétation mécanistique, enfin des vérifications expérimentales.

\courssub{Constat expérimental : les dosages sanguins au cours du cycle}

\textbf{Problème scientifique.} Les taux sanguins des hormones sexuelles chez la femme sont-ils constants, comme chez l'homme, ou variables ?

\textbf{Protocole.} On prélève chaque jour, pendant 28 jours, un échantillon de sang chez une femme volontaire non enceinte et ne prenant aucun contraceptif. On dose dans le plasma quatre hormones : la \cle{FSH} et la \cle{LH} (hormones hypophysaires, en UI/L), les \cle{œstrogènes} (principalement l'estradiol, en ng/L ou pg/mL) et la \cle{progestérone} (hormone ovarienne, en µg/L ou ng/mL). En parallèle, on note les jours des règles et on suit l'ovaire par échographie.

\textbf{Résultats.} Les dosages, représentés sur la figure ci-dessous, montrent que les quatre hormones présentent des \emph{variations cycliques spectaculaires}, loin d'être constantes.

\begin{popfigure}
\begin{center}
\begin{tikzpicture}
\begin{axis}[
  width=\linewidth, height=8.2cm,
  xlabel={Jours du cycle},
  ylabel={Taux plasmatique (unités arbitraires)},
  xmin=0, xmax=28, ymin=0, ymax=10,
  xtick={0,7,14,21,28},
  legend style={at={(0.02,0.97)}, anchor=north west, font=\small},
  grid=both, grid style={popline!40},
  axis line style={popdark},
]
% FSH : petit pic vers J3 (recrutement), pic J14, début J28
\addplot[popblue, very thick, smooth, domain=0:28, samples=200]
  {1.2 + 0.8*exp(-((x-3)^2)/8) + 3.2*exp(-((x-14)^2)/2.5) + 0.6*exp(-((x-26)^2)/12)};
\addlegendentry{FSH}
% LH : grand pic à J14
\addplot[poporange, very thick, smooth, domain=0:28, samples=200]
  {1.0 + 9.0*exp(-((x-14)^2)/1.6)};
\addlegendentry{LH}
% Oestrogènes : pic pré-ovulatoire J13 + bosse lutéale J21
\addplot[poppink, very thick, smooth, domain=0:28, samples=200]
  {0.8 + 5.2*exp(-((x-13)^2)/3.5) + 2.6*exp(-((x-21)^2)/14)};
\addlegendentry{Œstrogènes}
% Progestérone : bosse lutéale J21
\addplot[popgreen, very thick, smooth, domain=0:28, samples=200]
  {0.3 + 6.5*exp(-((x-21)^2)/12)};
\addlegendentry{Progestérone}
% Repères
\draw[dashed, popdark] (axis cs:14,0) -- (axis cs:14,10);
\node[popdark, font=\small, rotate=90, anchor=south] at (axis cs:14.4,0.3) {ovulation};
\draw[dashed, popmuted] (axis cs:1,0) -- (axis cs:1,10);
\node[popmuted, font=\small, rotate=90, anchor=south] at (axis cs:1.4,0.3) {règles};
\end{axis}
\end{tikzpicture}
\end{center}
\legende{Évolution des taux plasmatiques de FSH, LH, œstrogènes et progestérone au cours d'un cycle de 28 jours (unités arbitraires pour permettre la comparaison sur un même graphe). On distingue deux phases : la phase folliculaire (J1 à J14) et la phase lutéale (J15 à J28), séparées par l'ovulation (J14).}
\end{popfigure}

\textbf{Interprétation.} La lecture de la courbe révèle quatre faits majeurs :
\begin{enumerate}
  \item En début de cycle (J1 à J5), la \textbf{FSH} est légèrement élevée : elle coïncide avec la croissance d'un groupe de follicules ovariens (recrutement folliculaire).
  \item Les \textbf{œstrogènes} augmentent progressivement jusqu'à un maximum vers J13, juste avant un événement brutal.
  \item Vers J13--J14, la \textbf{LH} (et dans une moindre mesure la FSH) présente un \emph{pic brutal et bref} ; l'échographie montre que c'est à ce moment qu'a lieu l'\cle{ovulation}.
  \item Après l'ovulation, la \textbf{progestérone} (et secondairement les œstrogènes) s'élèvent pour former un large plateau entre J18 et J24, puis chutent brutalement : les \textbf{règles} surviennent 2 à 3 jours plus tard, et un nouveau cycle commence.
\end{enumerate}

\begin{methode}[Exploiter la courbe du cycle ovarien]
Pour analyser rigoureusement un tel document au Baccalauréat :
\begin{enumerate}
  \item \textbf{Repérer les pics} : identifier les maxima de chaque courbe et noter leur date (pic d'œstrogènes vers J13, pic de LH vers J14, pic de progestérone vers J21).
  \item \textbf{Associer chaque pic à un événement} : le pic de LH précède immédiatement l'ovulation ; la chute de progestérone précède les règles.
  \item \textbf{Raisonner en causalité} : quand deux événements sont corrélés, se demander lequel provoque l'autre, et le vérifier par une expérience (ablation, injection).
  \item \textbf{Conclure en distinguant les phases} : phase folliculaire (dominée par FSH et œstrogènes) et phase lutéale (dominée par progestérone).
\end{enumerate}
\end{methode}

\courssub{La phase folliculaire (J1 à J14) : croissance folliculaire et œstrogènes}

Au premier jour du cycle (premier jour des règles), la FSH hypophysaire stimule la croissance d'une dizaine de \cle{follicules ovariens}, petites structures de l'ovaire contenant chacune un ovocyte en maturation. En général, un seul follicule — le \emph{follicule dominant} (ou follicule de De Graaf) — poursuit son développement jusqu'à maturité ; les autres régressent (atrésie).

Or les cellules de la \emph{granulosa} qui entourent l'ovocyte, sous l'action combinée de la FSH (qui induit l'aromatase) et de la LH (qui stimule les cellules de la thèque pour fournir les précurseurs androgènes), synthétisent et libèrent des \textbf{œstrogènes} (essentiellement l'estradiol). Ainsi, plus le follicule grossit, plus il sécrète d'œstrogènes : c'est exactement ce que montre la montée progressive de la courbe rose entre J5 et J13. Les œstrogènes agissent alors sur l'utérus : ils provoquent la \emph{prolifération de l'endomètre} (la muqueuse utérine s'épaissit et se vascularise).

\begin{exemplebox}[Un pic hormonal massif et bref]
Le pic de LH pré-ovulatoire peut atteindre \textbf{3 à 5 fois le taux de base en 24 à 36 heures} : par exemple, le taux plasmatique de LH peut passer d'environ \SI{10}{UI/L} à \SI{40}{UI/L} entre J13 et J14, avant de retomber presque aussi vite. C'est cette brusque décharge qui déclenche l'ovulation. C'est d'ailleurs ce pic que détectent les \emph{tests d'ovulation} vendus en pharmacie, qui révèlent la LH dans les urines.
\end{exemplebox}

\begin{savaistu}[Théorie des deux cellules, deux gonadotrophines]
La synthèse d'œstrogènes par le follicule nécessite une coopération : les \textbf{cellules de la thèque} (sous l'action de la LH) produisent des androgènes (androstènedione) qui diffusent vers les \textbf{cellules de la granulosa} ; celles-ci (sous l'action de la FSH) expriment l'aromatase qui convertit les androgènes en œstrogènes (estradiol). C'est un exemple classique de compartimentation métabolique.
\end{savaistu}

\courssub{Le pic de LH et l'ovulation (vers J14)}

\begin{definition}[Ovulation]
L'\cle{ovulation} est la \textbf{libération de l'ovocyte} (plus précisément de l'ovocyte II bloqué en métaphase II de la méiose, entouré de ses cellules folliculaires : le cumulus oophorus) par le follicule ovarien mûr, qui se rompt à la surface de l'ovaire. Elle est \textbf{déclenchée par le pic de LH}, vers le 14\textsuperscript{e} jour d'un cycle de 28 jours. L'ovocyte libéré est capté par le pavillon de la trompe utérine, où la fécondation pourra avoir lieu.
\end{definition}

\textbf{Pourquoi le pic de LH survient-il ?} C'est le point le plus délicat du chapitre. En fin de phase folliculaire, le taux d'œstrogènes dépasse un seuil élevé (environ \SI{200}{pg/mL}) et se maintient pendant 24 à 36 heures. Or, à cette concentration élevée et soutenue, les œstrogènes exercent sur l'axe hypothalamo-hypophysaire un \cle{rétrocontrôle positif} : au lieu de freiner la sécrétion de LH, ils la \emph{stimulent} brutalement (en augmentant la fréquence des pulsations de GnRH et la sensibilité de l'hypophyse à la GnRH). C'est l'un des rares exemples de rétroaction positive en physiologie humaine. La conséquence est la décharge de LH, qui provoque la rupture du follicule : l'ovulation.

\begin{attention}[Erreur fréquente : la date de l'ovulation]
Beaucoup d'élèves (et d'adultes !) situent l'ovulation « pendant les règles » ou « juste après ». C'est faux. L'ovulation a lieu \textbf{environ 14 jours avant les règles suivantes}, car la phase lutéale a une durée relativement constante de 14 jours. Sur un cycle de 28 jours, elle tombe donc vers J14 ; sur un cycle de 32 jours, vers J18. Les règles, elles, marquent le \emph{début} du cycle (J1), pas son milieu.
\end{attention}

\courssub{La phase lutéale (J15 à J28) : le corps jaune et la progestérone}

Après l'ovulation, le follicule rompu ne disparaît pas : sous l'action de la LH (lutéinisation), ses cellules restantes (granulosa et thèque) se transforment en une structure jaunâtre et très vascularisée, le \cle{corps jaune} (\emph{corpus luteum}). Le corps jaune sécrète massivement de la \textbf{progestérone} et, dans une moindre mesure, des œstrogènes : c'est le large plateau vert de la courbe entre J18 et J24.

La progestérone agit sur l'endomètre déjà épaissi par les œstrogènes : elle le rend \emph{sécrétoire} — ses glandes s'enroulent et libèrent un mucus nutritif riche en glycogène, ses vaisseaux se développent en spirales. L'endomètre est alors \textbf{prêt à la nidation}, c'est-à-dire à accueillir un embryon (stade blastocyste) vers J20--J21.

En l'absence de fécondation, le corps jaune régresse spontanément après 12 à 14 jours (il devient le \emph{corps blanc} ou corps albicans). Les taux de progestérone et d'œstrogènes s'effondrent : privée de son soutien hormonal, la couche fonctionnelle de l'endomètre se nécrose et s'élimine avec un saignement : ce sont les \textbf{règles} (ou menstrues), qui durent 3 à 5 jours. Le cycle recommence.

\begin{savaistu}[Pourquoi « corps jaune » ?]
Le corps jaune doit son nom à sa couleur jaune orangée, visible à l'œil nu sur une coupe d'ovaire. Cette couleur est due à l'accumulation de \textbf{pigments caroténoïdes} (dont le bêta-carotène, le même pigment que la carotte !) dans ses cellules, riches en lipides. Si une grossesse débute, le corps jaune est maintenu par l'hormone \emph{hCG} (chorionique gonadotrophine) sécrétée par le trophoblaste de l'embryon et continue à produire la progestérone indispensable au maintien de l'endomètre pendant le premier trimestre, jusqu'à la prise de relais par le placenta.
\end{savaistu}

\begin{popfigure}
\begin{center}
\begin{tikzpicture}[font=\small]
% Bandeau des jours
\draw[->, thick, popdark] (0,0) -- (14.6,0) node[right]{jours};
\foreach \x/\j in {0/1, 3.5/7, 7/14, 10.5/21, 14/28}{
  \draw[popdark] (\x,0.08) -- (\x,-0.08);
  \node[below] at (\x,-0.1) {\j};
}
% Phases
\draw[popblue, very thick] (0,0.35) -- (7,0.35);
\node[popblue, above] at (3.5,0.4) {Phase folliculaire};
\draw[popgreen, very thick] (7,0.35) -- (14,0.35);
\node[popgreen, above] at (10.5,0.4) {Phase lutéale};
\fill[poporange] (7,0.35) circle (0.09);
\node[poporange, above right=0.02cm and -0.6cm] at (7,0.55) {\textbf{ovulation}};
% Ovaire : follicules
\node[popdark, anchor=west] at (-0.2,1.6) {\textbf{Ovaire}};
\foreach \x/\r in {0.7/0.16, 1.8/0.24, 3.0/0.33, 4.4/0.44, 5.8/0.56}{
  \draw[poppink, fill=poppinkL] (\x,1.6) circle (\r);
  \fill[popdark] (\x,1.6) circle (0.05);
}
\draw[popgreen, fill=popgreenL] (8.2,1.6) circle (0.55);
\node[popgreen] at (8.2,1.6) {\tiny corps};
\node[popgreen] at (8.2,1.28) {\tiny jaune};
\draw[popgreen, fill=popgreenL] (9.9,1.6) circle (0.42);
\draw[popmuted, fill=popline!50] (11.4,1.6) circle (0.28);
\node[popmuted, font=\tiny] at (11.4,2.05) {régression};
\draw[popmuted, fill=popline!30] (12.8,1.6) circle (0.16);
% Utérus : endomètre
\node[popdark, anchor=west] at (-0.2,3.0) {\textbf{Utérus}};
\foreach \x/\h in {0.7/0.10, 2.2/0.22, 3.8/0.38, 5.4/0.52, 7/0.60, 8.6/0.78, 10.2/0.92, 11.8/0.95, 13.2/0.30}{
  \draw[poporange, fill=poporangeL] (\x-0.5,2.6) rectangle (\x+0.5,2.6+\h);
}
\node[poporange, font=\tiny] at (3,3.6) {prolifération (œstrogènes)};
\node[popgreen, font=\tiny] at (10,3.85) {endomètre sécrétoire (progestérone)};
\draw[->, popmuted, thick] (13.2,2.5) -- (13.2,2.1) node[below, font=\tiny]{règles};
% Hormones dominantes
\node[popdark, anchor=west] at (-0.2,4.6) {\textbf{Hormones}};
\node[poppink] at (3.5,4.6) {FSH puis œstrogènes $\nearrow$};
\node[popgreen] at (10.5,4.6) {progestérone $\nearrow$ puis chute};
\end{tikzpicture}
\end{center}
\legende{Synchronisation des événements au cours d'un cycle de 28 jours : croissance puis rupture du follicule, formation et régression du corps jaune (ovaire) ; prolifération puis état sécrétoire puis desquamation de l'endomètre (utérus), sous le contrôle des hormones ovariennes.}
\end{popfigure}

\courssub{Les rétrocontrôles : démonstration expérimentale}

Les corrélations observées sur la courbe suggèrent que les hormones ovariennes contrôlent les sécrétions hypophysaires. Pour le \emph{prouver}, réalisons l'expérience classique.

\begin{experience}[Ovariectomie et injections d'œstrogènes (modèle animal)]
\textbf{Matériel et protocole.} On utilise des rates (ou des singes) femelles, dont on suit les cycles et dont on dose FSH et LH dans le sang.
\begin{itemize}
  \item \textbf{Lot 1 (témoin)} : animaux intacts. Les cycles sont réguliers, FSH et LH oscillent à des taux modérés.
  \item \textbf{Lot 2} : \textbf{ovariectomie} (ablation chirurgicale des deux ovaires). \emph{Observation} : les cycles s'arrêtent, et les taux de FSH et de LH \textbf{augmentent fortement et durablement} (perte du frein ovarien).
  \item \textbf{Lot 3} : animaux ovariectomisés recevant des \textbf{injections d'œstrogènes} à doses physiologiques (modérées). \emph{Observation} : les taux de FSH et de LH \textbf{diminuent} et rejoignent des valeurs basses (restauration du frein).
  \item \textbf{Lot 4} : injection d'une dose \emph{élevée} (supra-physiologique) d'œstrogènes à des animaux intacts en début de cycle. \emph{Observation} : une décharge brutale de LH est provoquée (mimant le pic pré-ovulatoire).
\end{itemize}
\textbf{Interprétation.} Sans ovaires, plus rien ne freine l'hypophyse : FSH et LH s'élèvent (levée d'inhibition). Les œstrogènes à dose modérée restaurent ce frein : c'est un \textbf{rétrocontrôle négatif}. Mais à dose élevée et soutenue, les œstrogènes déclenchent au contraire une décharge de LH : c'est le \textbf{rétrocontrôle positif}, responsable du pic ovulatoire.\\
\textbf{Conclusion.} Les hormones ovariennes (œstrogènes et progestérone) contrôlent l'activité de l'axe hypothalamo-hypophysaire selon leur concentration et leur durée d'exposition : l'ovaire et l'hypophyse dialoguent en boucle fermée.
\end{experience}

On résume ce double rétrocontrôle ainsi :
\begin{itemize}
  \item \textbf{Fin de phase folliculaire} : pic d'œstrogènes (taux élevé soutenu) $\Rightarrow$ \textbf{rétrocontrôle positif} sur l'hypophyse (et l'hypothalamus) $\Rightarrow$ pic de LH $\Rightarrow$ ovulation.
  \item \textbf{Phase lutéale} : progestérone (majoritaire) + œstrogènes $\Rightarrow$ \textbf{rétrocontrôle négatif} $\Rightarrow$ FSH et LH restent basses $\Rightarrow$ aucun nouveau follicule ne se développe pendant cette phase (blocage de la folliculogenèse).
\end{itemize}

\begin{popfigure}
\begin{center}
\begin{tikzpicture}[font=\small, node distance=1.1cm,
  glande/.style={ellipse, draw=popdark, fill=popblueL, minimum width=2.8cm, minimum height=1.1cm, align=center},
  cible/.style={ellipse, draw=popdark, fill=poppinkL, minimum width=2.6cm, minimum height=1.1cm, align=center},
  flecheplus/.style={->, very thick, popgreen},
  flechemoins/.style={->, very thick, popink},
]
\node[glande, align=center] (hypo) at (0,4.5){Hypothalamus\\ \footnotesize GnRH};
\node[glande, fill=popgoldL, align=center] (hyp) at (0,2.2){Hypophyse antérieure\\ \footnotesize FSH, LH};
\node[cible, align=center] (ovaire) at (0,-0.3){Ovaire\\ \footnotesize Follicule / Corps jaune};
\draw[->, very thick, popblue] (hypo) -- (hyp) node[midway, right, popdark, font=\footnotesize]{stimulation};
\draw[->, very thick, popblue] (hyp) -- (ovaire) node[midway, right, popdark, font=\footnotesize]{stimulation};
% Rétrocontrôle négatif (phase lutéale) - gauche
\draw[flechemoins] (ovaire.west) .. controls (-4.5,-0.3) and (-4.5,2.2) .. (hyp.west)
  node[midway, left, align=center, popink, font=\footnotesize]{rétrocontrôle \textbf{négatif}\\ progestérone + œstrogènes\\ (phase lutéale)};
% Rétrocontrôle positif (fin phase folliculaire) - droite
\draw[flecheplus] (ovaire.east) .. controls (4.5,-0.3) and (4.5,2.2) .. (hyp.east)
  node[midway, right, align=center, popgreen, font=\footnotesize]{rétrocontrôle \textbf{positif}\\ pic d'œstrogènes\\ (fin phase folliculaire) $\Rightarrow$ pic de LH};
\end{tikzpicture}
\end{center}
\legende{Boucle de régulation de l'axe hypothalamo-hypophyso-ovarien chez la femme. Flèche verte : rétrocontrôle positif des œstrogènes à taux élevé (déclenche le pic de LH ovulatoire). Flèche foncée : rétrocontrôle négatif de la progestérone et des œstrogènes pendant la phase lutéale (maintient FSH et LH à un taux bas).}
\end{popfigure}

\begin{aretenir}[L'essentiel du cycle ovarien]
\begin{itemize}
  \item \textbf{Phase folliculaire (J1--J14)} : la FSH recrute et fait croître les follicules ; le follicule dominant sécrète des œstrogènes (prolifération de l'endomètre). Rétrocontrôle négatif des œstrogènes sur FSH (sélection du dominant).
  \item \textbf{J13--J14} : le pic d'œstrogènes provoque, par rétrocontrôle positif, le pic de LH, qui déclenche l'\textbf{ovulation} et la lutéinisation.
  \item \textbf{Phase lutéale (J15--J28)} : le corps jaune sécrète progestérone et œstrogènes (endomètre sécrétoire, prêt pour la nidation) et exerce un rétrocontrôle négatif sur FSH/LH.
  \item \textbf{Sans fécondation} : chute des hormones ovariennes (régression du corps jaune) $\Rightarrow$ règles (desquamation endomètre) $\Rightarrow$ levée du frein sur FSH $\Rightarrow$ nouveau cycle.
\end{itemize}

\begin{center}
\begin{tabularx}{\linewidth}{|l|X|X|}\hline
\rowcolor{popblueL} \textbf{Paramètre} & \textbf{Phase folliculaire} & \textbf{Phase lutéale} \\ \hline
\textbf{Durée} & Variable (environ 14 j) & Fixe (14 j) \\ \hline
\textbf{Ovaire} & Croissance folliculaire $\rightarrow$ Ovulation & Corps jaune (formation $\rightarrow$ activité $\rightarrow$ régression) \\ \hline
\textbf{Hormone ovarienne dominante} & Œstrogènes (estradiol) & Progestérone (+ œstrogènes) \\ \hline
\textbf{Endomètre} & Prolifération (épaississement) & Sécrétion (maturation, vascularisation) \\ \hline
\textbf{Rétrocontrôle ovarien} & Négatif (début) $\rightarrow$ Positif (fin, pic œstrogènes) & Négatif (progestérone + œstrogènes) \\ \hline
\textbf{FSH / LH} & FSH $\nearrow$ puis $\searrow$ ; LH bas puis pic brutal & FSH et LH bas (inhibés) \\ \hline
\end{tabularx}
\end{center}
\end{aretenir}

\courssub{Application : le principe de la pilule contraceptive (complément Bac)}

\begin{exoresolu}[Comment la pilule empêche-t-elle l'ovulation ?]
\textbf{Énoncé.} La pilule œstroprogestative (combinée), prise quotidiennement pendant 21 jours, apporte des œstrogènes (éthinylestradiol) et un progestatif (dérivé de la progestérone) à taux constant. À l'aide des connaissances sur la régulation du cycle, expliquer son mode d'action contraceptif.
\tcblower
\textbf{Solution détaillée.} Les hormones apportées par la pilule maintiennent en permanence des taux d'œstrogènes et de progestatif comparables à ceux de la \emph{phase lutéale}. Or, à ces concentrations, ces hormones exercent un \textbf{rétrocontrôle négatif} continu sur l'axe hypothalamo-hypophysaire : les sécrétions endogènes de FSH et de LH sont bloquées à un niveau bas (suppression de l'axe). Deux conséquences majeures :
\begin{enumerate}
  \item Sans FSH suffisante, \textbf{aucun follicule ne se développe} (pas de recrutement, pas de follicule dominant) : il n'y a donc ni follicule mûr, ni pic d'œstrogènes endogènes.
  \item Sans pic d'œstrogènes endogènes, il n'y a pas de rétrocontrôle positif, donc \textbf{pas de pic de LH} et par conséquent \textbf{pas d'ovulation}.
\end{enumerate}
L'arrêt de la pilule pendant 7 jours (ou prise de comprimés placebo) provoque une chute hormonale artificielle, d'où des saignements de privation (hémorragie de sevrage) ressemblant à des règles. \textbf{Conclusion} : la pilule détourne le rétrocontrôle négatif naturel pour « faire croire » à l'axe hypothalamo-hypophysaire que la phase lutéale est permanente, bloquant ainsi l'ovulation. (Le progestatif épaissit aussi la glaire cervicale, ce qui freine le passage des spermatozoïdes : double sécurité.)
\end{exoresolu}

\begin{attention}[Pièges de rédaction au Bac]
\begin{itemize}
  \item Ne pas écrire que « la LH fabrique le corps jaune » de façon vague : préciser que la LH \emph{déclenche l'ovulation} puis \emph{entretient la lutéinisation} (transformation du follicule rompu en corps jaune).
  \item Ne pas confondre les deux rétrocontrôles : le rétrocontrôle \textbf{positif} n'existe qu'en \emph{fin de phase folliculaire} (œstrogènes à taux élevé et soutenu) ; le rétrocontrôle \textbf{négatif} domine en \emph{phase lutéale} (progestérone + œstrogènes) et en début de phase folliculaire.
  \item Dans l'expérience d'ovariectomie, la hausse de FSH/LH s'explique par la \emph{disparition du frein} ovarien (levée d'inhibition), pas par une « stimulation » active de l'hypophyse par l'absence d'ovaire.
  \item Préciser que l'ovocyte libéré est en \textbf{métaphase II} (blocage méiotique) et achèvera sa méiose \textbf{seulement en cas de fécondation}.
\end{itemize}
\end{attention}

Ainsi, le cycle féminin apparaît comme un dialogue permanent entre l'hypothalamus, l'hypophyse, l'ovaire et l'utérus, où chaque hormone est à la fois un message et une réponse. Quand ce dialogue est rompu — ovaires paresseux (insuffisance ovarienne), trompes obstruées, sécrétions hypophysaires insuffisantes (hyperprolactinémie, anorexie), endométriose —, la reproduction devient impossible : c'est le problème de l'\emph{infertilité}, que nous aborderons dans la section suivante.

\coursec{L'infertilité : causes possibles chez l'homme et chez la femme}

Dans les sections précédentes, nous avons vu que la fertilité repose sur un équilibre hormonal délicat : chez l'homme, l'axe hypothalamo-hypophysaire contrôle en permanence la spermatogenèse via la FSH, la LH et la testostérone ; chez la femme, le dialogue ovaires--hypophyse--hypothalamus orchestre le cycle ovarien et l'ovulation. Mais que se passe-t-il lorsque cet équilibre est rompu, ou lorsqu'un obstacle anatomique empêche la rencontre des gamètes ? C'est toute la question de l'\cle{infertilité}, un problème de santé publique majeur : au Cameroun comme ailleurs, de nombreux couples consultent dans les hôpitaux de district et les cliniques spécialisées (Yaoundé, Douala, Bafoussam) après des mois ou des années d'attente d'un enfant. Cette section te donne les clés scientifiques pour comprendre les causes de l'infertilité et pour construire des messages de sensibilisation efficaces.

\courssub{Définitions et données épidémiologiques}

\begin{definition}[Infertilité et stérilité]
L'\cle{infertilité} d'un couple est l'\textbf{absence de conception après 12 à 24 mois de rapports sexuels réguliers sans contraception} (définition de l'Organisation mondiale de la santé, OMS). On parle d'infertilité \emph{primaire} si le couple n'a jamais eu d'enfant, et d'infertilité \emph{secondaire} si le couple a déjà eu un enfant mais n'arrive plus à concevoir. La \cle{stérilité}, en revanche, désigne une \textbf{incapacité définitive} à procréer (par exemple après ablation des deux ovaires ou des deux testicules). L'infertilité est une difficulté à concevoir, souvent traitable ; la stérilité est irréversible.
\end{definition}

\begin{attention}[Une idée reçue à combattre]
Dans beaucoup de familles, lorsqu'un couple n'a pas d'enfant, on accuse spontanément la femme. C'est une \textbf{erreur scientifique grave} : les études épidémiologiques montrent que les causes sont réparties en trois tiers à peu près égaux. L'homme est en cause, seul ou associé à la femme, dans \textbf{environ un cas sur deux}. Le bilan d'infertilité doit donc \emph{toujours} concerner les deux partenaires, et la consultation doit se faire \textbf{en couple}.
\end{attention}

\begin{popfigure}
\begin{center}
\begin{tikzpicture}[scale=1]
  % Diagramme circulaire : 1/3 homme, 1/3 femme, 1/3 mixte/inexpliquée
  \fill[popblue!70] (0,0) -- (90:2.4) arc (90:210:2.4) -- cycle;
  \fill[poppink!70] (0,0) -- (210:2.4) arc (210:330:2.4) -- cycle;
  \fill[popgold!80] (0,0) -- (330:2.4) arc (330:450:2.4) -- cycle;
  \draw[popdark, thick] (0,0) circle (2.4);
  \node[white, font=\small\bfseries] at (150:1.5) {Homme};
  \node[white, font=\small\bfseries] at (150:1.1) {$\approx 1/3$};
  \node[white, font=\small\bfseries] at (270:1.5) {Femme};
  \node[white, font=\small\bfseries] at (270:1.1) {$\approx 1/3$};
  \node[popdark, font=\small\bfseries] at (30:1.55) {Mixte ou};
  \node[popdark, font=\small\bfseries] at (30:1.15) {inexpliquée};
  \node[popdark, font=\small\bfseries] at (30:0.8) {$\approx 1/3$};
\end{tikzpicture}
\end{center}
\legende{Répartition approximative des causes d'infertilité dans un couple : environ un tiers des cas relève d'une cause masculine, un tiers d'une cause féminine, et un tiers d'une cause mixte (les deux partenaires) ou inexpliquée. Cette répartition justifie l'examen systématique des deux partenaires.}
\end{popfigure}

\begin{aretenir}[À retenir]
L'infertilité est un problème \textbf{de couple}, pas un problème de femme. La répartition $1/3$ -- $1/3$ -- $1/3$ est un savoir admis qu'il faut savoir citer au Baccalauréat et utiliser dans tout outil de sensibilisation.
\end{aretenir}

\courssub{Les causes masculines de l'infertilité}

\textbf{Le spermogramme, examen de première intention}\par

Comment savoir si l'homme est en cause ? L'examen de base est le spermogramme.

\begin{definition}[Spermogramme]
Le \cle{spermogramme} est l'examen de laboratoire du sperme, recueilli après 3 à 5 jours d'abstinence. Il évalue trois paramètres essentiels : le \textbf{nombre} de spermatozoïdes (concentration), leur \textbf{mobilité} et leur \textbf{morphologie} (forme normale ou anormale). Selon les normes OMS (5\textsuperscript{e} édition, valeurs de référence du 5\textsuperscript{e} centile), un sperme fertile contient au moins \SI{15e6}{mL^{-1}} spermatozoïdes, avec au moins \SI{32}{\percent} de spermatozoïdes à mobilité progressive et \SI{4}{\percent} de formes normales.
\end{definition}

\begin{exemplebox}[Exemple chiffré : lecture d'un spermogramme]
Un homme de 32 ans consulte au centre hospitalier d'Essos (Yaoundé). Son spermogramme donne : volume \SI{2,8}{mL} ; concentration \SI{8e6}{mL^{-1}} ; mobilité progressive \SI{25}{\percent}. La concentration étant inférieure au seuil OMS de \SI{15e6}{mL^{-1}}, on conclut à une \textbf{oligospermie}. Une échographie révèle un \textbf{varicocèle} (dilatation des veines du cordon spermatique), cause fréquente et opérable.
\end{exemplebox}

\textbf{Les principales anomalies du sperme}\par

\begin{definition}[Anomalies quantitatives et qualitatives du sperme]
\begin{itemize}
  \item \cle{Oligospermie} : concentration de spermatozoïdes inférieure à la norme ($< \SI{15e6}{mL^{-1}}$).
  \item \cle{Azoospermie} : absence totale de spermatozoïdes dans le sperme (par défaut de production testiculaire ou par obstruction des voies excrétrices).
  \item \cle{Asthénospermie} (ou asthénzoospermie) : mobilité progressive insuffisante des spermatozoïdes ($< \SI{32}{\percent}$), qui ne peuvent pas remonter les voies génitales femelles jusqu'à l'ovocyte.
  \item \cle{Tératospermie} (ou tératozoospermie) : proportion excessive de spermatozoïdes de forme anormale ($< \SI{4}{\percent}$ de formes normales).
\end{itemize}
\end{definition}

\textbf{Les autres causes masculines}\par

\begin{itemize}
  \item \textbf{Le varicocèle} : dilatation variqueuse des veines du cordon spermatique. Elle élève la température du testicule (or la spermatogenèse exige une température d'environ \SI{34}{\celsius}, inférieure à la température corporelle de \SI{37}{\celsius}) et altère la production des spermatozoïdes. C'est la cause masculine la plus fréquente et elle se traite chirurgicalement.
  \item \textbf{Les infections sexuellement transmissibles (IST)} : les infections à \emph{Chlamydia trachomatis} et la gonococcie (blennorragie) non traitées provoquent des épididymites et des obstructions des canaux déférents, empêchant l'émission des spermatozoïdes.
  \item \textbf{Les troubles hormonaux} : un déficit en FSH ou en LH (panhypopituitarisme) prive les testicules de leur stimulation ; la spermatogenèse s'arrête. On retrouve ici le rôle de l'axe hypothalamo-hypophysaire étudié dans la section II.
  \item \textbf{La cryptorchidie} : non-descente d'un ou des deux testicules dans les bourses pendant la vie fœtale ou néonatale. Le testicule resté dans l'abdomen, exposé à \SI{37}{\celsius}, ne produit pas de spermatozoïdes ; non corrigée tôt (orchidopexie avant 2 ans idéalement), elle entraîne une infertilité.
  \item \textbf{Les facteurs toxiques} : alcoolisme chronique, tabagisme, exposition à la chaleur excessive (bains très chauds, postes de travail près de fours), certains médicaments et pesticides.
\end{itemize}

\courssub{Les causes féminines de l'infertilité}

\textbf{Les troubles de l'ovulation}\par

Sans ovulation, pas de fécondation possible. L'\cle{anovulation} (absence d'ovulation) résulte souvent d'un dérèglement de l'axe hypothalamo-hypophyso-ovarien : déficit en FSH ou en LH, hyperprolactinémie (excès de prolactine), dysthyroïdie. La cause la plus fréquente est le \textbf{syndrome des ovaires polykystiques (SOPK)} : les follicules démarrent leur croissance mais ne parviennent pas à maturité, l'ovulation n'a pas lieu ; on observe des cycles irréguliers (oligoménorrhée ou aménorrhée), un excès d'androgènes (hirsutisme, acné) et, à l'échographie, de nombreux petits follicules périphériques (aspect « en collier de perles »).

\textbf{L'obstruction des trompes de Fallope}\par

Même si l'ovulation est normale, l'ovocyte doit rencontrer les spermatozoïdes dans le tiers externe de la trompe (ampoule). Si la trompe est \textbf{bouchée}, la fécondation est impossible ; si elle est partiellement altérée (perte des cils, paroi rigide), l'œuf peut s'implanter dans la trompe : c'est la \textbf{grossesse extra-utérine}, urgence chirurgicale vitale. La cause dominante de l'obstruction tubaire est l'\textbf{infection génitale haute} (salpingite) consécutive à une IST (\emph{Chlamydia trachomatis}, gonocoque) mal traitée ou non traitée. Les avortements pratiqués dans de mauvaises conditions d'asepsie sont aussi une cause majeure d'infections et de séquelles tubaires.

\begin{popfigure}
\begin{center}
\begin{tikzpicture}[scale=1.05, every node/.style={font=\small}]
  % Ovaires
  \draw[fill=poppink!60, draw=popdark] (-3.6,0.6) ellipse (0.55 and 0.8);
  \draw[fill=poppink!60, draw=popdark] (3.6,0.6) ellipse (0.55 and 0.8);
  % Utérus
  \draw[fill=poppurple!25, draw=popdark, thick] (-1.1,-1.6) .. controls (-1.5,-0.4) and (-1.3,0.5) .. (-0.9,0.9)
        -- (0.9,0.9) .. controls (1.3,0.5) and (1.5,-0.4) .. (1.1,-1.6)
        .. controls (0.5,-2.0) and (-0.5,-2.0) .. (-1.1,-1.6);
  \draw[draw=popdark, thick] (-0.35,-1.9) rectangle (0.35,-2.9);
  % Trompe gauche : normale
  \draw[draw=popblue, very thick] (-0.9,0.9) .. controls (-1.6,1.5) and (-2.6,1.6) .. (-3.1,1.1);
  \draw[draw=popblue, very thick] (-3.1,1.1) .. controls (-3.35,0.85) .. (-3.2,0.55);
  % Trompe droite : obstruée
  \draw[draw=poporange, very thick] (0.9,0.9) .. controls (1.6,1.5) and (2.3,1.55) .. (2.7,1.3);
  \draw[fill=poporange, draw=popdark] (2.75,1.28) circle (0.16);
  \draw[draw=poporange, very thick, dashed] (2.91,1.28) .. controls (3.2,1.05) .. (3.2,0.6);
  % Légendes fléchées
  \node[popdark, anchor=east] at (-4.6,0.6) {\textbf{1. Ovaire}};
  \draw[->, popdark] (-4.55,0.6) -- (-4.15,0.6);
  \node[popblue, anchor=south west] at (-2.9,1.75) {\textbf{2. Trompe perméable}};
  \draw[->, popblue] (-2.2,1.72) -- (-2.2,1.45);
  \node[popdark, anchor=west] at (4.3,0.6) {\textbf{4. Utérus}};
  \draw[->, popdark] (4.3,0.55) -- (1.6,-0.2);
  \node[poporange, anchor=west] at (3.15,1.95) {\textbf{3. Obstruction tubaire}};
  \draw[->, poporange] (3.6,1.9) -- (2.8,1.4);
  \node[popdark, anchor=north] at (0,-3.0) {\textbf{5. Col de l'utérus}};
\end{tikzpicture}
\end{center}
\legende{Schéma de l'appareil génital féminin (vue de face). 1 : ovaire (production des ovocytes). 2 : trompe de Fallope gauche perméable, siège normal de la fécondation. 3 : obstruction de la trompe droite (point orange) consécutive à une salpingite : les spermatozoïdes ne peuvent plus atteindre l'ovocyte. 4 : utérus. 5 : col de l'utérus. L'obstruction tubaire se diagnostique par hystérosalpingographie.}
\end{popfigure}

\textbf{Les anomalies de l'utérus et l'endométriose}\par

\begin{itemize}
  \item \textbf{Les fibromes (myomes)} : tumeurs bénignes du muscle utérin, très fréquentes chez la femme africaine. Selon leur taille et leur siège (sous-muqueux, intramural, sous-séreux), ils peuvent déformer la cavité utérine et empêcher la nidation.
  \item \textbf{Les synéchies utérines} (syndrome d'Asherman) : brides cicatricielles dans la cavité utérine, consécutives à des curetages traumatiques (souvent après un avortement mal encadré) ou des infections post-partum ; elles réduisent ou oblitèrent la cavité utérine.
  \item \textbf{L'endométriose} : présence de muqueuse utérine (endomètre) en dehors de l'utérus (ovaires : kystes « chocolat », trompes, péritoine, ligaments utéro-sacrés). Ces foyers saignent à chaque cycle, provoquant douleurs pelviennes chroniques (dysménorrhée, dyspareunie), adhérences et altération de la fonction tubaire et ovarienne.
  \item \textbf{L'âge} : la réserve ovarienne diminue avec l'âge ; la qualité des ovocytes se dégrade nettement après 35 ans (augmentation des anomalies chromosomiques), et la fertilité chute fortement après 38 ans. Le recul de l'âge de la première grossesse est un facteur d'infertilité croissant en milieu urbain.
\end{itemize}

\courssub{Causes communes et facteurs favorisants}

Certaines causes concernent les deux sexes ou le couple dans son ensemble :

\begin{itemize}
  \item les \textbf{IST non traitées} (\emph{Chlamydia trachomatis}, gonococcie) : cause évitable n$^{\circ}$1 d'infertilité acquise, chez l'homme (épididymite, obstruction déférentielle) comme chez la femme (salpingite, hydrosalpinx) ;
  \item les \textbf{avortements pratiqués dans de mauvaises conditions}, sources d'infections pelviennes et de synéchies utérines ;
  \item le \textbf{tabac} et l'\textbf{alcool}, qui altèrent la spermatogenèse (fragmentation ADN, morphologie) et la qualité ovocytaire (réserve ovarienne) ;
  \item le \textbf{stress} chronique, qui peut perturber l'axe hypothalamo-hypophysaire des deux partenaires (hypercortisolémie inhibitrice sur la GnRH) ;
  \item l'\textbf{âge tardif} du projet de grossesse.
\end{itemize}

\courssub{Le bilan d'infertilité : une démarche de couple}

\begin{methode}[Démarche du bilan d'infertilité]
\begin{enumerate}
  \item \textbf{Interrogatoire du couple} : durée de l'infertilité, antécédents d'IST, de chirurgie pelvienne ou testiculaire, mode de vie (tabac, alcool, exposition professionnelle), régularité des cycles, fréquence des rapports.
  \item \textbf{Chez l'homme} : \textbf{spermogramme} (examen de première intention, simple, non invasif et peu coûteux), réalisé après 3 à 5 jours d'abstinence. Complété si besoin par des dosages hormonaux (FSH, LH, testostérone, prolactine) et une échographie des bourses (recherche de varicocèle, de kyste de l'épididyme, d'atrophie testiculaire).
  \item \textbf{Chez la femme} : \textbf{courbe de température} (la température basale au réveil s'élève d'environ \SI{0,4}{\celsius} à \SI{0,5}{\celsius} après l'ovulation sous l'effet thermogène de la progestérone : une courbe biphasique témoigne d'une ovulation) ; \textbf{dosages hormonaux} (FSH, LH, œstradiol en début de cycle J2-J3 ; progestérone en milieu de phase lutéale J21 pour confirmer l'ovulation ; AMH pour évaluer la réserve ovarienne) ; \textbf{échographie pelvienne} (ovaires, follicules en cours de cycle, utérus, fibromes, kystes d'endométriose) ; \textbf{hystérosalpingographie (HSG)} : radiographie de l'utérus et des trompes après injection d'un produit de contraste iodé, qui vérifie la \textbf{perméabilité tubaire} et la morphologie de la cavité utérine.
  \item \textbf{Synthèse} : identification de la ou des causes, puis proposition d'un traitement étiologique (chirurgie du varicocèle, induction de l'ovulation, chirurgie tubaire) ou orientation vers une procréation médicalement assistée (section V).
\end{enumerate}
\end{methode}

\begin{center}
\begin{tabularx}{\linewidth}{|l|X|X|}
\hline
\rowcolor{popblueL} \textbf{Cause} & \textbf{Chez l'homme} & \textbf{Chez la femme} \\
\hline
Production des gamètes & Oligospermie, azoospermie, asthénospermie, tératospermie (varicocèle, cryptorchidie, déficit FSH/LH, toxiques) & Anovulation (SOPK, hyperprolactinémie, insuffisance ovarienne), déficit hypophysaire \\
\hline
Transport des gamètes & Obstruction des canaux déférents (IST, chirurgie) & Obstruction des trompes (salpingite post-IST), endométriose (adhérences) \\
\hline
Accueil de l'œuf & --- & Fibromes sous-muqueux, synéchies utérines, malformations utérines \\
\hline
Facteurs généraux & Alcool, tabac, chaleur, stress, âge & Tabac, alcool, stress, âge $> 35$ ans \\
\hline
\rowcolor{popgoldL} \textbf{Examen clé de diagnostic} & Spermogramme, dosages hormonaux, échographie scrotale & Courbe de température, dosages hormonaux, échographie pelvienne, hystérosalpingographie \\
\hline
\end{tabularx}
\end{center}

\courssub{Sensibilisation : prévenir l'infertilité}

L'infertilité n'est pas une fatalité : une large part des causes est \textbf{évitable}. La prévention repose sur le dépistage et le traitement précoce des IST (dépistage gratuit du VIH et prise en charge des IST dans de nombreux centres de santé au Cameroun, ex. Centre Pasteur, CePREF), le recours à des soins encadrés en cas d'interruption de grossesse (légalement encadrée), une hygiène de vie saine (limitation de l'alcool et du tabac, alimentation équilibrée), et surtout la \textbf{consultation précoce du couple} dès 12 mois d'essais infructueux (6 mois si la femme a plus de 35 ans) — et non de la femme seule.

\begin{methode}[Élaborer un outil de sensibilisation sur l'infertilité]
\begin{enumerate}
  \item Choisir un \textbf{public cible} précis (élèves de lycée, jeunes couples, chefs de famille, leaders d'opinion).
  \item Formuler un \textbf{message clair et unique} : « L'infertilité est un problème de couple : consultons ensemble dès 12 mois ».
  \item Illustrer les \textbf{causes} par un schéma simple (répartition en tiers, IST et leurs conséquences sur les trompes et les canaux déférents).
  \item Donner les \textbf{moyens de prévention} concrets : dépistage et traitement systématique des IST chez les deux partenaires, hygiène de vie, éviter l'automédication et les avortements à risque.
  \item Terminer par un \textbf{appel à l'action} : « Couple infertile ? Rendez-vous ensemble à la formation sanitaire la plus proche pour un bilan. »
\end{enumerate}
\end{methode}

\begin{exoresolu}[Interpréter un bilan d'infertilité]
Énoncé. Un couple consulte après 2 ans de rapports réguliers non protégés sans grossesse. Chez la femme : cycles réguliers de 28 jours, courbe de température biphasique, dosages hormonaux normaux (progestérone J21 \SI{25}{ng.mL^{-1}}) ; l'hystérosalpingographie montre deux trompes obstruées (pas d'épandage péritonéal). Chez l'homme : spermogramme avec \SI{22e6}{mL^{-1}} spermatozoïdes et \SI{45}{\percent} de mobilité progressive. (1) Ce couple est-il infertile ? Justifier. (2) Identifier le ou les facteurs en cause. (3) Proposer une explication probable de l'anomalie observée et une attitude préventive.
\tcblower
Solution. (1) Oui : l'absence de grossesse après 24 mois de rapports réguliers sans contraception correspond à la définition OMS de l'infertilité du couple. (2) La femme ovule normalement (courbe biphasique, cycles réguliers, progestérone J21 $> \SI{10}{ng.mL^{-1}}$, hormones basales normales) mais ses trompes sont obstruées bilatéralement : cause féminine mécanique (facteur tubaire). L'homme présente un spermogramme normal (concentration $> \SI{15e6}{mL^{-1}}$, mobilité progressive $> \SI{32}{\percent}$) : il n'est pas en cause. (3) L'obstruction tubaire bilatérale évoque une salpingite ancienne, le plus souvent consécutive à une IST (\emph{Chlamydia trachomatis}, \emph{Neisseria gonorrhoeae}) non traitée ou mal traitée. Prévention : dépistage et traitement précoce des IST, y compris chez le partenaire (traitement unique dose d'azithromycine + céfriaxone pour la gonococcie), et consultation précoce en couple. Une FIVETE (section V) pourra être proposée à ce couple en raison de l'obstruction tubaire bilatérale irréversible.
\end{exoresolu}

\begin{attention}[Pièges de rédaction au Baccalauréat]
\begin{itemize}
  \item Ne confonds pas \textbf{infertilité} (difficulté à concevoir, potentiellement traitable) et \textbf{stérilité} (incapacité définitive, absolue).
  \item N'écris jamais que « la femme est responsable de l'infertilité » : cite la répartition en trois tiers (masculin, féminin, mixte/inexpliqué).
  \item Le spermogramme concerne l'\textbf{homme} ; l'hystérosalpingographie concerne la \textbf{femme} (perméabilité des trompes) : n'inverse pas les examens.
  \item Une courbe de température \textbf{biphasique} indique une ovulation ; une courbe \textbf{monophasique} (plate) évoque une anovulation.
  \item L'oligospermie est une \textbf{diminution} du nombre, l'azoospermie est une \textbf{absence} totale : ne confonds pas les termes.
\end{itemize}
\end{attention}

\begin{exercice}[Pour t'entraîner]
Dans un service de gynécologie de Douala, on recense 90 couples infertiles. En appliquant la répartition admise, combien de cas environ relèvent d'une cause masculine ? d'une cause féminine ? d'une cause mixte ou inexpliquée ? Rédige ensuite le slogan et les trois idées forces d'une affiche de sensibilisation destinée aux élèves de Terminale.
\end{exercice}

\begin{corrige}[Exercice]
Environ $90/3 = 30$ cas pour chaque catégorie : 30 causes masculines, 30 causes féminines, 30 causes mixtes ou inexpliquées. Affiche attendue : slogan clair (ex. « Protège ta fertilité : dépiste et traite les IST ») ; idées forces : (1) l'infertilité est un problème de couple ($1/3$, $1/3$, $1/3$) ; (2) les IST non traitées bouchent les trompes et les canaux déférents ; (3) consulter en couple dès 12 mois d'essais, dans une formation sanitaire.
\end{corrige}

\coursec{V. La procréation médicalement assistée (PMA)}

Nous venons de voir, dans la section précédente, que l'infertilité peut avoir des causes variées : azoospermie ou oligospermie chez l'homme, obstruction des trompes, troubles de l'ovulation ou anomalies de la glaire cervicale chez la femme. Une question se pose alors naturellement : \emph{que peut faire la médecine lorsque la fécondation naturelle est impossible ou difficile ?} C'est précisément l'objet de la \cle{procréation médicalement assistée} (PMA). Dans un pays comme le Cameroun, où la parentalité occupe une place sociale et culturelle centrale, ces techniques répondent à une souffrance réelle de nombreux couples.

\begin{definition}[Procréation médicalement assistée]
La \cle{procréation médicalement assistée} (PMA) est l'ensemble des techniques médicales et biologiques visant à obtenir une grossesse lorsque la fécondation naturelle est impossible ou difficile. Les deux techniques principales au programme sont l'\cle{insémination artificielle} (IA) et la \cle{FIVETE} (fécondation \emph{in vitro} et transfert d'embryon).
\end{definition}

L'idée directrice est simple à retenir : selon la cause de l'infertilité, le médecin va soit \textbf{aider les gamètes à se rencontrer dans le corps de la femme} (insémination artificielle : fécondation \emph{in vivo}), soit \textbf{faire se rencontrer les gamètes hors du corps, au laboratoire} (FIVETE : fécondation \emph{in vitro}).

\courssub{L'insémination artificielle (IA)}

\textbf{Observation et problème.} Un couple consulte au centre hospitalier : les examens montrent une oligospermie modérée chez l'homme (concentration en spermatozoïdes faible, mobilité réduite), mais des trompes perméables et une ovulation normale chez la femme. Le sperme contient bien des spermatozoïdes, mais ils sont trop peu nombreux ou trop peu mobiles pour franchir seuls le col de l'utérus et rejoindre l'ovocyte. \emph{Hypothèse} : si l'on dépose directement les spermatozoïdes, préalablement sélectionnés et concentrés, à proximité de l'ovocyte au bon moment, la fécondation devient possible. C'est le principe de l'insémination artificielle.

\begin{definition}[Insémination artificielle]
L'\cle{insémination artificielle} (IA) consiste à introduire, à l'aide d'une fine sonde, des \textbf{spermatozoïdes préparés} (lavés, sélectionnés et concentrés en laboratoire) dans les voies génitales féminines, généralement \textbf{dans l'utérus} (insémination intra-utérine), \textbf{au moment de l'ovulation}. La fécondation a lieu \emph{in vivo}, c'est-à-dire dans le corps de la femme, dans la trompe, comme lors d'une fécondation naturelle.
\end{definition}

\begin{methode}[Déroulement d'une insémination artificielle]
\begin{enumerate}
\item \textbf{Repérage de l'ovulation} : suivi échographique de la croissance du follicule et dosage de l'hormone LH (le pic de LH annonce l'ovulation, environ 24 à 36 h plus tard) ; parfois une légère stimulation ovarienne est associée.
\item \textbf{Recueil et préparation du sperme} : le sperme est recueilli le jour de l'insémination, puis traité au laboratoire pour éliminer le plasma séminal et ne garder que les spermatozoïdes les plus mobiles, concentrés dans un petit volume.
\item \textbf{Insémination} : les spermatozoïdes préparés sont déposés dans la cavité utérine à l'aide d'un cathéter fin, geste indolore de quelques minutes.
\item \textbf{Fécondation naturelle} : les spermatozoïdes remontent dans la trompe où l'un d'eux peut féconder l'ovocyte ; la nidation a lieu dans l'utérus 6 à 7 jours plus tard.
\end{enumerate}
\end{methode}

\textbf{Indications et condition indispensable.} L'IA est indiquée en cas d'\textbf{infertilité masculine modérée} (oligospermie ou asthénospermie légères), de \textbf{troubles de la glaire cervicale} (glaire hostile qui bloque la progression des spermatozoïdes) ou de certaines infertilités inexpliquées. La condition absolue est que \textbf{les trompes soient perméables} : en effet, la fécondation ayant lieu dans la trompe, une obstruction tubaire rend l'IA inutile. On distingue deux variantes :
\begin{itemize}
\item l'\cle{IAC} : insémination artificielle avec le sperme du \textbf{conjoint} ;
\item l'\cle{IAD} : insémination artificielle avec le sperme d'un \textbf{donneur}, lorsque l'homme ne produit aucun spermatozoïde utilisable (azoospermie totale) ou porte une maladie génétique transmissible grave.
\end{itemize}

\courssub{La FIVETE : fécondation in vitro et transfert d'embryon}

\textbf{Observation et problème.} Une femme présente une obstruction bilatérale des trompes, séquelle d'infections génitales mal soignées. Ses ovaires produisent des ovocytes normaux et son utérus est sain, mais jamais les spermatozoïdes ne pourront rencontrer l'ovocyte : le « pont » est coupé. \emph{Hypothèse} : si l'on prélève les ovocytes et qu'on les met en contact avec les spermatozoïdes \textbf{hors du corps}, dans une boîte de laboratoire, on peut obtenir un embryon qu'il suffira ensuite de replacer dans l'utérus. C'est exactement le raisonnement qui a conduit les chercheurs britanniques Edwards et Steptoe à la naissance historique de Louise Brown.

\begin{exemplebox}[Une naissance historique]
Le premier « bébé FIVETE » du monde, \textbf{Louise Brown}, est né le 25 juillet 1978 en Angleterre. Sa mère souffrait d'une obstruction des trompes. Depuis, plusieurs millions d'enfants sont nés grâce à cette technique dans le monde, et Robert Edwards a reçu le prix Nobel de physiologie ou médecine en 2010. Le taux de succès actuel est d'environ \SI{25}{\percent} par cycle, toutes causes et tous âges confondus.
\end{exemplebox}

\begin{definition}[FIVETE]
La \cle{FIVETE} (fécondation \emph{in vitro} et transfert d'embryon) est une technique de procréation médicalement assistée consistant à réaliser la \textbf{fécondation en laboratoire} (\emph{in vitro}, littéralement « dans le verre », c'est-à-dire en boîte de culture), à laisser se développer l'embryon quelques jours, puis à \textbf{transférer l'embryon dans l'utérus} de la femme où se poursuivra la grossesse.
\end{definition}

\begin{methode}[Les cinq étapes de la FIVETE]
\begin{enumerate}
\item \textbf{Stimulation ovarienne} : la femme reçoit des gonadotrophines (principalement FSH recombinante) pendant une dizaine de jours pour obtenir la maturation simultanée de plusieurs follicules, au lieu d'un seul par cycle ; le suivi se fait par échographie et dosages hormonaux (estradiol).
\item \textbf{Déclenchement de l'ovulation et ponction des ovocytes} : injection d'hCG (qui mime le pic de LH) pour déclencher la maturation finale ; ponction des ovocytes par voie transvaginale sous guidage échographique et anesthésie locale ou générale, environ 36 h après l'injection d'hCG.
\item \textbf{Fécondation in vitro} : les ovocytes sont mis en contact avec les spermatozoïdes préparés dans une boîte de culture (FIVETE classique) ; en cas d'infertilité masculine sévère, un spermatozoïde unique est injecté directement dans l'ovocyte : technique \cle{ICSI} (injection intracytoplasmique de spermatozoïde).
\item \textbf{Culture embryonnaire} : les embryons obtenus se développent en incubateur (37 °C, 5 \% CO\textsubscript{2}) pendant \textbf{2 à 5 jours} (stades 2 à 8 cellules, voire blastocyste) ; les embryons de meilleure qualité sont sélectionnés selon des critères morphologiques et cinétiques.
\item \textbf{Transfert d'embryon(s)} : un ou deux embryons (transfert unique de plus en plus recommandé) sont déposés dans la cavité utérine à l'aide d'un cathéter fin ; l'embryon s'implante alors dans l'endomètre (nidation) et la grossesse se déroule normalement. Les embryons surnuméraires de bonne qualité peuvent être vitrifiés (congelés) pour une tentative ultérieure.
\end{enumerate}
\end{methode}

\begin{popfigure}
\begin{center}
\begin{tikzpicture}[font=\small, node distance=0.55cm and 0.9cm,
  etape/.style={draw=popblue, thick, fill=popblueL, rounded corners=3pt, text width=2.5cm, minimum height=1.5cm, align=center},
  num/.style={circle, draw=popdark, fill=poporange, text=white, font=\bfseries\small, inner sep=1.5pt},
  fleche/.style={-{Stealth[length=2.5mm]}, thick, popdark}]
\node[etape, align=center] (e1){\textbf{1. Stimulation ovarienne}\\ (FSH)};
\node[etape, right=of e1, align=center] (e2){\textbf{2. Ponction des ovocytes}\\ (36 h post-hCG)};
\node[etape, right=of e2, align=center] (e3){\textbf{3. Fécondation \emph{in vitro}}\\ (FIVETE ou ICSI)};
\node[etape, below=of e3, align=center] (e4){\textbf{4. Culture embryonnaire}\\ (2 à 5 jours)};
\node[etape, left=of e4, align=center] (e5){\textbf{5. Transfert dans l'utérus}\\ (nidation)};
\draw[fleche] (e1) -- (e2);
\draw[fleche] (e2) -- (e3);
\draw[fleche] (e3) -- (e4);
\draw[fleche] (e4) -- (e5);
\draw[fleche, dashed, popgreen] (e5.west) -- ++(-1.1,0) |- (e1.west) node[midway, left, align=center, text=popgreen]{Échec :\\ nouvelle tentative\\ (embryons congelés)};
\end{tikzpicture}
\end{center}
\legende{Les cinq étapes de la FIVETE : (1) stimulation ovarienne par FSH, (2) déclenchement par hCG et ponction des ovocytes à J+36 h, (3) fécondation \emph{in vitro} (classique ou ICSI), (4) culture embryonnaire de 2 à 5 jours en incubateur, (5) transfert de l'embryon dans l'utérus. En cas d'échec de nidation, un nouveau cycle peut être tenté avec des embryons vitrifiés.}
\end{popfigure}

\textbf{Indications de la FIVETE.} La FIVETE est proposée lorsque l'IA est impossible ou a échoué :
\begin{itemize}
\item \textbf{obstruction ou absence des trompes} (cause tubaire) : c'est l'indication historique ;
\item \textbf{infertilité masculine sévère} : oligo-asthéno-tératospermie grave, azoospermie avec spermatozoïdes prélevés dans le testicule (TESE) ou l'épididyme (MESA) — technique ICSI associée ;
\item \textbf{infertilité inexpliquée} : lorsque tous les examens sont normaux mais qu'aucune grossesse ne survient ;
\item endométriose sévère ou échecs répétés d'inséminations artificielles.
\end{itemize}

\begin{attention}[Erreur fréquente : le « bébé éprouvette »]
Contrairement à une idée répandue, la FIVETE ne produit \textbf{pas} un enfant qui se développerait entièrement dans une éprouvette ou hors du corps de la mère. Seules la \textbf{fécondation} et les \textbf{2 à 5 premiers jours} du développement embryonnaire ont lieu \emph{in vitro}. Dès le transfert, l'embryon s'implante dans l'\textbf{utérus de la femme} et la grossesse (environ 9 mois) se déroule exactement comme une grossesse naturelle. Au baccalauréat, écris toujours : « fécondation \emph{in vitro}, développement \emph{in utero} ».
\end{attention}

\courssub{Comparaison entre insémination artificielle et FIVETE}

Le critère fondamental qui distingue les deux techniques est le \textbf{lieu de la fécondation} : \emph{in vivo} (dans la trompe) pour l'IA, \emph{in vitro} (en laboratoire) pour la FIVETE. Le schéma suivant résume cette différence essentielle.

\begin{popfigure}
\begin{center}
\begin{tikzpicture}[font=\small, scale=0.95, every node/.style={transform shape}]
% --- IA : fécondation in vivo ---
\node[draw=popgreen, thick, fill=popgreenL, rounded corners=3pt, minimum width=6.2cm, minimum height=3.4cm] (ia) at (0,0) {};
\node[font=\bfseries, text=popdark] at (0,1.45) {IA : fécondation \emph{in vivo}};
\draw[fill=poppinkL, draw=poppink, thick] (-1.6,0.55) ellipse (0.35 and 0.28);
\node[font=\scriptsize] at (-1.6,0.55) {ovaire};
\draw[poppink, thick] (-1.25,0.55) .. controls (-0.5,0.9) and (0.2,0.9) .. (0.9,0.4);
\node[font=\scriptsize, text=poppink] at (-0.2,1.0) {trompe (perméable)};
\draw[fill=poporangeL, draw=poporange, thick] (0.9,0.1) ellipse (0.55 and 0.75);
\node[font=\scriptsize] at (0.9,-0.05) {utérus};
\draw[fill=popblue, draw=popblue] (-0.35,0.72) circle (0.09);
\node[font=\scriptsize, text=popblue] at (-0.35,0.35) {ovocyte};
\foreach \x/\y in {1.6/0.5,1.9/0.3,2.2/0.6,1.7/0.05,2.35/0.2}
  \draw[popblue, thick] (\x,\y) -- ++(-0.22,0.06);
\node[font=\scriptsize, text=popblue, align=center] at (2.1,-0.35) {spermatozoïdes\\ déposés dans l'utérus};
\draw[-{Stealth}, thick, popblue] (1.55,0.45) .. controls (1.0,0.75) and (0.3,0.85) .. (-0.2,0.78);
\node[font=\scriptsize, align=center, text=popdark] at (0,-1.35) {rencontre des gamètes\\ \textbf{dans la trompe}};
% --- FIVETE : fécondation in vitro ---
\node[draw=poppurple, thick, fill=poppurpleL, rounded corners=3pt, minimum width=6.2cm, minimum height=3.4cm] (fiv) at (8.6,0) {};
\node[font=\bfseries, text=popdark] at (8.6,1.45) {FIVETE : fécondation \emph{in vitro}};
\draw[fill=poppinkL, draw=poppink, thick] (6.7,0.55) ellipse (0.35 and 0.28);
\node[font=\scriptsize] at (6.7,0.55) {ovaire};
\draw[-{Stealth}, thick, popdark] (7.1,0.55) -- (8.0,0.55) node[midway, above, font=\scriptsize]{ponction};
\draw[draw=popblue, thick, fill=white] (8.6,0.45) ellipse (0.75 and 0.5);
\node[font=\scriptsize, text=popblue] at (8.6,1.02) {boîte de culture};
\draw[fill=popblue, draw=popblue] (8.45,0.45) circle (0.09);
\foreach \x/\y in {8.85/0.6,8.95/0.35,8.8/0.25}
  \draw[popblue, thick] (\x,\y) -- ++(-0.2,0.05);
\draw[-{Stealth}, thick, popdark] (9.4,0.45) -- (10.3,0.3) node[midway, above, font=\scriptsize]{transfert};
\draw[fill=poporangeL, draw=poporange, thick] (10.85,0.05) ellipse (0.5 and 0.7);
\node[font=\scriptsize] at (10.85,-0.1) {utérus};
\draw[fill=popgold, draw=popgold] (10.7,0.25) circle (0.1);
\node[font=\scriptsize, text=popdark] at (11.6,0.45) {embryon};
\node[font=\scriptsize, align=center, text=popdark] at (8.6,-1.35) {rencontre des gamètes\\ \textbf{au laboratoire}};
\end{tikzpicture}
\end{center}
\legende{Comparaison des deux techniques de PMA. À gauche, insémination artificielle : les spermatozoïdes préparés sont déposés dans l'utérus et la fécondation a lieu \emph{in vivo}, dans la trompe (qui doit être perméable). À droite, FIVETE : les ovocytes sont ponctionnés, fécondés \emph{in vitro} en boîte de culture, puis l'embryon est transféré dans l'utérus.}
\end{popfigure}

\begin{center}
\begin{tabularx}{\linewidth}{|l|X|X|}
\hline
\rowcolor{popblueL} \textbf{Critère} & \textbf{Insémination artificielle} & \textbf{FIVETE} \\
\hline
Lieu de la fécondation & \emph{In vivo} (trompe de la femme) & \emph{In vitro} (boîte de culture) \\
\hline
Trompes perméables & Indispensables & Non nécessaires \\
\hline
Indications principales & Infertilité masculine modérée, glaire cervicale hostile & Obstruction tubaire, infertilité masculine sévère, infertilité inexpliquée \\
\hline
Stimulation ovarienne & Absente ou légère & Systématique et forte (FSH) \\
\hline
Taux de succès par tentative & Environ 10 à \SI{15}{\percent} & Environ 20 à \SI{30}{\percent} \\
\hline
Coût et lourdeur & Modérés, geste simple & Élevés, technique lourde \\
\hline
\end{tabularx}
\end{center}

\courssub{Limites, données chiffrées et questions éthiques}

\textbf{Un succès qui dépend fortement de l'âge.} Les données des centres de PMA montrent que le taux de réussite par tentative diminue nettement avec l'âge de la femme, car la qualité des ovocytes décline, en particulier après 35 ans. L'histogramme ci-dessous illustre des ordres de grandeur classiquement observés.

\begin{popfigure}
\begin{center}
\begin{tikzpicture}
\begin{axis}[
  ybar, bar width=0.55cm,
  width=0.85\linewidth, height=6.2cm,
  ymin=0, ymax=40,
  ylabel={Taux de naissance par tentative (\%)},
  xlabel={Âge de la femme},
  symbolic x coords={{moins de 35 ans},{35--37 ans},{38--40 ans},{plus de 40 ans}},
  xtick=data, x tick label style={font=\small},
  ytick={0,10,20,30,40},
  nodes near coords, nodes near coords style={font=\small\bfseries},
  axis lines=left, enlarge x limits=0.2,
  every axis plot/.append style={fill=popblue, draw=popdark}
]
\addplot coordinates {({moins de 35 ans},30) ({35--37 ans},25) ({38--40 ans},15) ({plus de 40 ans},7)};
\end{axis}
\end{tikzpicture}
\end{center}
\legende{Taux de naissance par tentative de FIVETE selon l'âge de la femme (ordres de grandeur). Le succès passe d'environ \SI{30}{\percent} avant 35 ans à moins de \SI{10}{\percent} après 40 ans, en raison du vieillissement des ovocytes.}
\end{popfigure}

\begin{exoresolu}[Exploiter des données sur la FIVETE]
\textbf{Énoncé.} Dans un centre de PMA, on observe les taux de naissance par tentative de FIVETE suivants : \SI{30}{\percent} chez les femmes de moins de 35 ans et \SI{7}{\percent} chez les femmes de plus de 40 ans. (1) Calcule le nombre moyen de naissances attendues pour 100 tentatives dans chaque groupe. (2) Propose une interprétation biologique de cette différence. (3) Explique pourquoi la stimulation ovarienne et le transfert de plusieurs embryons ont été historiquement utilisés pour compenser ce faible rendement, et quel risque cela comporte.
\tcblower
\textbf{Solution.} (1) Pour 100 tentatives : $100 \times \num{0,30} = 30$ naissances attendues avant 35 ans ; $100 \times \num{0,07} = 7$ naissances après 40 ans. (2) Avec l'âge, la réserve ovarienne diminue et surtout la qualité des ovocytes se dégrade : les anomalies chromosomiques (comme les trisomies dues à des non-disjonctions en méiose I ou II) deviennent plus fréquentes, ce qui réduit la fécondabilité et provoque davantage d'échecs de nidation ou de fausses couches précoces. (3) La stimulation ovarienne permet d'obtenir plusieurs ovocytes en un seul cycle, donc plusieurs embryons, ce qui augmente les chances qu'au moins un embryon de bonne qualité soit disponible. Transférer plusieurs embryons augmente le taux de grossesse, mais au prix d'un risque élevé de \textbf{grossesse multiple} (jumeaux, triplés), source de prématurité et de complications pour la mère et les enfants ; c'est pourquoi les centres modernes limitent le nombre d'embryons transférés (souvent un seul).
\end{exoresolu}

\textbf{Limites pratiques.} Même dans les meilleures conditions, le taux de réussite reste limité à environ \textbf{20 à \SI{30}{\percent} par tentative} : plusieurs cycles sont souvent nécessaires, ce qui représente une épreuve physique (stimulation hormonale, ponctions) et psychologique pour le couple. Le coût est élevé et l'accès reste très inégal.

\begin{savaistu}[La PMA au Cameroun]
Au Cameroun, des centres de procréation médicalement assistée existent à \textbf{Yaoundé} et à \textbf{Douala}, et des centaines d'enfants y sont déjà nés grâce à la FIVETE. Cependant, le coût d'une seule tentative de FIVETE dépasse souvent \num{1500000} FCFA, sans prise en charge par une assurance maladie, ce qui la rend inaccessible à la majorité des couples. La prévention des infections sexuellement transmissibles — première cause d'obstruction des trompes — demeure donc le moyen le plus efficace et le moins coûteux de lutter contre l'infertilité.
\end{savaistu}

\textbf{Questions éthiques (débat citoyen).} La PMA soulève des questions de société sur lesquelles tu peux être amené à réfléchir : le \textbf{don de gamètes} (sperme ou ovocytes d'un donneur) pose la question de l'anonymat et de l'accès de l'enfant à ses origines ; la \textbf{vitrification d'embryons} surnuméraires interroge sur leur devenir (don, recherche, destruction) ; le coût élevé pose le problème de l'\textbf{équité d'accès} entre les couples selon leurs revenus. Ces questions n'ont pas de réponse scientifique unique : elles relèvent d'un débat citoyen éclairé par les connaissances biologiques que tu viens d'acquérir. (Complément utile : renseigne-toi sur la réglementation en vigueur dans ton pays.)

\begin{exercice}[Élaborer un outil de sensibilisation]
À partir des connaissances acquises sur les causes de l'infertilité (section IV) et les techniques de PMA (section V), réalise une affiche ou une fiche d'information grand public (format A5) destinée à une salle d'attente de centre de santé au Cameroun. Ton support doit : (1) rappeler les principales causes évitables d'infertilité (IST, avortements à risque, exposition à des toxiques) ; (2) expliquer simplement la différence entre IA et FIVETE ; (3) indiquer les structures de référence au Cameroun ; (4) insister sur la prévention comme priorité. Tu pourras utiliser des schémas simplifiés et un langage accessible.
\end{exercice}

\begin{aretenir}[L'essentiel de la section]
\begin{itemize}
\item La PMA regroupe les techniques visant à obtenir une grossesse lorsque la fécondation naturelle est impossible ou difficile.
\item \textbf{Insémination artificielle} : dépôt de spermatozoïdes préparés dans l'utérus au moment de l'ovulation ; fécondation \emph{in vivo} ; exige des trompes perméables ; variantes IAC (sperme du conjoint) et IAD (sperme de donneur).
\item \textbf{FIVETE} : stimulation ovarienne (FSH) $\rightarrow$ déclenchement (hCG) + ponction $\rightarrow$ fécondation \emph{in vitro} (classique ou ICSI) $\rightarrow$ culture embryonnaire (2 à 5 jours) $\rightarrow$ transfert dans l'utérus ; indiquée en cas d'obstruction tubaire, d'infertilité masculine sévère ou inexpliquée.
\item Le taux de succès est d'environ 20 à \SI{30}{\percent} par tentative et diminue fortement avec l'âge de la femme (qualité ovocytaire) ; le coût élevé limite l'accès au Cameroun.
\item La prévention des IST reste la stratégie la plus efficace contre l'infertilité tubaire.
\end{itemize}
\end{aretenir}

\coursec{VI. Synthèse et intégration}

Après avoir découvert les techniques de procréation médicalement assistée (insémination artificielle, FIVETE), il est temps de prendre du recul et de relier tous les savoirs de cette leçon. L'objectif de cette synthèse est double : d'une part, organiser les connaissances autour d'un principe unificateur — l'\cle{axe hypothalamo-hypophyso-gonadique} — ; d'autre part, te donner les outils méthodologiques et citoyens pour traiter un sujet de Baccalauréat et pour agir dans ta communauté.

\courssub{Bilan comparatif : régulation continue chez l'homme, cyclique chez la femme}

\textbf{Intuition.} Chez l'homme, la production de spermatozoïdes commence à la puberté et se poursuit de façon quasi ininterrompue toute la vie. Chez la femme, en revanche, tout est rythmé : un ovocyte est libéré environ tous les 28 jours, entre la puberté et la ménopause. Cette différence fondamentale s'explique par le \emph{type de régulation hormonale}.

Chez l'homme, la régulation est \cle{continue} (tonique) : l'hypothalamus sécrète la GnRH de façon régulière, l'hypophyse libère FSH et LH à débit constant, et les testicules répondent en permanence. Chez la femme, la régulation est \cle{cyclique} : les taux de FSH, LH, œstrogènes et progestérone varient au cours du cycle ovarien, avec un pic de LH déclenchant l'ovulation vers le 14\textsuperscript{e} jour d'un cycle de 28 jours.

\begin{propriete}[Principe unificateur de la leçon (admis)]
Toute hormone gonadique en excès (testostérone chez l'homme ; œstrogènes et progestérone chez la femme) exerce un \cle{rétrocontrôle négatif} sur l'hypothalamus et l'hypophyse, freinant la libération de GnRH, de FSH et de LH. Ce principe explique à la fois la stabilité de la fonction reproductive masculine, les phases du cycle féminin, le fonctionnement de la pilule contraceptive et certaines causes d'infertilité.
\end{propriete}

\begin{attention}[Une exception à connaître : le rétrocontrôle positif]
Chez la femme, en fin de phase folliculaire, le taux élevé d'\cle{œstrogènes} sécrétés par le follicule mûr exerce exceptionnellement un \cle{rétrocontrôle positif} sur l'axe hypothalamo-hypophysaire : il provoque le \emph{pic de LH} (et de FSH) qui déclenche l'ovulation. C'est le seul cas de rétrocontrôle positif du programme ; tous les autres sont négatifs.
\end{attention}

\begin{center}
\begin{tabularx}{\linewidth}{|l|X|X|}
\hline
\rowcolor{popblueL} \textbf{Critère} & \textbf{Homme} & \textbf{Femme} \\
\hline
Hormones hypothalamique et hypophysaires & GnRH ; FSH (spermatogenèse), LH (sécrétion de testostérone) & GnRH ; FSH (croissance folliculaire), LH (ovulation, corps jaune) \\
\hline
Hormones gonadiques & Testostérone (cellules de Leydig, interstitielles) & Œstrogènes (follicules), progestérone (corps jaune) \\
\hline
Type de régulation & Continue (tonique), de la puberté à la vieillesse & Cyclique (environ 28 jours), de la puberté à la ménopause \\
\hline
Gamétogenèse & Spermatogenèse continue (des millions de spermatozoïdes par jour) & Ovogenèse cyclique (1 ovocyte libéré par cycle en général) \\
\hline
Causes d'infertilité fréquentes & Oligospermie, azoospermie, varicocèle, IST non traitées, troubles hormonaux & Anovulation, obstruction des trompes (salpingite), endométriose, anomalies utérines, troubles hormonaux \\
\hline
Examens de première intention & Spermogramme, dosages de la testostérone, de la FSH et de la LH & Dosages hormonaux (FSH, LH, œstradiol, progestérone), échographie pelvienne, hystérosalpingographie \\
\hline
\end{tabularx}
\end{center}

\begin{attention}[Contraception et PMA : deux démarches opposées]
Erreur fréquente à l'examen : confondre \cle{contraception} et \cle{PMA}. La contraception (pilule, préservatif, stérilet) \emph{empêche} la fécondation ou la nidation — la pilule œstroprogestative exploite d'ailleurs le rétrocontrôle négatif pour bloquer le pic de LH et donc l'ovulation. La PMA (insémination artificielle, FIVETE) \emph{favorise} la rencontre des gamètes chez un couple infertile. Les deux agissent sur la reproduction, mais dans des sens opposés. Une phrase comme « la FIVETE est un moyen de contraception » est une faute grave.
\end{attention}

\courssub{Le lien logique : une perturbation de l'axe peut causer une infertilité}

\textbf{Intuition.} L'axe hypothalamo-hypophyso-gonadique fonctionne comme une chaîne de commande à trois maillons : l'hypothalamus (GnRH) commande l'hypophyse (FSH, LH), qui commande les gonades (testicules, ovaires), lesquelles informent en retour les étages supérieurs (rétrocontrôle). Si un seul maillon est défaillant, toute la chaîne est perturbée et la fertilité peut être compromise.

\begin{itemize}
\item \textbf{Niveau hypothalamique} : stress intense, anorexie, tumeur $\Rightarrow$ déficit en GnRH $\Rightarrow$ FSH et LH insuffisantes $\Rightarrow$ arrêt de la gamétogenèse ou anovulation.
\item \textbf{Niveau hypophysaire} : adénome hypophysaire $\Rightarrow$ sécrétion anormale de FSH/LH $\Rightarrow$ dysfonctionnement gonadique.
\item \textbf{Niveau gonadique} : varicocèle, orchite (testicules) ; ovaires polykystiques, insuffisance ovarienne prématurée (ovaires) $\Rightarrow$ production de gamètes ou d'hormones altérée.
\item \textbf{Niveau des voies génitales} : obstruction des trompes après une IST non traitée (Chlamydiae, gonocoques), obstruction des canaux déférents $\Rightarrow$ les gamètes ne se rencontrent plus.
\end{itemize}

\begin{popfigure}
\begin{tikzpicture}[mindmap, grow cyclic, every node/.style={concept, font=\small}, root concept/.append style={concept color=popblue, font=\small\bfseries}, level 1/.append style={sibling angle=120, level distance=3.6cm}, level 2/.append style={level distance=2.6cm, font=\scriptsize}]
\node[root concept]{Santé reproductive}
  child[concept color=popgreen] { node {Régulation hormonale}
    child { node {Homme : continue (FSH, LH, testostérone)} }
    child { node {Femme : cyclique (FSH, LH, œstrogènes, progestérone)} }
    child { node {Rétrocontrôle négatif} }
  }
  child[concept color=poporange] { node {Infertilité}
    child { node {Causes masculines (spermogramme anormal)} }
    child { node {Causes féminines (anovulation, trompes)} }
    child { node {Diagnostic : interrogatoire, dosages, imagerie} }
  }
  child[concept color=poppurple] { node {PMA}
    child { node {Insémination artificielle} }
    child { node {FIVETE} }
  };
\end{tikzpicture}
\legende{Carte mentale récapitulative de la leçon : la régulation hormonale (axe hypothalamo-hypophyso-gonadique) fonde la fertilité ; sa perturbation explique l'infertilité ; la PMA propose des solutions adaptées à chaque cause.}
\end{popfigure}

\courssub{La démarche diagnostique : de l'interrogatoire aux examens}

\textbf{Observation.} Au Cameroun comme ailleurs, un couple est dit \cle{infertile} lorsqu'il n'obtient pas de grossesse après environ 12 mois de rapports sexuels réguliers sans contraception. L'enquête médicale porte toujours sur \emph{les deux partenaires} : dans environ un tiers des cas la cause est masculine, un tiers féminine, un tiers mixte ou inexpliquée.

\begin{methode}[Démarche diagnostique d'une infertilité de couple]
\begin{enumerate}
\item \textbf{Interrogatoire du couple} : âge, antécédents (IST, oreillons, chirurgie pelvienne), mode de vie (tabac, alcool), fréquence des rapports.
\item \textbf{Examen clinique} des deux partenaires (organes génitaux, recherche de varicocèle, signes de trouble hormonal).
\item \textbf{Examens paracliniques chez l'homme} : \cle{spermogramme} (volume, concentration, mobilité, morphologie des spermatozoïdes), dosages de FSH, LH et testostérone.
\item \textbf{Examens paracliniques chez la femme} : dosages hormonaux (FSH, LH, œstradiol, progestérone en deuxième partie de cycle), échographie pelvienne (ovaires, utérus), hystérosalpingographie (perméabilité des trompes), courbe de température.
\item \textbf{Synthèse} : identification de la ou des causes, puis orientation vers le traitement adapté (médical, chirurgical ou PMA).
\end{enumerate}
\end{methode}

\courssub{De la cause au traitement : choisir la technique de PMA}

\begin{center}
\begin{tabularx}{\linewidth}{|X|X|}
\hline
\rowcolor{popgreenL} \textbf{Cause d'infertilité identifiée} & \textbf{Technique adaptée} \\
\hline
Oligospermie modérée, glaire cervicale hostile & Insémination artificielle avec sperme du conjoint (IAC) \\
\hline
Azoospermie sécrétoire, absence de spermatozoïdes & Insémination avec sperme de donneur (IAD) \\
\hline
Obstruction ou absence des trompes & FIVETE (fécondation in vitro puis transfert d'embryon) \\
\hline
Anovulation (ovaires polykystiques) & Induction de l'ovulation (traitement hormonal) avant toute PMA \\
\hline
Endométriose sévère, échecs d'insémination & FIVETE \\
\hline
Infertilité inexpliquée persistante & FIVETE en dernier recours \\
\hline
\end{tabularx}
\end{center}

\begin{exoresolu}[Analyse d'un spermogramme et orientation diagnostique]
Un couple consulte au Centre hospitalier d'Essos (Yaoundé) après 3 ans d'attente d'un enfant. Le spermogramme de monsieur N. donne : concentration en spermatozoïdes $= \SI{8}{millions/mL}$ (norme $\geq \SI{15}{millions/mL}$), mobilité $= \SI{25}{\%}$ (norme $\geq \SI{40}{\%}$). Les dosages hormonaux de madame N. et son hystérosalpingographie sont normaux. (1) Analyser le spermogramme. (2) Conclure sur l'origine probable de l'infertilité. (3) Proposer une technique de PMA adaptée.
\tcblower
\textbf{(1)} La concentration ($\SI{8}{millions/mL} < \SI{15}{millions/mL}$) est inférieure à la norme : c'est une \cle{oligospermie}. La mobilité ($\SI{25}{\%} < \SI{40}{\%}$) est aussi réduite : c'est une \cle{asthénospermie}. On parle d'\emph{oligoasthénospermie}.
\textbf{(2)} Les examens de madame étant normaux, l'infertilité du couple est d'origine \emph{masculine} : les spermatozoïdes, trop peu nombreux et peu mobiles, ont une faible probabilité d'atteindre et de féconder l'ovocyte.
\textbf{(3)} Une \cle{insémination artificielle} avec sperme du conjoint peut être tentée (les spermatozoïdes sont sélectionnés, concentrés et déposés directement dans l'utérus) ; en cas d'échec, une \cle{FIVETE} sera proposée. On recherchera aussi une cause traitable (varicocèle, IST) avant la PMA.
\end{exoresolu}

\courssub{Méthode pour résoudre un exercice type Baccalauréat}

Les sujets de Bac sur ce chapitre présentent presque toujours des \emph{documents expérimentaux} : courbes de dosages hormonaux, résultats de castration/greffe, spermogrammes. Voici la méthode en quatre étapes.

\begin{methode}[Analyser un document expérimental sur les hormones sexuelles]
\begin{enumerate}
\item \textbf{Identifier le document} : quelle grandeur est mesurée (hormone, nombre de gamètes), chez qui (homme, femme, animal), en fonction de quoi (temps, traitement) ? Lis attentivement les axes, les unités et la légende.
\item \textbf{Extraire les données chiffrées} : cite les valeurs remarquables avec leurs unités (ex. « le taux de LH passe de 5 à \SI{80}{UI/L} au 14\textsuperscript{e} jour »).
\item \textbf{Comparer témoin et expérience} : que devient le paramètre quand on supprime (castration, ablation) ou qu'on ajoute (injection d'hormone, greffe) un facteur ? La comparaison révèle le rôle du facteur.
\item \textbf{Conclure en reliant au savoir du cours} : formule le rôle de l'hormone ou de l'organe en mobilisant le vocabulaire exact (rétrocontrôle négatif, pic de LH, ovulation, spermatogenèse).
\end{enumerate}
\end{methode}

\begin{attention}[Pièges de rédaction au Bac]
\begin{itemize}
\item Ne jamais écrire « la courbe monte » sans citer de valeurs : toute affirmation doit s'appuyer sur des données chiffrées du document.
\item Ne pas confondre FSH et LH : la FSH agit sur la gamétogenèse (tubules séminifères, follicules), la LH sur la sécrétion des hormones stéroïdes (testostérone) et sur l'ovulation.
\item Ne pas attribuer à l'hypophyse la production de GnRH : la GnRH est \emph{hypothalamique}.
\end{itemize}
\end{attention}

\courssub{Rôle citoyen : sensibiliser, déstigmatiser, prévenir}

\begin{aretenir}[L'essentiel de la leçon]
La fertilité repose sur l'axe hypothalamo-hypophyso-gonadique : continu chez l'homme, cyclique chez la femme, et contrôlé dans les deux cas par un rétrocontrôle négatif des hormones gonadiques. Toute perturbation d'un maillon de cet axe peut entraîner une infertilité, dont le diagnostic repose sur l'interrogatoire du couple, le spermogramme, les dosages hormonaux et l'imagerie. La PMA (insémination artificielle, FIVETE) offre des solutions adaptées à chaque cause.
\end{aretenir}

\textbf{Agir en citoyen.} L'infertilité touche environ un couple sur dix ; au Cameroun, elle est encore trop souvent imputée à la seule femme, source de stigmatisation et de ruptures familiales. Or les causes masculines représentent environ un tiers des cas. Ton rôle de citoyen éclairé est de :

\begin{itemize}
\item \textbf{Sensibiliser} : élaborer des outils (affiches, sketches, causeries éducatives dans les clubs de santé scolaire) expliquant que l'infertilité est un problème \emph{de couple}, médical et non mystique.
\item \textbf{Combattre les idées reçues} : l'infertilité n'est ni une malédiction ni une faute morale ; elle a des causes identifiées et parfois traitables.
\item \textbf{Promouvoir le dépistage et le traitement précoce des IST} (infections à Chlamydiae, gonococcie, VIH) : non traitées, elles provoquent salpingites et obstructions des voies génitales, première cause évitable d'infertilité. Le préservatif et le dépistage régulier sont des gestes de protection de la fertilité future.
\item \textbf{Encourager la consultation précoce} : plus le bilan est réalisé tôt, plus les chances de traitement sont élevées.
\end{itemize}

\begin{savaistu}[La PMA au Cameroun]
Depuis les années 1990-2000, plusieurs centres hospitaliers camerounais (à Yaoundé et à Douala) pratiquent l'insémination artificielle et la FIVETE, permettant à des couples infertiles d'avoir des enfants sans se rendre à l'étranger. Ces techniques, encadrées par la loi et l'éthique médicale, illustrent concrètement l'application des savoirs de cette leçon à la santé publique nationale.
\end{savaistu}

\begin{exercice}[Pour t'entraîner]
Un couple consulte pour infertilité de 2 ans. Chez madame, les dosages montrent une progestérone effondrée en deuxième partie de cycle et l'échographie ne montre aucun follicule mature ; chez monsieur, le spermogramme est normal. (1) Interpréter les résultats de madame. (2) Identifier le maillon de l'axe probablement en cause. (3) Proposer une prise en charge.
\end{exercice}

\begin{corrige}[Exercice 1]
(1) L'absence de follicule mature et la progestérone très faible indiquent une \cle{anovulation} : sans ovulation, pas de corps jaune, donc pas de sécrétion de progestérone en deuxième partie de cycle.
(2) Le maillon en cause peut être hypothalamo-hypophysaire (déficit en GnRH ou en FSH/LH, par exemple lié au stress ou à un trouble nutritionnel) ou ovarien (ovaires polykystiques) ; des dosages de FSH et de LH permettraient de trancher : des taux bas orientent vers une cause centrale, des taux élevés vers une cause ovarienne.
(3) Une \cle{induction de l'ovulation} par traitement hormonal est proposée en première intention ; en cas d'échec, une insémination artificielle puis une FIVETE pourront être envisagées.
\end{corrige}

\newpage

\coursec{Activité d'intégration}

\coursec{Leçon 07 : La santé reproductive}

\courssub{Activité d'intégration : Le cas du couple Mbarga}

\courssub{Objectifs de l'activité}

À l'issue de cette activité d'intégration, tu seras capable de :

\begin{itemize}
\item \cle{interpréter} une courbe de dosages hormonaux au cours du cycle ovarien et d'y déceler une anomalie (absence de pic de LH, absence de phase lutéale) ;
\item \cle{analyser} les résultats d'un examen biologique de fertilité masculine (spermogramme) en les comparant aux valeurs de référence de l'OMS ;
\item \cle{identifier} les causes possibles de l'infertilité d'un couple, chez l'homme comme chez la femme ;
\item \cle{choisir} et \cle{justifier} la technique de procréation médicalement assistée (PMA) la plus adaptée à une situation clinique donnée ;
\item \cle{concevoir} un outil de sensibilisation sur les causes évitables de l'infertilité, destiné à un public de jeunes.
\end{itemize}

Cette activité mobilise donc l'ensemble des savoir-faire de la leçon dans une \cle{situation-problème authentique} : le suivi médical d'un couple camerounais confronté à l'infertilité.

\courssub{Situation}

M. et Mme \textbf{Mbarga}, mariés depuis 4 ans à Douala, désirent un enfant. Malgré des rapports sexuels réguliers et non protégés, aucune grossesse n'est survenue. Soumis à une forte pression familiale, le couple décide enfin de consulter ensemble le service de gynécologie-obstétrique de l'Hôpital Laquintinie de Douala. Le médecin prescrit un bilan complet de fertilité. Trois documents du dossier médical te sont confiés.

\textbf{Document 1 — Dosages hormonaux de Mme Mbarga au cours d'un cycle de 28 jours.}

Les taux sanguins de LH (hormone lutéinisante), de FSH (hormone folliculo-stimulante), d'{\oe}stradiol et de progestérone ont été mesurés tous les deux jours pendant un cycle. Les résultats sont résumés dans le tableau et la courbe ci-dessous.

\begin{center}
\begin{tabularx}{\linewidth}{|l|X|X|X|X|}
\hline
\rowcolor{popblueL} \textbf{Hormone} & \textbf{Jour 1} & \textbf{Jour 13} & \textbf{Jour 21} & \textbf{Valeurs attendues (cycle normal 28 j)} \\
\hline
LH (UI/L) & 5 & 6 & 5 & Pic de 25 à 40 UI/L vers J13--J14 \\
\hline
FSH (UI/L) & 6 & 7 & 5 & Pic modéré de 10 à 15 UI/L vers J13 \\
\hline
{\OE}stradiol (pg/mL) & 40 & 90 & 60 & Pic de 200 à 300 pg/mL avant l'ovulation (J12--J13) \\
\hline
Progestérone (ng/mL) & 0,4 & 0,5 & 0,6 & Montée à 10--20 ng/mL en phase lutéale (J21) \\
\hline
\end{tabularx}
\end{center}

\begin{popfigure}
\begin{center}
\begin{tikzpicture}
\begin{axis}[
width=\linewidth, height=7cm,
xlabel={Jour du cycle}, ylabel={Taux sanguin (échelle relative)},
xmin=0, xmax=29, ymin=0, ymax=10,
xtick={0,7,14,21,28}, xticklabels={0, 7, 14, 21, 28},
ytick={0,2,4,6,8,10},
legend style={at={(0.02,0.98)}, anchor=north west, font=\small, fill=white, fill opacity=0.8, draw=popline},
grid=both, grid style={popline!30},
clip=false
]
% Patient data points (markers)
\addplot[only marks, mark=*, popblue, mark size=2.5] coordinates {(1,1.5) (13,1.8) (21,1.5)};
\addlegendentry{LH (Mme Mbarga)}
\addplot[only marks, mark=square*, poporange, mark size=2.5] coordinates {(1,1.2) (13,1.4) (21,1.0)};
\addlegendentry{FSH (Mme Mbarga)}
\addplot[only marks, mark=triangle*, popgreen, mark size=2.5] coordinates {(1,0.8) (13,1.0) (21,0.9)};
\addlegendentry{Progestérone (Mme Mbarga)}
\addplot[only marks, mark=diamond*, poppurple, mark size=2.5] coordinates {(1,0.5) (13,2.5) (21,1.0)};
\addlegendentry{{\OE}stradiol (Mme Mbarga)}

% Theoretical normal LH peak (dashed)
\addplot[popcurve, popblue, dashed, thick, domain=1:28, samples=200] {1.5 + 7*exp(-((x-13)^2)/8)};
\addlegendentry{LH attendue (cycle normal)}

% Annotations
\node[popblue, anchor=south, font=\small] at (axis cs:13,8.5) {Pic LH normal};
\draw[popblue, dashed, thick, ->] (axis cs:13,8.2) -- (axis cs:13,7.5);
\node[popmuted, anchor=north, font=\small, align=center] at (axis cs:14,1.2) {Absence de pic\\chez Mme Mbarga};
\end{axis}
\end{tikzpicture}
\end{center}
\legende{Évolution des taux hormonaux chez Mme Mbarga (symboles pleins) comparée au pic de LH théorique d'un cycle ovulatoire normal (courbe pointillée). Chez la patiente, la LH et la FSH restent basales, l'{\oe}stradiol n'atteint pas le seuil de rétroaction positive, et la progestérone ne monte pas en phase lutéale.}
\end{popfigure}

\textbf{Document 2 — Spermogramme de M. Mbarga.}

L'analyse du sperme, recueilli après \SI{4}{jours} d'abstinence sexuelle, donne les résultats suivants :

\begin{center}
\begin{tabularx}{\linewidth}{|l|X|X|}
\hline
\rowcolor{popblueL} \textbf{Paramètre} & \textbf{Résultat de M. Mbarga} & \textbf{Valeurs de référence (OMS, 2021)} \\
\hline
Volume de l'éjaculat & \SI{3,0}{mL} & $\ge$ \SI{1,5}{mL} \\
\hline
Concentration en spermatozoïdes & \SI{12}{\million\per\mL} & $\ge$ \SI{15}{\million\per\mL} \\
\hline
Mobilité progressive (a+b) & \SI{22}{\%} & $\ge$ \SI{32}{\%} \\
\hline
Formes typiques (morphologie normale) & \SI{6}{\%} & $\ge$ \SI{4}{\%} \\
\hline
\end{tabularx}
\end{center}

\textbf{Document 3 — Hystérosalpingographie de Mme Mbarga.}

La radiographie de l'utérus et des trompes après injection d'un produit de contraste iodé montre une cavité utérine de morphologie normale et des \cle{trompes de Fallope parfaitement perméables} bilatéralement : le produit de contraste s'écoule librement dans la cavité péritonéale des deux côtés (épanchement péritonéal symétrique).

\textbf{Problème posé :} Quelles sont les causes de l'infertilité du couple Mbarga, et quelle technique de procréation médicalement assistée peut-on leur proposer en première intention ?

\courssub{Consignes}

Par groupes de 4 à 6 élèves, vous traiterez les tâches suivantes, puis vous présenterez vos conclusions à la classe.

\begin{enumerate}
\item \textbf{Analyse du document 1.} Décris l'évolution de chaque hormone au cours du cycle de Mme Mbarga. Compare avec un cycle ovulatoire normal. Quel événement majeur du cycle ovarien ne se produit vraisemblablement pas ? Justifie ta réponse en citant \textbf{deux arguments chiffrés} tirés des dosages.
\item \textbf{Diagnostic chez Mme Mbarga.} Nomme précisément l'anomalie observée et explique, à l'aide de tes connaissances sur la régulation du cycle ovarien (axe hypothalamo-hypophyso-ovarien), pourquoi elle empêche la conception.
\item \textbf{Analyse du document 2.} Repère les paramètres anormaux du spermogramme de M. Mbarga. Nomme la (ou les) anomalie(s) correspondante(s) à l'aide du vocabulaire médical approprié.
\item \textbf{Rôle du document 3.} Pourquoi le résultat de l'hystérosalpingographie est-il une information capitale pour le choix de la technique de PMA ?
\item \textbf{Synthèse diagnostique.} L'infertilité du couple est-elle d'origine féminine, masculine ou mixte ? Justifie.
\item \textbf{Choix thérapeutique.} Parmi l'insémination artificielle (IA) et la fécondation in vitro avec transfert d'embryon (FIVETE), propose la technique la plus adaptée à ce couple et justifie ton choix par \textbf{au moins trois arguments} tirés du dossier. Précise quel traitement préalable devra recevoir Mme Mbarga.
\item \textbf{Affiche de sensibilisation.} Conçois une affiche destinée aux élèves de ton lycée, intitulée « Protégeons notre fertilité », présentant au moins quatre causes évitables de l'infertilité et les comportements préventifs correspondants.
\end{enumerate}

\courssub{Ressources à mobiliser}

\begin{aretenir}[Ce qu'il faut avoir en tête pour réussir]
\begin{itemize}
\item \textbf{Cycle ovarien normal (28 jours) :} 
      \begin{itemize}
      \item \textit{Phase folliculaire (J1--J14) :} La FSH stimule la croissance folliculaire $\rightarrow$ sécrétion d'{\oe}stradiol croissante.
      \item \textit{Rétroaction positive :} Le pic d'{\oe}stradiol ($>$ 200 pg/mL) déclenche le \cle{pic de LH} (et de FSH) vers J13--J14.
      \item \textit{Ovulation :} Le pic de LH provoque la rupture du follicule de De Graaf $\rightarrow$ libération de l'ovocyte (J14).
      \item \textit{Phase lutéale (J15--J28) :} Le follicule rompu se luteinise en \cle{corps jaune} $\rightarrow$ sécrétion massive de \cle{progestérone} (10--20 ng/mL vers J21) et d'{\oe}stradiol.
      \end{itemize}
\item \textbf{Signes biologiques d'un cycle anovulatoire :} Absence de pic de LH \textbf{et} taux de progestérone restant bas ($<$ \SI{3}{ng/mL}) en phase lutéale (J21).
\item \textbf{Vocabulaire du spermogramme (référentiel OMS) :} 
      \begin{itemize}
      \item \cle{Oligospermie} (ou oligozoospermie) : concentration $<$ \SI{15}{\million\per\mL} ;
      \item \cle{Asthénospermie} (ou asthénzoospermie) : mobilité progressive $<$ \SI{32}{\%} ;
      \item \cle{Tératospermie} (ou tératozoospermie) : formes typiques $<$ \SI{4}{\%} ;
      \item \cle{Azoospermie} : absence totale de spermatozoïdes dans l'éjaculat.
      \end{itemize}
\item \textbf{Techniques de PMA au programme :}
      \begin{itemize}
      \item \cle{Insémination artificielle (IA)} : Dépôt de spermatozoïdes préparés (capacités, concentrés) dans la cavité utérine au moment de l'ovulation. \textbf{Conditions :} Trompes perméables + Ovulation (spontanée ou induite) + Facteur masculin modéré.
      \item \cle{FIVETE (Fécondation In Vitro et Transfert d'Embryon)} : Fécondation en laboratoire (mise en contact ovocyte/spermatozoïdes ou ICSI) puis transfert de l'embryon (J2--J3 ou J5--J6) dans l'utérus. \textbf{Indications :} Trompes obstruées/absentes, infertilité masculine sévère (oligo-asthéno-tératospermie sévère, azoospermie), échecs répétés d'IA, infertilité inexpliquée.
      \end{itemize}
\item \textbf{Principales causes évitables d'infertilité (contexte camerounais) :}
      \begin{itemize}
      \item Infections sexuellement transmissibles (IST) non traitées ou mal traitées : \emph{Chlamydia trachomatis}, \emph{Neisseria gonorrhoeae} $\rightarrow$ salpingites, épididymites, séquelles obstructives ;
      \item Tabagisme (actif/passif) et consommation excessive d'alcool : altération de la spermatogenèse, diminution de la réserve ovarienne, qualité ovocytaire ;
      \item Retard de consultation / automédication / pratiques traditionnelles à risque ;
      \item Obésité / malnutrition sévère / expositions toxiques professionnelles (pesticides, métaux lourds).
      \end{itemize}
\end{itemize}
\end{aretenir}

\courssub{Démarche et corrigé commenté}

\begin{exoresolu}[Le couple Mbarga : diagnostic étiologique et proposition thérapeutique]
Reprends les sept consignes et rédige, pour chacune, une réponse argumentée, structurée et chiffrée.
\tcblower

\textbf{Étape 1 — Analyse des courbes hormonales (Consignes 1 et 2).}

Sur l'ensemble du cycle observé, les taux de LH et de FSH demeurent bas et stables : LH entre \SI{5}{UI/L} et \SI{6}{UI/L} (J1, J13, J21) ; FSH entre \SI{5}{UI/L} et \SI{7}{UI/L}. Aucun pic brutal n'est visible vers le 13\ieme{} jour. L'{\oe}stradiol, bien que légèrement supérieur au J13 (\SI{90}{pg/mL} vs \SI{40}{pg/mL} au J1), reste très en deçà du seuil physiologique de la rétroaction positive (\SIrange{200}{300}{pg/mL}). 

\emph{Deux arguments chiffrés convergent vers l'absence d'ovulation :}
\begin{enumerate}
\item \textbf{L'absence totale de pic de LH} (max \SI{6}{UI/L} au J13 au lieu de \SIrange{25}{40}{UI/L}) prive l'ovaire du signal déclencheur de la maturation finale de l'ovocyte et de la rupture folliculaire.
\item \textbf{L'absence de montée progestéronique en phase lutéale} : au J21, la progestérone est à \SI{0,6}{ng/mL} alors qu'elle devrait atteindre \SIrange{10}{20}{ng/mL}. Cela prouve l'absence de formation d'un corps jaune fonctionnel, donc l'absence de rupture folliculaire.
\end{enumerate}

\emph{Interprétation physiopathologique :} Il s'agit d'un cycle \cle{anovulatoire}. L'anovulation chez Mme Mbarga traduit un dysfonctionnement de l'axe hypothalamo-hypophyso-ovarien : soit une insuffisance de la sécrétion pulsatile de GnRH (hypothalamus), soit une insuffisance de réponse hypophysaire (LH/FSH), soit une résistance ovarienne aux gonadotrophines. Sans ovocyte libéré dans la trompe, la fécondation est impossible : c'est une cause féminine \textbf{majeure et mécanique} d'infertilité, mais potentiellement réversible par induction pharmacologique de l'ovulation.

\textbf{Étape 2 — Analyse du spermogramme (Consigne 3).}

Deux paramètres sont en-deçà des seuils de référence OMS 2021 :
\begin{itemize}
\item Concentration = \SI{12}{\million\per\mL} $<$ \SI{15}{\million\per\mL} $\rightarrow$ \cle{oligospermie} modérée.
\item Mobilité progressive = \SI{22}{\%} $<$ \SI{32}{\%} $\rightarrow$ \cle{asthénospermie} modérée.
\end{itemize}
Le volume (\SI{3,0}{mL}) et la morphologie (\SI{6}{\%} de formes typiques, seuil à \SI{4}{\%}) sont normaux. M. Mbarga présente donc une \textbf{oligo-asthénospermie modérée}. Sa fertilité est diminuée (moins de spermatozoïdes, moins mobiles) mais non nulle : il produit des gamètes mâles viables et morphologiquement normaux.

\textbf{Étape 3 — Rôle de l'hystérosalpingographie (Consigne 4).}

La perméabilité tubaire bilatérale confirmée par l'hystérosalpingographie est \textbf{la condition sine qua non de l'insémination artificielle (IA)}. En effet, l'IA dépose les spermatozoïdes dans l'utérus ; la fécondation devant se produire dans l'ampoule de la trompe, les trompes doivent être perméables pour permettre la rencontre gamétique et le transport de l'embryon vers l'utérus. Si les trompes étaient bouchées (hydrosalpinx, synéchies), l'IA serait vouée à l'échec et la FIVETE s'imposerait d'emblée. Ce document écarte donc une cause tubaire et valide la faisabilité de l'IA.

\textbf{Étape 4 — Synthèse diagnostique (Consigne 5).}

Le couple Mbarga présente une infertilité \cle{mixte} (ou conjugale) :
\begin{itemize}
\item \textbf{Facteur féminin majeur :} Anovulation (cycle anovulatoire documenté).
\item \textbf{Facteur masculin modéré :} Oligo-asthénospermie.
\end{itemize}
Cette association est fréquente (environ 30 à 40\% des couples infertiles). Elle souligne l'impératif médical d'un \textbf{bilan systématique des deux partenaires} dès le début de la prise en charge, évitant ainsi la stigmatisation injuste de l'un des époux (souvent la femme dans nos contextes socioculturels).

\textbf{Étape 5 — Choix de la technique de PMA (Consigne 6).}

La technique la plus adaptée en \textbf{première intention} est l'\textbf{Insémination Artificielle (IA) avec stimulation ovarienne contrôlée}, pour les trois raisons majeures suivantes :
\begin{enumerate}
\item \textbf{Trompes perméables} (Doc 3) : condition anatomique indispensable à l'IA.
\item \textbf{Anovulation traitable} : L'induction de l'ovulation par citrate de clomifène (anti-œstrogène, 1er choix) ou par gonadotrophines exogènes (FSH recombinante, hMG, 2e choix si échec clomifène) permet d'obtenir une croissance folliculaire monofolliculaire idéale et de déclencher l'ovulation (par injection d'hCG, analogue du pic de LH). Cela corrige la cause féminine majeure.
\item \textbf{Facteur masculin modéré compensable} : La préparation du sperme au laboratoire (gradient de densité, swim-up) permet de sélectionner et concentrer les spermatozoïdes les plus mobiles (récupération de 5 à 20 millions de spermatozoïdes mobiles progressifs), compensant l'oligo-asthénospermie. Le dépôt intra-utérin au moment précis de l'ovulation induite raccourcit le parcours et augmente la probabilité de rencontre.
\end{enumerate}
La FIVETE, plus invasive, plus coûteuse (souvent > \SI{1\,500\,000}{FCFA} au Cameroun vs \SI{150\,000}{FCFA} pour une IA) et source d'hyperstimulation ovarienne, ne sera envisagée qu'en \textbf{deuxième intention} après 3 à 4 échecs d'IA bien conduits, ou si la stimulation ovarienne s'avère inefficace (insuffisance ovarienne).

\textbf{Étape 6 — Affiche de sensibilisation (Consigne 7).}

\begin{exemplebox}[Proposition de structure pour l'affiche « Protégeons notre fertilité »]
\textbf{TITRE : PROTÉGEONS NOTRE FERTILITÉ — AGISSONS AUJOURD'HUI POUR DEMAIN}

\begin{center}
\begin{tabularx}{\linewidth}{|>{\bfseries}X|X|}
\hline
\rowcolor{popblueL} \textbf{CAUSE ÉVITABLE} & \textbf{COMPORTEMENT PRÉVENTIF} \\
\hline
\textbf{1. Infections sexuellement transmissibles (IST)} \\
Chlamydiose, Gonococcie, Syphilis $\rightarrow$ Salpingites, Épididymites, Stérilité tubaire/obstructive. & \textbf{Préservatif masculin/féminin} à chaque rapport à risque. \\
& \textbf{Dépistage systématique} (urines, frottis) chez les deux partenaires avant tout projet de grossesse. \\
& \textbf{Traitement précoce, complet et simultané} des deux partenaires (antibiotiques adaptés). \\
\hline
\textbf{2. Tabac et Alcool} \\
Tabac : toxicité ovocytaire, ménopause précoce, altération ADN spermatozoïdes. \\
Alcool : perturbation axe hormonal, tératogénèse, asthénospermie. & \textbf{Arrêt total du tabac} (sevrage aidé si besoin). \\
& \textbf{Consommation d'alcool modérée ou nulle} (zéro alcool en période de conception). \\
\hline
\textbf{3. Retard de consultation / Automédication / Pratiques à risque} \\
Infections génitales traitées par plantes non dosées, antibiotiques inadaptés $\rightarrow$ chronicisation, séquelles. & \textbf{Consultation précoce} (médecin, sage-femme) devant tout écoulement, douleur pelvienne, brûlure mictionnelle. \\
& \textbf{Respect strict des ordonnances} (doses, durée). Pas d'automédication. \\
\hline
\textbf{4. Mode de vie et Environnement} \\
Obésité / Maigreur extrême (aménorrhée), expositions pesticides/métaux lourds (agriculteurs, soudeurs), stress chronique. & \textbf{Poids santé} (IMC 18,5--25). Activité physique régulière. \\
& \textbf{Protection individuelle} (gants, masques) au travail. \\
& \textbf{Gestion du stress}, sommeil suffisant. \\
\hline
\textbf{5. Âge maternel avancé} \\
Baisse de la réserve ovarienne et qualité ovocytaire après 35 ans. & \textbf{Information sur l'horloge biologique} : ne pas différer indéfiniment le 1er projet de grossesse si possible. \\
& \textbf{Consultation préconceptionnelle} dès 30--35 ans si projet différé. \\
\hline
\end{tabularx}
\end{center}

\textbf{BAS DE L'AFFICHE (Message clé) :}
\begin{center}
\fcolorbox{popblue}{popcream}{\parbox{0.9\linewidth}{\centering
\textbf{APRÈS 12 MOIS} de rapports réguliers, non protégés, sans grossesse : \textbf{CONSULTEZ EN COUPLE} (Gynécologue / Urologue / Biologiste). \\
L'infertilité est un \textbf{problème médical}, pas une faute, pas une malédiction, pas une honte. \\
\textit{« Ensemble, brisons le silence, préservons la vie. »}
}}
\end{center}
\end{exemplebox}
\end{exoresolu}

\courssub{Bilan de l'activité}

Cette enquête autour du couple Mbarga t'a permis de mettre en évidence plusieurs acquis essentiels pour le Baccalauréat et pour ta future vie de citoyen responsable :

\begin{itemize}
\item La lecture croisée des dosages de \cle{LH} (déclencheur) et de \cle{progestérone} (témoin de la luteinisation) est la \cle{clé diagnostique} d'un cycle anovulatoire : pas de pic LH $\rightarrow$ pas d'ovulation $\rightarrow$ pas de corps jaune $\rightarrow$ progestérone basse.
\item L'infertilité d'un couple est souvent \cle{mixte} : le bilan standardisé (dosages hormonaux J3--J5 + progestérone J21 + imagerie tubaire chez la femme ; spermogramme après 3--5 j d'abstinence chez l'homme) doit \textbf{toujours} porter sur les deux partenaires.
\item Le choix entre \cle{Insémination Artificielle} et \cle{FIVETE} repose sur un algorithme clinique rigoureux : \textbf{Perméabilité tubaire ?} Oui $\rightarrow$ IA possible. \textbf{Ovulation induisible ?} Oui $\rightarrow$ IA possible. \textbf{Facteur masculin sévère ?} Non $\rightarrow$ IA en 1ère intention. FIVETE = recours en cas d'obstruction tubaire, de facteur masculin sévère (ICSI), ou d'échecs répétés.
\item De nombreuses causes d'infertilité sont \cle{évitables} : la prévention primaire (préservatif, hygiène de vie, vaccination VHB/VPH) et la prévention secondaire (dépistage/traitement précoce des IST, consultation préconceptionnelle) sont des messages de santé publique majeurs que tu peux relayer efficacement auprès de tes pairs.
\end{itemize}

Tu as ainsi mobilisé, dans une démarche d'investigation scientifique complète (observation $\rightarrow$ hypothèses $\rightarrow$ analyse de documents $\rightarrow$ interprétation $\rightarrow$ conclusion $\rightarrow$ communication), l'ensemble des savoir-faire de la leçon : identifier les hormones sexuelles, interpréter des résultats expérimentaux et biologiques, décrire les techniques de PMA et élaborer un outil de sensibilisation.

\begin{attention}[Pièges fréquents au Bac — Vigilance !]
\begin{itemize}
\item \textbf{Confondre FIVETE et IA :} L'IA = fécondation \textbf{in vivo} (dans la trompe) après dépôt utérin. La FIVETE = fécondation \textbf{in vitro} (en laboratoire). Ne pas écrire « fécondation in vitro » pour l'IA.
\item \textbf{Oublier l'unité :} Concentration en \SI{}{\million\per\mL} (pas « millions » tout court). Taux hormonaux en UI/L, pg/mL, ng/mL. Toujours écrire l'unité.
\item \textbf{Diagnostiquer l'anovulation sur un seul argument :} Il faut \textbf{deux} arguments : absence de pic LH \textbf{ET} progestérone basse en phase lutéale. Un taux d'{\oe}stradiol bas seul ne suffit pas (peut être bas en début de cycle).
\item \textbf{Vocabulaire spermogramme :} « Oligospermie » = concentration basse. « Asthénospermie » = mobilité basse. « Tératospermie » = morphologie basse. Ne pas les confondre. « Normospermie » = tout normal.
\item \textbf{Indication FIVETE :} Ne pas proposer la FIVETE pour une simple oligo-asthénospermie modérée avec trompes perméables. C'est une indication d'IA. La FIVETE (ou ICSI) est pour les cas sévères ou tubaires.
\end{itemize}
\end{attention}

\begin{savaistu}[Santé reproductive au Cameroun : Repères contextuels]
\begin{itemize}
\item \textbf{Prévalence :} L'infertilité touche environ \SI{15}{\%} à \SI{20}{\%} des couples en âge de procréer au Cameroun (données OMS/Ministère de la Santé Publique). L'infertilité secondaire (après une 1ère grossesse) est fréquente, souvent post-infectieuse (IST mal soignées, avortements à risque).
\item \textbf{Accès à la PMA :} Quelques centres spécialisés existent (Yaoundé : CHU, Hôpital Central, cliniques privées ; Douala : CHU, Laquintinie, cliniques privées). Le coût reste un frein majeur (non remboursé par la CNPS/assurance maladie). L'IA est plus accessible financièrement.
\item \textbf{Cadre légal :} La loi camerounaise encadre la PMA (Loi n°2002/004 du 19 avril 2002 relative à la procréation médicalement assistée). Don de gamètes autorisé sous conditions strictes (anonymat, non-rémunération). Gestation pour autrui (GPA) interdite.
\item \textbf{Prévention IST :} Le Cameroun a intégré la prise en charge syndromique des IST dans les formations sanitaires de base. Le dépistage du VIH/syphilis est systématique en consultation prénatale (PTME). Le vaccin anti-VPH (cancer du col) est déployé chez les jeunes filles (9--14 ans) : prévention majeure de l'infertilité par conservation du col utérin.
\end{itemize}
\end{savaistu}

\newpage

\coursec{Méthodes et exercices résolus}

\coursec{Rappels et mise en place expérimentale : les hormones sexuelles}

\begin{definition}[Hormones sexuelles et gonades]
Les \cle{hormones sexuelles} sont des messagers chimiques produits principalement par les \cle{gonades} (testicules chez l'homme, ovaires chez la femme) et par l'axe hypothalamo-hypophysaire. Elles assurent le développement des caractères sexuels, la gamétogenèse et, chez la femme, le cycle ovarien et la gestation.
\end{definition}

\begin{propriete}[Classification des hormones de la reproduction]
\begin{center}
\begin{tabularx}{\linewidth}{|l|X|X|X|}\hline
\rowcolor{popblueL} \textbf{Hormone} & \textbf{Glande source} & \textbf{Nature chimique} & \textbf{Cible principale / Rôle clé} \\ \hline
GnRH (Gonadolibérine) & Hypothalamus & Peptide (protéique) & Hypophyse antérieure : stimule la sécrétion de FSH et LH \\ \hline
FSH (Folliculo-stimulante) & Hypophyse antérieure & Glycoprotéique & \textbf{Homme :} Tubes séminifères (cellules de Sertoli) $\rightarrow$ spermatogenèse \\
& & & \textbf{Femme :} Follicules ovariens $\rightarrow$ croissance, sécrétion d'\oe{}strogènes \\ \hline
LH (Lutéinisante) & Hypophyse antérieure & Glycoprotéique & \textbf{Homme :} Cellules de Leydig $\rightarrow$ sécrétion de testostérone \\
& & & \textbf{Femme :} Ovulation (pic LH), luteinisation $\rightarrow$ corps jaune (progestérone) \\ \hline
Testostérone & Testicules (cellules de Leydig) & Stéroïdienne (cholestérol) & Caractères sexuels secondaires, spermatogenèse (synergie FSH), rétrocontrôle $-$ \\ \hline
\oe{}strogènes (principalement \oe{}stradiol) & Ovaires (follicule, cellules de la thèque + granulosa) & Stéroïdienne & Prolifération endomètre, rétrocontrôle $-$ (faible) puis $+$ (pic), caractères sexuels \\ \hline
Progestérone & Ovaires (corps jaune) puis Placenta & Stéroïdienne & Endomètre sécrétoire, maintien grossesse, rétrocontrôle $-$ (phase lutéale), thermogenèse \\ \hline
Inhibine & Testicules (Sertoli) / Ovaires (granulosa) & Protéique & Rétrocontrôle $-$ sélectif sur la FSH (pas sur LH) \\ \hline
\end{tabularx}
\end{center}
\end{propriete}

\begin{methode}[Identifier et classer les hormones sexuelles (Savoir-faire 1)]
\begin{enumerate}
\item \textbf{Repérer l'organe producteur} : hypothalamus (GnRH), hypophyse antérieure (FSH, LH), gonades (testostérone chez l'homme ; \oe{}strogènes et progestérone chez la femme).
\item \textbf{Classer par nature chimique} : hormones \cle{protéiques/peptidiques} (GnRH, FSH, LH, inhibine — récepteurs membranaires, seconde messager cAMP) et hormones \cle{stéroïdiennes} (testostérone, \oe{}strogènes, progestérone — diffusion membranaire, récepteurs nucléaires intracellulaires).
\item \textbf{Associer chaque hormone à sa cible et à son effet} : FSH $\rightarrow$ tubes séminifères (spermatogenèse) ou follicules ovariens ; LH $\rightarrow$ cellules de Leydig (testostérone) ou ovulation/corps jaune ; testostérone $\rightarrow$ caractères sexuels secondaires et spermatogenèse ; \oe{}strogènes $\rightarrow$ prolifération de l'endomètre ; progestérone $\rightarrow$ endomètre sécrétoire, maintien de la grossesse.
\item \textbf{Utiliser un tableau de synthèse} (hormone / glande / nature / cible / effet) : c'est la présentation la plus sûre en copie d'examen.
\end{enumerate}
\end{methode}

\begin{exoresolu}[QCM de synthèse sur les hormones sexuelles]
Pour chaque affirmation, répondre par vrai ou faux en justifiant.
\begin{enumerate}
\item La testostérone est sécrétée par l'hypophyse.
\item La FSH stimule la croissance des follicules ovariens chez la femme.
\item La progestérone est une hormone protéique.
\item La GnRH est commune à l'homme et à la femme.
\end{enumerate}
\tcblower
\textbf{Solution détaillée.}
\begin{enumerate}
\item \textbf{Faux.} La testostérone est une hormone stéroïdienne sécrétée par les \cle{cellules de Leydig} des testicules (tissu interstitiel), sous le contrôle de la LH hypophysaire. L'hypophyse sécrète les gonadostimulines FSH et LH, pas de stéroïdes.
\item \textbf{Vrai.} La FSH (hormone folliculo-stimulante) agit sur les follicules ovariens : elle stimule leur croissance et leur maturation pendant la phase folliculaire du cycle ovarien. Chez l'homme, elle agit sur les tubes séminifères (spermatogenèse via les cellules de Sertoli).
\item \textbf{Faux.} La progestérone est une hormone \cle{stéroïdienne} (dérivée du cholestérol), sécrétée par le corps jaune puis par le placenta. Les hormones protéiques de la régulation sont la GnRH, la FSH et la LH.
\item \textbf{Vrai.} La GnRH (gonadolibérine), neurohormone hypothalamique, stimule la sécrétion de FSH et de LH chez les deux sexes ; seuls les organes cibles et les gonadostéroïdes produits diffèrent.
\end{enumerate}
\textbf{Conclusion :} la distinction glande productrice / nature chimique / organe cible permet de lever toute ambiguïté.
\end{exoresolu}

\begin{savaistu}[Contexte camerounais : Dépistage hormonal]
Au Cameroun, les dosages hormonaux (FSH, LH, testostérone, progestérone, \oe{}stradiol, AMH) sont réalisés dans les laboratoires des hôpitaux de référence (HGOPED, HGY, CHU Yaoundé/Douala) et certains laboratoires privés. Ils sont prescrits dans le bilan d'infertilité du couple, le suivi de la puberté, ou le diagnostic de pathologies hypophysaires. La connaissance des valeurs de référence selon l'âge et la phase du cycle est indispensable pour l'interprétation.
\end{savaistu}

\coursec{Régulation des hormones sexuelles chez l'homme}

\begin{propriete}[Axe hypothalamo-hypophyso-testiculaire chez l'homme]
La régulation est un \cle{rétrocontrôle négatif} simple (pas de rétrocontrôle positif chez l'homme adulte).
\begin{itemize}
\item \textbf{Hypothalamus} $\xrightarrow{\text{GnRH (pulsatile)}}$ \textbf{Hypophyse antérieure} $\xrightarrow{\text{FSH + LH}}$ \textbf{Testicules}.
\item \textbf{FSH} $\rightarrow$ Cellules de Sertoli (tubes séminifères) $\rightarrow$ Spermatogenèse + sécrétion d'\textbf{inhibine}.
\item \textbf{LH} $\rightarrow$ Cellules de Leydig (interstitielles) $\rightarrow$ Sécrétion de \textbf{testostérone}.
\item \textbf{Testostérone} : effets périphériques (caractères sexuels, anabolisme, libido) + \textbf{rétrocontrôle négatif} sur hypothalamus (GnRH $\downarrow$) et hypophyse (LH $\downarrow$, FSH $\downarrow$).
\item \textbf{Inhibine} : rétrocontrôle négatif \emph{sélectif} sur la FSH hypophysaire (affinage de la spermatogenèse).
\end{itemize}
\end{propriete}

\begin{experience}[Expériences classiques : castration, greffe, injection chez le rat (Savoir-faire 2)]
\textbf{Objectif :} Mettre en évidence l'origine et le rôle de la testostérone et le rétrocontrôle.
\textbf{Matériel :} Rats mâles, lots témoins et expérimentaux.
\textbf{Protocole et résultats types :}
\begin{center}
\begin{tabularx}{\linewidth}{|l|X|X|X|}\hline
\rowcolor{popblueL} Lot & Traitement & Vésicules séminales / Caractères sexuels & Taux sanguin de LH \\ \hline
A (Témoin) & Intact & Normaux & Normal \\ \hline
B & Castré (ablation testicules) & Atrophiés & \textbf{Élevé} (perte rétrocontrôle $-$ ) \\ \hline
C & Castré + Injection testostérone & \textbf{Restaurés} & \textbf{Faible} (rétrocontrôle $-$ rétabli) \\ \hline
D & Intact + Forte dose testostérone & Normaux (saturés) & \textbf{Très faible} (rétrocontrôle $-$ maximal) \\ \hline
\end{tabularx}
\end{center}
\textbf{Interprétation :}
\begin{enumerate}
\item Comparaison A/B : La castration supprime la source de testostérone $\rightarrow$ atrophie des organes cibles. La LH augmente : preuve du rétrocontrôle négatif de la testostérone sur l'hypophyse.
\item Comparaison B/C : L'injection de testostérone restaure les caractères sexuels (effet périphérique) et abaisse la LH (effet central rétrocontrôle négatif). \textbf{La spermatogenèse n'est pas restaurée} car elle nécessite l'action locale combinée de FSH (absente car inhibée par la testostérone injectée) et d'une concentration testiculaire en testostérone très élevée (non atteinte par voie sanguine).
\item Lot D : L'excès de testostérone exogène effondre la LH endogène $\rightarrow$ atrophie testiculaire et arrêt de la spermatogenèse (principe de la contraception masculine hormonale et risque du dopage aux stéroïdes anabolisants).
\end{enumerate}
\textbf{Conclusion :} La testostérone a un double rôle : (1) effet périphérique (maintien caractères sexuels, organes annexes) ; (2) effet central inhibiteur (rétroaction négative sur axe hypothalamo-hypophysaire).
\end{experience}

\begin{methode}[Exploiter une expérience de castration / greffe / injection (Savoir-faire 2)]
\begin{enumerate}
\item \textbf{Décrire les protocoles} en comparant les lots : lot témoin (intact) et lots expérimentaux (castré, castré + injection d'hormone, greffé...). Une seule variable doit changer entre deux lots comparés.
\item \textbf{Comparer les résultats} lot par lot : castration $\rightarrow$ disparition des caractères sexuels et arrêt de la gamétogenèse ; injection d'extrait testiculaire $\rightarrow$ restauration.
\item \textbf{Raisonner par élimination} : si l'ablation de l'hypophyse supprime la sécrétion testiculaire mais que l'injection de LH la restaure, alors l'hypophyse contrôle le testicule via la LH.
\item \textbf{Mettre en évidence le rétrocontrôle} : injection de fortes doses de testostérone $\rightarrow$ chute de FSH/LH $\Rightarrow$ rétroaction négative du stéroïde sur l'axe hypothalamo-hypophysaire.
\item \textbf{Conclure} en reliant chaque effet à une hormone précise, et schématiser la boucle de régulation avec flèches $+$ (stimulation) et $-$ (inhibition).
\end{enumerate}
\end{methode}

\begin{popfigure}
\begin{tikzpicture}[node distance=1.8cm, every node/.style={font=\small, align=center}]
\node[draw=popblue, fill=popblueL, rounded corners, minimum width=3.8cm, align=center] (hypoth){Hypothalamus\\GnRH (pulsatile)};
\node[draw=poppurple, fill=poppurpleL, rounded corners, minimum width=3.8cm, below=of hypoth, align=center] (hypop){Hypophyse antérieure\\FSH \quad LH};
\node[draw=poporange, fill=poporangeL, rounded corners, minimum width=3.8cm, below=of hypop, align=center] (testic){Testicules\\Cellules de Leydig : Testostérone\\Cellules de Sertoli : Inhibine, Spermatogenèse};
\node[draw=popgreen, fill=popgreenL, rounded corners, minimum width=3.5cm, right=2.5cm of testic, align=center] (periph){Cibles périphériques\\Caractères sexuels secondaires\\Métabolisme, Libido};
\draw[-{Stealth}, thick, popblue] (hypoth) -- node[right]{$+$ GnRH} (hypop);
\draw[-{Stealth}, thick, poppurple] (hypop) -- node[right]{$+$ FSH / LH} (testic);
\draw[-{Stealth}, thick, poporange] (testic) -- node[above]{$+$ Testostérone} (periph);
\draw[-{Stealth}, thick, poppink] (testic.west) .. controls +(-2.5,1.5) and +(-2.5,-1.5) .. node[left, xshift=-1.5cm]{$-$ Testostérone} (hypoth.west);
\draw[-{Stealth}, thick, poppink] (testic.west) ++(-0.4,0.2) .. controls +(-1.6,0.8) and +(-1.6,-0.8) .. node[left, xshift=-1.2cm]{$-$ Testostérone} (hypop.west);
\draw[-{Stealth}, thick, poppink, dashed] (testic.west) ++(-0.4,-0.2) .. controls +(-1.2,0.4) and +(-1.2,-0.4) .. node[left, xshift=-1.0cm]{$-$ Inhibine} (hypop.west);
\end{tikzpicture}
\legende{Boucle de régulation hormonale chez l'homme : double rétrocontrôle négatif de la testostérone (sur hypothalamus et hypophyse) et rétrocontrôle négatif sélectif de l'inhibine (sur FSH seulement).}
\end{popfigure}

\coursec{Régulation des hormones sexuelles chez la femme}

\begin{propriete}[Axe hypothalamo-hypophyso-ovarien et cycle ovarien]
La régulation féminine est cyclique (environ 28 jours) et fait intervenir un \cle{rétrocontrôle négatif} permanent et un \cle{rétrocontrôle positif} transitoire (unique en physiologie humaine).
\begin{itemize}
\item \textbf{Phase folliculaire (J1--J14)} : FSH $\rightarrow$ recrutement/croissance follicules $\rightarrow$ sécrétion croissante d'\oe{}strogènes (par cellules granulosa + thèque).
\begin{itemize}
\item \oe{}strogènes bas/moyens $\rightarrow$ rétrocontrôle $-$ sur FSH/LH (sélection du follicule dominant).
\item \oe{}strogènes élevés ($> 200$ pg/mL) \textbf{pendant 48 h} $\rightarrow$ rétrocontrôle \textbf{positif} sur hypothalamus/hypophyse $\rightarrow$ \textbf{pic de LH (et FSH)} au J14.
\end{itemize}
\item \textbf{Ovulation (J14 $\pm$ 1)} : Pic de LH $\rightarrow$ rupture follicule de De Graaf $\rightarrow$ libération ovocyte II + cellule polaire.
\item \textbf{Phase lutéale (J14--J28)} : Follicule rompu $\rightarrow$ \textbf{corps jaune} (sous action LH) $\rightarrow$ sécrétion massive de \textbf{progestérone} + \oe{}strogènes.
\begin{itemize}
\item Progestérone + \oe{}strogènes $\rightarrow$ rétrocontrôle $-$ fort sur hypothalamus/hypophyse $\rightarrow$ FSH/LH bas (blocage nouvelle ovulation).
\item Progestérone $\rightarrow$ endomètre sécrétoire (préparation nidation).
\end{itemize}
\item \textbf{Régression du corps jaune (J24--J28 si pas grossesse)} : Chute progestérone/\oe{}strogènes $\rightarrow$ levée rétrocontrôle $-$ $\rightarrow$ remontée FSH (début cycle suivant) + nécrose endomètre $\rightarrow$ \textbf{règles}.
\end{itemize}
\end{propriete}

\begin{methode}[Lire et interpréter les courbes du cycle ovarien et utérin (Savoir-faire 3)]
\begin{enumerate}
\item \textbf{Identifier les axes} : abscisse = jours du cycle (0 à 28) ; ordonnées = concentrations de FSH, LH, \oe{}strogènes, progestérone (unités : UI/L ou pg/mL ou ng/mL).
\item \textbf{Repérer le pic de LH} vers le 14\textsuperscript{e} jour : il déclenche l'\cle{ovulation} (c'est le seul événement daté avec certitude).
\item \textbf{Découper le cycle} : phase folliculaire (J1--J14, dominance \oe{}strogènes, prolifération de l'endomètre) et phase lutéale (J14--J28, dominance progestérone, endomètre sécrétoire).
\item \textbf{Expliquer les rétroactions} : \oe{}strogènes modérés $\rightarrow$ rétrocontrôle négatif ; \oe{}strogènes élevés prolongés ($>200$ pg/mL pendant 48 h) $\rightarrow$ rétrocontrôle \emph{positif} déclenchant le pic de LH ; progestérone $\rightarrow$ rétrocontrôle négatif.
\item \textbf{Relier ovaire et utérus} : chute des deux stéroïdes en fin de cycle $\rightarrow$ nécrose de l'endomètre $\rightarrow$ règles.
\end{enumerate}
\end{methode}

\begin{exoresolu}[Analyse d'un cycle de 28 jours (Savoir-faire 3)]
Le document ci-dessous donne les variations des taux plasmatiques de LH et de progestérone au cours d'un cycle de 28 jours.

\begin{center}
\begin{tikzpicture}
\begin{axis}[width=0.9\linewidth,height=6cm,xlabel={Jours du cycle},ylabel={Concentration (unités arbitraires)},xmin=0,xmax=28,ymin=0,ymax=10,xtick={0,7,14,21,28},legend style={at={(0.02,0.98)},anchor=north west,font=\small},grid=both,grid style={popline!40}]
\addplot[popcurve,popblue,domain=0:28,samples=200]{1.5+7*exp(-((x-14)^2)/2)};
\addlegendentry{LH (pic J14)}
\addplot[popcurve,poporange,domain=0:28,samples=200]{0.5+6*exp(-((x-21)^2)/18)};
\addlegendentry{Progestérone (pic J21)}
\end{axis}
\end{tikzpicture}
\end{center}

1. Dater l'ovulation en justifiant. 2. Expliquer l'origine de la progestérone après le 14\textsuperscript{e} jour. 3. Prévoir l'évolution de l'endomètre entre le 14\textsuperscript{e} et le 28\textsuperscript{e} jour.
\tcblower
\textbf{Solution détaillée.}
\begin{enumerate}
\item La courbe de LH présente un \cle{pic brutal} centré au 14\textsuperscript{e} jour. Or le pic de LH (déclenché par le rétrocontrôle positif des \oe{}strogènes sécrétés par le follicule mûr) est la cause directe de la rupture du follicule de De Graaf. \textbf{L'ovulation a donc lieu au 14\textsuperscript{e} jour}, environ 24 à 36 h après le début du pic.
\item Après l'ovulation, le follicule rompu se transforme en \cle{corps jaune} sous l'action de la LH résiduelle. Le corps jaune sécrète abondamment la progestérone, d'où la montée de la courbe orange entre J14 et J21 (pic vers J21), puis sa chute quand le corps jaune régresse (en l'absence de fécondation, vers J24--J26).
\item Entre J14 et J28, sous l'action de la progestérone, l'endomètre devient \cle{sécrétoire} : épaississement, vascularisation, développement des glandes utérines sécrétant le glycogène — préparation à la nidation. La chute de progestérone en fin de cycle supprime ce maintien : l'endomètre fonctionnel se nécrose et s'élimine : ce sont les \textbf{règles} (J28 $\rightarrow$ J1 du cycle suivant).
\end{enumerate}
\textbf{Conclusion :} le cycle utérin est entièrement piloté par les sécrétions ovariennes, elles-mêmes contrôlées par l'axe hypothalamo-hypophysaire.
\end{exoresolu}

\begin{methode}[Construire le schéma fonctionnel d'une boucle de rétrocontrôle (Savoir-faire 4)]
\begin{enumerate}
\item \textbf{Placer les trois étages} verticalement : hypothalamus (GnRH) $\rightarrow$ hypophyse antérieure (FSH, LH) $\rightarrow$ gonade (stéroïdes + gamètes).
\item \textbf{Fléches montantes} = stimulations, annotées « $+$ » et nommant l'hormone transportée par le sang.
\item \textbf{Fléches descendantes de retour} = rétroactions, annotées « $-$ » (inhibition) ou « $+$ » (cas unique du rétrocontrôle positif des \oe{}strogènes avant l'ovulation).
\item \textbf{Indiquer les effets périphériques} des stéroïdes (caractères sexuels, endomètre) en dehors de la boucle.
\item \textbf{Vérifier la cohérence} : toute boucle fermée doit être stable ; préciser si le rétrocontrôle est court (sur l'hypothalamus) ou long (sur hypothalamus + hypophyse).
\end{enumerate}
\end{methode}

\begin{exoresolu}[Schéma de la régulation chez la femme : phase folliculaire tardive (rétrocontrôle positif) et phase lutéale (rétrocontrôle négatif)]
Construire les deux schémas fonctionnels contrastés.
\tcblower
\textbf{Solution détaillée.}

\textbf{A. Phase folliculaire tardive (vers J12--J14) : Rétrocontrôle POSITIF des \oe{}strogènes}
\begin{center}
\begin{tikzpicture}[node distance=1.6cm, every node/.style={font=\small, align=center}]
\node[draw=popblue, fill=popblueL, rounded corners, minimum width=3.6cm] (hypoth) {Hypothalamus : GnRH};
\node[draw=poppurple, fill=poppurpleL, rounded corners, minimum width=3.6cm, below=of hypoth] (hypop) {Hypophyse : FSH + LH};
\node[draw=poporange, fill=poporangeL, rounded corners, minimum width=3.6cm, below=of hypop, align=center] (ovaire){Ovaire (Follicule mûr)\\ \oe{}strogènes $\uparrow\uparrow$};
\draw[-{Stealth}, thick, popblue] (hypoth) -- node[right]{$+$ GnRH} (hypop);
\draw[-{Stealth}, thick, poppurple] (hypop) -- node[right]{$+$ FSH / LH} (ovaire);
\draw[-{Stealth}, thick, popgreen!70!black] (ovaire.east) .. controls +(2,1) and +(2,-1) .. node[right]{$+$ \oe{}strogènes (seuil haut)} (hypop.east);
\draw[-{Stealth}, thick, popgreen!70!black] (ovaire.east) ++(0.1,0.3) .. controls +(2.5,1.5) and +(2.5,-1.5) .. node[right, xshift=1.5cm]{$+$ \oe{}strogènes} (hypoth.east);
\legende{Rétrocontrôle positif : les \oe{}strogènes élevés et durables stimulent l'axe $\rightarrow$ pic de LH $\rightarrow$ ovulation.}
\end{tikzpicture}
\end{center}

\textbf{B. Phase lutéale (J15--J28) : Rétrocontrôle NÉGATIF de la progestérone (et \oe{}strogènes)}
\begin{center}
\begin{tikzpicture}[node distance=1.6cm, every node/.style={font=\small, align=center}]
\node[draw=popblue, fill=popblueL, rounded corners, minimum width=3.6cm] (hypoth) {Hypothalamus : GnRH};
\node[draw=poppurple, fill=poppurpleL, rounded corners, minimum width=3.6cm, below=of hypoth] (hypop) {Hypophyse : FSH + LH (bas)};
\node[draw=poporange, fill=poporangeL, rounded corners, minimum width=3.6cm, below=of hypop, align=center] (ovaire){Ovaire (Corps jaune)\\ Progestérone $\uparrow\uparrow$ + \oe{}strogènes};
\node[draw=popgreen, fill=popgreenL, rounded corners, minimum width=3.2cm, right=2.2cm of ovaire] (uterus) {Endomètre sécrétoire};
\draw[-{Stealth}, thick, popblue] (hypoth) -- node[right]{$+$ GnRH} (hypop);
\draw[-{Stealth}, thick, poppurple] (hypop) -- node[right]{$+$ LH (maintien CJ)} (ovaire);
\draw[-{Stealth}, thick, poporange] (ovaire) -- node[above]{$+$ Progestérone} (uterus);
\draw[-{Stealth}, thick, poppink] (ovaire.west) .. controls +(-2.2,1.2) and +(-2.2,-1.2) .. node[left]{$-$ Progestérone / \oe{}strogènes} (hypoth.west);
\draw[-{Stealth}, thick, poppink] (ovaire.west) ++(-0.4,0.15) .. controls +(-1.4,0.7) and +(-1.4,-0.7) .. node[left, xshift=-1.2cm]{$-$} (hypop.west);
\legende{Phase lutéale : la progestérone (et \oe{}strogènes) du corps jaune inhibe l'hypothalamus et l'hypophyse (flèches $-$), maintenant FSH et LH à taux bas (pas de nouvelle ovulation), tout en stimulant l'endomètre ($+$).}
\end{tikzpicture}
\end{center}

\textbf{Conclusion :} l'alternance rétrocontrôle positif (ovulation) / négatif (phase lutéale) structure le cycle. La régression du corps jaune lève l'inhibition et permet le redémarrage cyclique (remontée FSH).
\end{exoresolu}

\coursec{L'infertilité : causes possibles chez l'homme et chez la femme}

\begin{definition}[Infertilité]
L'\cle{infertilité} (ou stérilité) d'un couple est définie par l'absence de grossesse après 12 à 24 mois de rapports sexuels réguliers (2 à 3 par semaine) et non protégés. Elle concerne environ 15 à 20 \% des couples en âge de procréer. La répartition étiologique approximative est : 1/3 causes masculines, 1/3 causes féminines, 1/3 causes mixtes ou idiopathiques (inexpliquées).
\end{definition}

\begin{propriete}[Causes principales de l'infertilité classées par mécanisme]
\begin{center}
\begin{tabularx}{\linewidth}{|l|X|X|}\hline
\rowcolor{popblueL} \textbf{Mécanisme} & \textbf{Chez l'homme} & \textbf{Chez la femme} \\ \hline
\textbf{Production des gamètes} & \textbf{Azoospermie} (absence spermatozoïdes) ou \textbf{Oligospermie} (nombre $< 15 \times 10^6$/mL) : causes génétiques (microdélétion Y, Klinefelter), hormonales (hypogonadisme), varicocèle, cryptorchidie non opérée, toxiques (tabac, alcool, chaleur, chimiothérapie). & \textbf{Anovulation} (absence ovulation) : Syndrome des ovaires polykystiques (SOPK, 1\textsuperscript{re} cause), insuffisance ovarienne précoce, hyperprolactinémie, troubles thyroïdiens, anorexie/obésité, stress intense. \\ \hline
\textbf{Transport / Rencontre des gamètes} & Obstruction des voies excrétrices (épididymes, canaux déférents) : séquelles d'IST (chlamydiae, gonocoques), malformations congénitales (absence canaux déférents chez mucoviscidose), vasectomie. & \textbf{Obstruction tubaire} (trompes de Fallope) : \textbf{1\textsuperscript{re} cause féminine évitable}. Séquelles d'IST (salpingites), endométriose, adhérences post-chirurgicales, ligature des trompes. \\ \hline
\textbf{Qualité / Fonction des gamètes} & Asthénospermie (mobilité réduite), tératospermie (morphologie anormale), fragmentation ADN spermatique. & Mauvaise qualité ovocytaire (âge maternel avancé $> 35$ ans), zona pellucida anormal. \\ \hline
\textbf{Nidation / Milieu utérin} & -- & Anomalies utérines (fibromes sous-muqueux, synéchies, malformations müllériennes), endométriose (adénomyose), endomètre atrophique ou non réceptif. \\ \hline
\textbf{Immunologique / Inexpliquée} & Anticorps anti-spermatozoïdes. & Anticorps anti-spermatozoïdes, facteurs immunologiques complexes. \\ \hline
\end{tabularx}
\end{center}
\end{propriete}

\begin{attention}[Facteurs de risque évitables au Cameroun]
\begin{itemize}
\item \textbf{IST non traitées ou mal traitées :} Chlamydia trachomatis, Neisseria gonorrhoeae $\rightarrow$ salpingites (femme) / épididymites (homme) $\rightarrow$ obstruction définitive. \textbf{Prévention :} dépistage systématique, traitement précoce du couple (prescription partenaire), préservatif.
\item \textbf{Avortements clandestins / non sécurisés :} cause majeure d'infertilité tubaire et utérine (infections, perforations, synéchies).
\item \textbf{Mode de vie :} Tabac (baisse qualité sperme, ménopause précoce), alcool, obésité, exposition professionnelle aux pesticides/métaux lourds (agriculture cotonnière/cacaoyère).
\item \textbf{Âge maternel :} fertilité féminine maximale 20--25 ans, chute nette après 35 ans, critique après 40 ans. Sensibilisation à la planification familiale.
\end{itemize}
\end{attention}

\begin{methode}[Concevoir un message de sensibilisation rigoureux (Savoir-faire 5)]
\begin{enumerate}
\item \textbf{Définir l'infertilité} : absence de grossesse après 12 à 24 mois de rapports non protégés ; elle concerne le \emph{couple} (environ 1/3 causes masculines, 1/3 féminines, 1/3 mixtes ou inexpliquées).
\item \textbf{Classer les causes} en deux colonnes (homme / femme) et par mécanisme : production de gamètes, transport, rencontre gamètes / nidation.
\item \textbf{Identifier les facteurs évitables} : IST non traitées (chlamydiae, gonococcies $\rightarrow$ obstruction des trompes ou des épididymes), varicocèle, tabac, alcool, âge, perturbateurs endocriniens.
\item \textbf{Formuler des messages} clairs, non culpabilisants, avec un appel à l'action (dépistage et traitement des IST, consultation au centre de santé).
\item \textbf{Adapter le support} au public visé (affiche, sketch, spot radio en langue locale : français, anglais, fulfulde, ewondo, bassa...).
\end{enumerate}
\end{methode}

\begin{exoresolu}[Conception d'une affiche de sensibilisation : « L'infertilité : comprendre pour prévenir »]
\textbf{Titre :} L'infertilité : comprendre pour prévenir.

\textbf{Qu'est-ce que l'infertilité ?} Un couple est dit infertile lorsqu'aucune grossesse ne survient après 12 mois de rapports réguliers non protégés. Elle touche environ 15 à 20 \% des couples et concerne l'homme autant que la femme.

\textbf{Chez l'homme :} azoospermie ou oligospermie (production absente ou insuffisante : varicocèle, infections des voies génitales, anomalie hormonale, toxiques) ; obstruction des épididymes ou canaux déférents (séquelle d'IST) ; anomalies de la mobilité (asthénospermie) ou morphologie (tératospermie).

\textbf{Chez la femme :} anovulation (troubles hormonaux, SOPK) ; \textbf{obstruction des trompes de Fallope} (séquelles d'IST non traitées, avortements non sécurisés = première cause évitable) ; anomalies de l'utérus ou de l'endomètre (fibromes, synéchies, endométriose) empêchant la nidation.

\textbf{Prévenir, c'est possible :}
\begin{itemize}
\item Se faire dépister et traiter rapidement les IST (centres de santé, CEDC) ; traitement obligatoire du partenaire.
\item Utiliser le préservatif (double protection : IST + grossesses non désirées).
\item Éviter tabac, alcool excessif, automédication hormonale, exposition aux toxiques.
\item Consulter tôt en cas de douleurs pelviennes, règles irrégulières, absence de grossesse désirée.
\item Ne pas retarder indéfiniment le projet de grossesse (fertilité féminine diminue nettement après 35 ans).
\end{itemize}

\textbf{Message final :} L'infertilité n'est ni une fatalité ni une honte : parlons-en et consultons ensemble au centre de santé le plus proche. \textbf{Ensemble, brisons le silence !}
\end{exoresolu}

\coursec{La procréation médicalement assistée (PMA)}

\begin{definition}[Procréation Médicalement Assistée (PMA)]
Ensemble des techniques cliniques et biologiques permettant d'obtenir une grossesse en manipulant les gamètes et/ou les embryons, lorsque la conception naturelle est impossible ou a échoué. Au Cameroun, la PMA est pratiquée dans quelques centres spécialisés (ex: Centre de Procréation Médicalement Assistée de l'Hôpital Central de Yaoundé, cliniques privées à Douala/Yaoundé) avec un cadre éthique strict (loi sur la santé de la reproduction).
\end{definition}

\begin{propriete}[Les deux techniques majeures du programme : Insémination Artificielle et FIVETE]
\begin{center}
\begin{tabularx}{\linewidth}{|l|X|X|}\hline
\rowcolor{popblueL} \textbf{Critère} & \textbf{Insémination Artificielle (IA / IAC / IAD)} & \textbf{FIVETE (Fécondation In Vitro et Transfert d'Embryon)} \\ \hline
\textbf{Indication principale} & Infertilité masculine modérée (oligo/asthénospermie), infertilité cervicale (glaires hostiles), infertilité inexpliquée, don de sperme. & Obstruction bilatérale des trompes (cause tubaire), infertilité masculine sévère (ICSI), échec IA, endométriose sévère, don d'ovocytes. \\ \hline
\textbf{Lieu de la fécondation} & \textbf{In vivo} (dans la trompe utérine). & \textbf{In vitro} (en laboratoire, boîte de Pétri / milieu de culture). \\ \hline
\textbf{Étapes clés} & 1. Stimulation ovarienne légère (optionnelle) + monitoring écho/hormonal. \\
2. Déclenchement ovulation (hCG). \\
3. Recueil sperme + \textbf{capacitation/préparation} (gradient densité, swim-up) $\rightarrow$ concentration spermatozoïdes mobiles. \\
4. \textbf{Dépôt intra-utérin} (cathéter) au moment de l'ovulation. & 1. \textbf{Stimulation ovarienne contrôlée} (gonadotrophines exogènes + agoniste/antagoniste GnRH) $\rightarrow$ multi-folliculaire. \\
2. \textbf{Ponction ovocytaire} (sous écho, anesthésie). \\
3. \textbf{Fécondation in vitro} : mise en contact ovocytes + spermatozoïdes (FIV classique) \textbf{OU} ICSI (microinjection spermatozoïde dans ovocyte). \\
4. \textbf{Culture embryonnaire} 2 à 5 jours (stade 4-8 cellules ou blastocyste). \\
5. \textbf{Transfert embryonnaire} (1 à 2 embryons) dans cavité utérine (cathéter souple). \\
6. Congélation embryons surnuméraires (vitrification). \\ \hline
\textbf{Taux de réussite / cycle} & $\approx$ 10--15 \% (proche fécondité naturelle mensuelle). & $\approx$ 20--30 \% (variable selon âge femme, cause, centre). \\ \hline
\textbf{Risques / Limites} & Grossesses multiples (si stimulation), infection rare. & Syndrome d'hyperstimulation ovarienne (SHSO), grossesses multiples, grossesse extra-utérine, coût élevé, charge psychologique. \\ \hline
\end{tabularx}
\end{center}
\end{propriete}

\begin{methode}[Présenter l'insémination artificielle et la FIVETE (Savoir-faire 6)]
\begin{enumerate}
\item \textbf{Partir de l'indication} : le choix de la technique dépend de la cause d'infertilité (défaut de dépôt/qualité spermatozoïdes + trompes perméables $\rightarrow$ insémination ; obstruction des trompes ou infertilité sévère $\rightarrow$ FIVETE).
\item \textbf{Décrire les étapes dans l'ordre chronologique}, en nommant précisément les structures (trompe, follicule, ovocyte, embryon, utérus).
\item \textbf{Distinguer les lieux de fécondation} : insémination artificielle = fécondation \emph{in vivo} (dans la trompe) ; FIVETE = fécondation \emph{in vitro} (en laboratoire) puis transfert d'embryon.
\item \textbf{Préciser les acteurs} : stimulation ovarienne, prélèvement, culture, transfert ; possibilité de don de sperme (IAD) ou d'ovocytes.
\item \textbf{Rester factuel} : citer les limites (taux de réussite de l'ordre de 20 à 30 \% par tentative de FIVETE) et le cadre éthique (respect de l'embryon, anonymat du don, non-commercialisation).
\end{enumerate}
\end{methode}

\begin{exoresolu}[Choisir et justifier une technique de PMA adaptée à la cause (Savoir-faire 6)]
\textbf{Couple 1 :} Monsieur présente une oligospermie modérée ($8 \times 10^6$ spz/mL, mobilité 30 \%) ; Madame a des trompes perméables (hystérosalpingographie normale) et ovule normalement (cycles réguliers, progestérone J21 $> 30$ nmol/L).
\textbf{Couple 2 :} Madame présente une obstruction bilatérale des trompes (séquelle de salpingite à Chlamydia, hystérosalpingographie : pas de passage du produit de contraste) ; les ovocytes (AMH normale) et le sperme (numération normale) sont de bonne qualité.
Pour chaque couple, proposer la technique de PMA la plus adaptée, décrire ses étapes et justifier le choix.
\tcblower
\textbf{Solution détaillée.}

\textbf{Couple 1 — Insémination Artificielle Intra-utérine avec sperme du Conjoint (IAC).}
\textbf{Étapes :} (1) Recueil du sperme par masturbation après 3--5 jours d'abstinence ; (2) \textbf{Préparation spermatique} en laboratoire (gradient de densité type PureSperm ou swim-up) pour sélectionner les spermatozoïdes les plus mobiles, les concentrer dans un faible volume (0,5 mL) et éliminer le plasma séminal (prostaglandines, bactéries) ; (3) Surveillance de l'ovulation (échographie + dosage LH urinaire ou sérique) $\pm$ légère stimulation (citrate de clomifène ou gonadotrophines faibles doses) ; (4) Déclenchement de l'ovulation par hCG (5000--10000 UI) quand follicule dominant $\ge 18$ mm ; (5) \textbf{Insémination} 36--40 h post-hCG : dépôt du sperme préparé dans la cavité utérine via cathéter souple ; la fécondation a lieu \emph{in vivo}, dans la trompe.
\textbf{Justification :} Les trompes sont perméables et l'ovulation normale ; le seul problème est la quantité/qualité des spermatozoïdes au niveau du col/utérus. La concentration en laboratoire et le dépôt intra-utérin raccourcissent le parcours et augmentent la densité de spermatozoïdes mobiles au site de fécondation. C'est la technique de 1\textsuperscript{er} intention.

\textbf{Couple 2 — FIVETE (avec ICSI si qualité sperme limite, sinon FIV classique).}
\textbf{Étapes :} (1) \textbf{Stimulation ovarienne contrôlée} : protocole antagoniste (court) ou agoniste (long) avec gonadotrophines recombinantes (FSH $\pm$ LH) pour obtenir 8--12 follicules matures ; monitoring écho/hormonal (E2) ; (2) Déclenchement ovulation (hCG ou agoniste GnRH) ; (3) \textbf{Ponction ovocytaire} transvaginale sous guidage échographique et anesthésie locale/générale ; (4) \textbf{Fécondation in vitro} : mise en contact ovocytes (dénudés des cellules cumulus) et spermatozoïdes préparés (50 000--100 000 spz/ovocyte) en milieu de culture (37°C, 5 \% CO2) \textbf{OU} ICSI (injection intracytoplasmique d'un spermatozoïde unique par ovocyte) si indication masculine ; (5) \textbf{Culture embryonnaire} : contrôle fécondation à J1 (2 pronucléus), division J2 (4 cellules), J3 (8 cellules) ou culture prolongée J5/J6 (blastocyste) ; (6) \textbf{Transfert d'embryon(s)} (1 ou 2 selon âge/qualité/réglementation) dans l'utérus sous contrôle échographique (vessie pleine) ; (7) Soutien lutéal (progestérone vaginale/IM) ; test $\beta$-hCG à J12--J14.
\textbf{Justification :} L'obstruction bilatérale des trompes empêche \textbf{toute rencontre naturelle} des gamètes (sperme dans utérus, ovocyte dans trompe) et rend l'insémination artificielle strictement inutile (le sperme déposé dans l'utérus ne peut pas traverser les trompes bouchées, l'ovocyte ne peut pas les rejoindre). La FIVETE court-circuite les trompes en réalisant la fécondation et le début du développement (2 à 5 jours) hors de l'organisme maternel, avant de déposer l'embryon directement dans la cavité utérine.

\textbf{Conclusion :} L'indication dépend du siège du dysfonctionnement : voie génitale franchissable (trompes perméables) + défaut spermatique modéré $\rightarrow$ insémination ; trompes obstruées ou infertilité sévère $\rightarrow$ FIVETE. Le taux de réussite de la FIVETE reste de l'ordre de 20 à 30 \% par cycle, fortement corrélé à l'âge maternel.
\end{exoresolu}

\begin{savaistu}[PMA au Cameroun : accès et éthique]
Au Cameroun, la PMA est encadrée par la loi n°2001/005 du 16 avril 2001 relative à la santé de la reproduction. Les centres autorisés (ex: CPMA de l'Hôpital Central de Yaoundé, Centre Mère-Enfant de la Fondation Chantal Biya, cliniques privées agréées) pratiquent l'IAC et la FIVETE/ICSI. Le don de gamètes est autorisé (anonyme, gratuit, non lucratif). La gestation pour autrui (GPA) est interdite. Le coût reste un frein majeur (FIVETE $\approx$ 1,5 à 3 millions FCFA non remboursés par la CNPS), limitant l'accès aux couples aisés. La sensibilisation à la prévention de l'infertilité (IST, avortements à risque) reste la priorité de santé publique.
\end{savaistu}

\coursec{Synthèse et intégration}

\begin{aretenir}[Bilan synthétique : Régulation hormonale et PMA]
\begin{itemize}
\item \textbf{Axe commun :} Hypothalamus (GnRH pulsatile) $\xrightarrow{+}$ Hypophyse (FSH, LH) $\xrightarrow{+}$ Gonades (Stéroïdes + Gamètes).
\item \textbf{Homme :} Régulation \textbf{linéaire stable} par rétrocontrôle \textbf{négatif} continu (Testostérone $-$ sur LH/GnRH ; Inhibine $-$ sur FSH). Pas de cyclicité.
\item \textbf{Femme :} Régulation \textbf{cyclique} (28 j) par alternance : Phase folliculaire (Rétrocontrôle $-$ puis \textbf{$+$} des \oe{}strogènes $\rightarrow$ Pic LH $\rightarrow$ Ovulation) ; Phase lutéale (Rétrocontrôle \textbf{$-$} fort Progestérone/\oe{}strogènes $\rightarrow$ Blocage FSH/LH $\rightarrow$ Régression corps jaune $\rightarrow$ Règles $\rightarrow$ Nouveau cycle).
\item \textbf{Infertilité :} Partage des causes (H/F/Mixte). Causes évitables majeures : IST, avortements non sécurisés, âge, toxiques.
\item \textbf{PMA :} \textbf{IA} = Fécondation \textbf{in vivo} (trompes perméables requises) ; \textbf{FIVETE} = Fécondation \textbf{in vitro} (trompes non requises, court-circuitage).
\end{itemize}
\end{aretenir}

\begin{exercice}[Entraînement Bac : Analyse expérimentale et schématisation]
\textbf{Exercice 1 (Régulation mâle).} On réalise une expérience sur le rat mâle adulte. Lot 1 : Témoin. Lot 2 : Hypophysectomisé. Lot 3 : Hypophysectomisé + injection quotidienne de LH. Lot 4 : Hypophysectomisé + injection quotidienne de FSH. Lot 5 : Hypophysectomisé + injection quotidienne de LH + FSH. On dose la testostérone sanguine et on observe la spermatogenèse après 3 semaines.
\begin{enumerate}
\item Prédire les résultats attendus pour chaque lot (testostérone : normale / basse / haute ; spermatogenèse : normale / absente / partielle).
\item Justifier le rôle respectif de la LH et de la FSH sur la fonction testiculaire.
\end{enumerate}

\textbf{Exercice 2 (Cycle féminin).} Une patiente de 28 ans, cycles réguliers de 28 jours, consulte pour infertilité. Le dosage hormonal au 3\textsuperscript{e} jour du cycle donne : FSH = 12 UI/L (N : 3--10), LH = 6 UI/L (N : 2--10), \oe{}stradiol = 45 pg/mL (N : 20--50). Au 21\textsuperscript{e} jour : Progestérone = 8 nmol/L (N phase lutéale $> 30$ nmol/L).
\begin{enumerate}
\item Interpréter le taux de FSH au J3. Quel diagnostic évoquer ?
\item Interpréter le taux de progestérone au J21. Quelle en est la conséquence sur l'endomètre et la nidation ?
\item Proposer une prise en charge thérapeutique adaptée.
\end{enumerate}

\textbf{Exercice 3 (PMA).} Un couple consulte après 2 ans d'infertilité. Bilan : Hystérosalpingographie = trompes bouchées bilatérales. Sperme = 2 millions/mL, mobilité 10 \%, formes normales 2 \% (oligo-asthéno-tératospermie sévère).
\begin{enumerate}
\item Quelle technique de PMA proposer en 1\textsuperscript{re} intention ? Justifier.
\item Décrire l'étape cruciale de la fécondation en laboratoire pour ce couple précis.
\item Citer deux risques majeurs de cette technique pour la femme et pour la grossesse.
\end{enumerate}
\end{exercice}

\begin{corrige}[Exercice 1]
\begin{enumerate}
\item Lot 1 : Testostérone normale, Spermatogenèse normale. Lot 2 : Testostérone très basse (pas de LH), Spermatogenèse absente. Lot 3 : Testostérone normale/restaurée (LH stimule Leydig), Spermatogenèse absente (pas de FSH sur Sertoli). Lot 4 : Testostérone basse (pas de LH), Spermatogenèse absente/partielle (FSH seul insuffisant sans testostérone locale haute). Lot 5 : Testostérone normale, Spermatogenèse normale/restaurée (synergie LH+FSH).
\item LH $\rightarrow$ Cellules de Leydig $\rightarrow$ Testostérone (effet endocrinien + paracrine haute concentration). FSH $\rightarrow$ Cellules de Sertoli $\rightarrow$ Facteurs de croissance (ABP, inhibine, nutriments) indispensables à la méiose et spermiogenèse. Les deux sont nécessaires : la testostérone initie/entretient, la FSH assure la complétion quantitative et qualitative.
\end{enumerate}
\end{corrige}

\begin{corrige}[Exercice 2]
\begin{enumerate}
\item FSH J3 = 12 UI/L $>$ 10 UI/L (seuil usuel). Signe d'\textbf{insuffisance ovarienne précoce} (réserve ovarienne diminuée) : le follicule ne sécrète pas assez d'inhibine/\oe{}stradiol pour freiner la FSH hypophysaire. Pronostic de fertilité altéré.
\item Progestérone J21 = 8 nmol/L $<$ 30 nmol/L. Signe d'\textbf{insuffisance lutéale} (corps jaune dysfonctionnel). Conséquence : endomètre mal sécrétoire, non réceptif $\rightarrow$ échec de nidation ou fausse couche précoce.
\item Stimulation ovarienne (citrate de clomifène ou gonadotrophines) pour améliorer recrutement folliculaire + \textbf{soutien lutéal} par progestérone naturelle (voie vaginale 200--400 mg/j ou IM) du J16 au test de grossesse. Si échec $\rightarrow$ FIVETE (meilleur contrôle, taux réussite supérieur malgré réserve diminuée).
\end{enumerate}
\end{corrige}

\begin{corrige}[Exercice 3]
\begin{enumerate}
\item \textbf{FIVETE avec ICSI (Injection Intra-Cytoplasmique de Spermatozoïde).} Justification : Double indication absolue. (1) Trompes bouchées $\rightarrow$ impossible de féconder in vivo (IA exclue). (2) Sperme très altéré (oligo-asthéno-tératospermie sévère) $\rightarrow$ FIV classique (mise en contact) taux de fécondation quasi nul ; ICSI permet de sélectionner \textbf{un} spermatozoïde morphologiquement normal et mobile par ovocyte.
\item Étape cruciale : \textbf{ICSI}. Sous microscope inversé (400x), immobilisation du spermatozoïde choisi (coupure queue), aspiration par micropipette d'injection, perforation de la zona pellucida et de l'oolème de l'ovocyte (métaphase II), injection du spermatozoïde dans le cytoplasme. Vérification de la survie de l'ovocyte et de la fécondation (2 pronucléus) à 16--18 h.
\item Risques : (Femme) \textbf{Syndrome d'Hyperstimulation Ovarienne (SHSO)} : ascite, hémoconcentration, risque thromboembolique, insuffisance rénale (prévention : agoniste GnRH pour déclenchement, congélation tous embryons). (Grossesse) \textbf{Grossesse multiple} (si transfert 2 embryons) : prématurité, faible poids, prééclampsie (politique de transfert électif d'un seul embryon - eSET - recommandée).
\end{enumerate}
\end{corrige}

\begin{attention}[Les 5 erreurs qui coûtent des points au Bac]
\begin{enumerate}
\item \textbf{Confondre les glandes productrices} : la testostérone et les \oe{}strogènes/progestérone sont sécrétées par les \emph{gonades}, jamais par l'hypophyse ; FSH et LH sont hypophysaires, la GnRH hypothalamique. Une inversion dans un tableau de synthèse est lourdement sanctionnée.
\item \textbf{Oublier de nommer l'hormone sur les flèches du schéma} et de préciser le signe ($+$ ou $-$) du rétrocontrôle : une flèche non légendée ne vaut rien. Le rétrocontrôle \emph{positif} des \oe{}strogènes (pic de LH) est l'exception à connaître et à schématiser.
\item \textbf{Dater l'ovulation sur le pic de progestérone} : c'est le \emph{pic de LH} (vers J14) qui déclenche l'ovulation ; la progestérone culmine une semaine plus tard (corps jaune, vers J21).
\item \textbf{Comparer des lots expérimentaux qui diffèrent par plusieurs variables} : toute interprétation valide exige une comparaison deux à deux avec une seule variable modifiée, et un lot témoin cité explicitement.
\item \textbf{Inverser insémination artificielle et FIVETE} : insémination = fécondation \emph{in vivo} (sperme déposé dans l'utérus) ; FIVETE = fécondation \emph{in vitro} puis \textbf{transfert d'embryon} (pas d'ovocyte, pas dans la trompe). Écrire « transfert d'ovocyte » ou « injection d'embryon dans la trompe » est une erreur de vocabulaire éliminatoire.
\end{enumerate}
\end{attention}

\newpage

\coursec{Exercices}

\courssub{Je vérifie mes connaissances}

\begin{exercice}[QCM : Identification des hormones sexuelles]
Pour chacune des affirmations suivantes, une seule réponse est correcte. Recopie le numéro et la lettre correspondante.
\begin{enumerate}
    \item L'hormone qui déclenche l'ovulation et transforme le follicule en corps jaune est :
    \begin{itemize}
        \item[A.] La FSH
        \item[B.] La LH
        \item[C.] L'œstradiol
        \item[D.] La progestérone
    \end{itemize}
    \item Chez l'homme, la testostérone est sécrétée par :
    \begin{itemize}
        \item[A.] Les cellules de Sertoli
        \item[B.] Les cellules de Leydig
        \item[C.] L'adénohypophyse
        \item[D.] L'hypothalamus
    \end{itemize}
    \item L'hormone qui inhibe la sécrétion de FSH par rétrocontrôle négatif chez l'homme est :
    \begin{itemize}
        \item[A.] La testostérone
        \item[B.] L'inhibine
        \item[C.] La LH
        \item[D.] La GnRH
    \end{itemize}
    \item Laquelle de ces hormones est une gonadulphine hypophysaire ?
    \begin{itemize}
        \item[A.] La progestérone
        \item[B.] L'œstradiol
        \item[C.] La LH
        \item[D.] La testostérone
    \end{itemize}
\end{enumerate}
\end{exercice}

\begin{exercice}[Vrai ou Faux justifié : Régulation hormonale]
Indique si chaque affirmation est VRAIE ou FAUSSE en justifiant brièvement ton choix par un mécanisme physiologique précis.
\begin{enumerate}
    \item La GnRH est sécrétée de façon continue et constante tout au long de la vie de la femme.
    \item Chez la femme, le pic de LH est déclenché par un rétrocontrôle positif des œstrogènes sur l'hypophyse.
    \item La progestérone sécrétée par le corps jaune inhibe la sécrétion de LH et de FSH pendant la phase lutéale.
    \item Chez l'homme, une destruction des cellules de Leydig entraîne une baisse de la testostérone et une augmentation de la FSH et de la LH.
\end{enumerate}
\end{exercice}

\begin{exercice}[Définitions et vocabulaire précis]
Donne une définition rigoureuse pour chacun des termes suivants en utilisant le vocabulaire scientifique attendu au Baccalauréat :
\begin{enumerate}
    \item \cle{Azoospermie}
    \item \cle{FIVETE} (Fécondation In Vitro et Transfert d'Embryon)
    \item \cle{Rétrocontrôle négatif} (appliqué à l'axe hypothalamo-hypophyso-gonal)
    \item \cle{Phase folliculaire} et \cle{Phase lutéale} (durée moyenne, structure sécrétrice principale, hormone dominante)
\end{enumerate}
\end{exercice}

\begin{exercice}[Lecture directe : Dosages hormonaux au jour 21 d'un cycle]
Le tableau ci-dessous donne les taux plasmatiques moyens de quatre hormones chez une femme au 21\ieme\ jour d'un cycle ovarien de 28 jours.
\begin{center}
\begin{tabularx}{\linewidth}{|l|X|}\hline
\textbf{Hormone} & \textbf{Taux plasmatique (UI/L ou nmol/L)} \\ \hline
FSH & 4,5 UI/L \\ \hline
LH & 6,2 UI/L \\ \hline
Œstradiol & 450 pmol/L \\ \hline
Progestérone & 45 nmol/L \\ \hline
\end{tabularx}
\end{center}
\begin{enumerate}
    \item Identifie la phase du cycle ovarien correspondant à ce jour. Justifie par deux arguments hormonaux.
    \item Quelle structure ovarienne est fonctionnelle à ce moment ? Quelle est son origine cellulaire ?
    \item Quel serait l'effet d'une administration exogène de progestérone à forte dose à cette période sur la sécrétion de LH ? Explique le mécanisme.
\end{enumerate}
\end{exercice}

\courssub{Je m'entraîne}

\begin{exercice}[Analyse de courbes hormonales : cycle ovarien de 28 jours]
On te fournit les courbes schématiques des variations des taux plasmatiques de FSH, LH, œstradiol et progestérone au cours d'un cycle ovarien standard de 28 jours (ovulation au jour 14).
\begin{enumerate}
    \item Associe chaque courbe (A, B, C, D) à l'hormone correspondante en justifiant par au moins un critère distinctif (pic, plateau, phase de sécrétion).
    \item Situe précisément le jour de l'ovulation sur l'axe des temps. Quel événement hormonal le déclenche ?
    \item Explique le rôle du \textbf{pic de LH} sur : (a) la rupture folliculaire, (b) la lutéinisation des cellules de la granulosa et de la thèque.
    \item Pourquoi le taux de FSH reste-t-il bas pendant la phase lutéale malgré la disparition du follicule dominant ?
\end{enumerate}
\end{exercice}

\begin{exercice}[Expérience classique : Rôle de l'hypophyse chez le rat mâle]
On réalise l'expérience suivante sur des rats mâles adultes :
\begin{itemize}
    \item \textbf{Lot 1 (Témoin) :} Rat intact.
    \item \textbf{Lot 2 :} Rat hypophysectomisé (ablation de l'hypophyse).
    \item \textbf{Lot 3 :} Rat hypophysectomisé + injection quotidienne d'extraits hypophysaires (contenant FSH et LH).
    \item \textbf{Lot 4 :} Rat hypophysectomisé + injection quotidienne de testostérone.
\end{itemize}
Au bout de 3 semaines, on mesure la masse des vésicules séminales et on observe la spermatogenèse au microscope.
\textbf{Résultats :}
\begin{center}
\begin{tabularx}{\linewidth}{|l|X|X|}\hline
\textbf{Lot} & \textbf{Masse vésicules séminales} & \textbf{Spermatogenèse} \\ \hline
1 (Témoin) & Normale & Complète (spermatozoïdes présents) \\ \hline
2 & Très atrophiée & Arrêtée au stade spermatogonie / spermatocyte I \\ \hline
3 & Normale & Complète \\ \hline
4 & Normale & Arrêtée au stade spermatogonie / spermatocyte I \\ \hline
\end{tabularx}
\end{center}
\begin{enumerate}
    \item Interprète les résultats du Lot 2 par rapport au Lot 1. Quel est le rôle de l'hypophyse sur l'appareil reproducteur mâle ?
    \item Compare les Lots 3 et 4. Pourquoi la testostérone seule (Lot 4) ne suffit-elle pas à restaurer une spermatogenèse complète ?
    \item Conclus sur les rôles respectifs de la LH (via la testostérone) et de la FSH dans la spermatogenèse.
\end{enumerate}
\end{exercice}

\begin{exercice}[Cas clinique : Infertilité masculine et rétrocontrôle]
Un couple consulte au Centre de Santé de la Maternité de l'Hôpital Central de Yaoundé pour infertilité de 3 ans. Le spermogramme de l'homme révèle une \textbf{azoospermie} (absence de spermatozoïdes dans l'éjaculat). Le dosage hormonal montre :
\begin{itemize}
    \item Testostérone : 2,5 nmol/L (Valeurs de référence : 10 -- 35 nmol/L)
    \item FSH : 45 UI/L (Valeurs de référence : 1,5 -- 12,4 UI/L)
    \item LH : 32 UI/L (Valeurs de référence : 1,7 -- 8,6 UI/L)
\end{itemize}
\begin{enumerate}
    \item Analyse les taux de FSH et LH par rapport aux normes. Quel type de rétrocontrôle est en jeu ici ?
    \item Déduis le niveau de l'anomalie : testiculaire, hypophysaire ou hypothalamique ? Justifie ton raisonnement en expliquant le mécanisme physiopathologique.
    \item Cette infertilité peut-elle être prise en charge par une simple injection de testostérone ? Pourquoi ?
    \item Propose une technique de PMA adaptée si une biopsie testiculaire retrouve des spermatozoïdes.
\end{enumerate}
\end{exercice}

\begin{exercice}[Savoir-faire : Élaboration d'un outil de sensibilisation]
Dans le cadre de la Journée Mondiale de la Santé de la Reproduction, tu dois concevoir un \textbf{dépliant (flyer) A5 recto-verso} destiné aux élèves de Première C et TI du Lycée de Bafoussam sur le thème : \og \textbf{Préserver sa fertilité dès l'adolescence : les gestes qui comptent} \fg.
\begin{enumerate}
    \item Structure ton dépliant en 4 parties logiques (ex : \og Savoir \fg, \og Risques évitables \fg, \og Signes d'alerte \fg, \og Où consulter ? \fg).
    \item Pour la partie \og Risques évitables \fg, liste 5 causes majeures d'infertilité \textbf{évitables} (infectieuses, toxiques, thermiques, comportementales, nutritionnelles) en précisant pour chacune le mécanisme lésif (ex : Chlamydia $\rightarrow$ salpingite $\rightarrow$ obstruction tubaire).
    \item Rédige un slogan accrocheur et une phrase d'accroche pour la couverture.
    \item Indique 2 structures de santé de référence au Cameroun où les jeunes peuvent consulter gratuitement et confidentiellement (ex : Centre de Santé Scolaire, CECOS, Association Camerounaise pour le Bien-Être Familial - ACBEF).
\end{enumerate}
\end{exercice}

\courssub{Je me prépare au Bac}

\begin{exercice}[Sujet type Bac : Analyse de document -- Régulation hormonale féminine et contraception]
\textbf{Document 1 :} Variations des taux plasmatiques d'hormones au cours d'un cycle ovarien de 28 jours chez une femme non contraceptée.
\begin{center}
\begin{tikzpicture}[scale=0.7]
\begin{axis}[
    width=\linewidth, height=8cm,
    xlabel={Jour du cycle}, ylabel={Taux relatif (arbitraire)},
    xmin=0, xmax=28, ymin=0, ymax=1.1,
    xtick={1,7,14,21,28}, ytick={0,0.2,0.4,0.6,0.8,1.0},
    legend pos=north east, grid=major,
    legend style={font=\small, fill opacity=0.8, draw opacity=1}
]
\addplot[popblue, thick, smooth, domain=0:28, samples=200] {0.3+0.4*exp(-(x-13)^2/8) + 0.1*exp(-(x-21)^2/18)}; \addlegendentry{Courbe 1}
\addplot[poporange, thick, smooth, domain=0:28, samples=200] {0.1+0.8*exp(-(x-13.5)^2/3)}; \addlegendentry{Courbe 2}
\addplot[popgreen, thick, smooth, domain=0:28, samples=200] {0.2+0.5*exp(-(x-12.5)^2/10) + 0.6*exp(-(x-21)^2/12)}; \addlegendentry{Courbe 3}
\addplot[poppurple, thick, smooth, domain=0:28, samples=200] {0.05+0.7*exp(-(x-21)^2/10)}; \addlegendentry{Courbe 4}
\draw[dashed, popdark] (14,0) -- (14,1.1) node[midway, right, font=\tiny] {Ovulation ?};
\end{axis}
\end{tikzpicture}
\end{center}
\textbf{Document 2 :} Mécanisme d'action d'une pilule contraceptive \og microprogestative \fg (contenant uniquement un progestatif de synthèse).
\begin{enumerate}
    \item \textbf{Identification :} Associe chaque courbe (1 à 4) aux hormones : FSH, LH, Œstradiol, Progestérone. Justifie chaque association par un argument tiré du document.
    \item \textbf{Ovulation :} Situe le jour probable de l'ovulation. Quel événement hormonal (visible sur le graphique) en est le déclencheur immédiat ?
    \item \textbf{Rétrocontrôles :} En t'appuyant sur les variations des courbes 1, 3 et 4, explique la succession des rétrocontrôles (négatif puis positif puis négatif) exercés par les hormones ovariennes sur l'hypophyse au cours du cycle.
    \item \textbf{Contraception :} En te basant sur le Document 2 et tes connaissances, explique par quels mécanismes la pilule microprogestative empêche la grossesse (cite au moins 3 mécanismes distincts : ovulation, glaire cervicale, endomètre, mobilité tubaire).
    \item \textbf{Contexte camerounais :} Au Cameroun, le taux de prévalence contraceptive moderne reste faible (\textasciitilde 20\% selon l'EDS-MICS 2018). Cite deux obstacles socio-culturels ou systémiques à l'accès à la contraception pour les adolescentes et propose une mesure concrète pour lever chaque obstacle.
\end{enumerate}
\end{exercice}

\begin{exercice}[Sujet type Bac : Cas clinique -- Choix d'une technique de PMA]
\textbf{Situation :} Mme M., 32 ans, et M. M., 35 ans, mariés depuis 5 ans, consultent au Service de Gynécologie-Obstétrique de l'Hôpital Général de Douala pour infertilité primaire.
\textbf{Bilan de la femme :} Cycles réguliers de 28 jours, température basale biphasique, dosage hormonal (J3) normal, hystérosalpingographie (HSG) montrant une \textbf{obstruction bilatérale des trompes} (hydrosalpinx) sans anomalie utérine.
\textbf{Bilan de l'homme :} Sperme normal (numération 60 millions/mL, mobilité 55\%, morphologie 15\% formes normales selon Kruger).
\textbf{Questionnaire :}
\begin{enumerate}
    \item Rappelle la définition de l'infertilité selon l'OMS. Pourquoi parle-t-on ici d'infertilité \og primaire \fg ?
    \item \textbf{Analyse étiologique :} Quelle est la cause de l'infertilité chez ce couple ? Pourquoi l'insémination artificielle (IA) \textbf{ne peut pas} être proposée en première intention ? Justifie par l'anatomie et la physiologie de la fécondation naturelle.
    \item \textbf{Choix de la technique :} La technique indiquée est la \cle{FIVETE}. Décris les \textbf{5 étapes principales} de cette technique en précisant pour chacune : le lieu (in vivo / in vitro), la durée approximative, et le contrôle biologique associé.
    \item \textbf{Stimulation ovarienne :} Avant la ponction, Mme M. reçoit un traitement par agoniste de la GnRH (protocol long) puis gonadotrophines (FSH recombinante). Explique l'intérêt de l'agoniste de la GnRH (effet \og flare-up \fg puis désensibilisation) pour obtenir une cohorte folliculaire synchronisée.
    \item \textbf{Aspects éthiques et légaux au Cameroun :} La loi camerounaise (Loi n°2005/015 du 29 décembre 2005) encadre la PMA. Cite deux principes éthiques fondamentaux qui régissent la FIVETE au Cameroun (ex : consentement libre et éclairé du couple, respect de l'embryon, interdiction du don de gamètes entre vifs ? -- \textit{vérifier la législation actuelle : don autorisé mais encadré}). Quelle est la limite d'âge souvent recommandée pour la prise en charge en FIVETE dans les centres publics camerounais ?
\end{enumerate}
\end{exercice}

\begin{exercice}[Sujet type Bac : Synthèse -- Régulation mâle vs femelle et applications thérapeutiques]
\textbf{Consigne :} Réponds aux questions suivantes en structurant ton raisonnement. La qualité de la rédaction, la précision du vocabulaire et la logique du raisonnement seront évaluées.
\begin{enumerate}
    \item \textbf{Schéma comparatif :} Réalise un tableau comparatif synthétique de la régulation hormonale chez l'homme et chez la femme (cycle) en 6 lignes : \og Hormone hypothalamique \fg, \og Gonadotrophines hypophysaires \fg, \og Cible gonadique \fg, \og Hormones gonadiques principales \fg, \og Type(s) de rétrocontrôle sur l'hypophyse \fg, \og Cyclicité \fg.
    \item \textbf{Mécanisme du rétrocontrôle positif :} Chez la femme, seul le rétrocontrôle \textbf{positif} des œstrogènes permet le pic de LH. Décris l'expérience classique (ex : ovariectomisée + implants d'œstradiol) qui a démontré la nécessité d'un \textbf{seuil critique} et d'une \textbf{durée minimale d'exposition} aux œstrogènes pour déclencher ce pic. Quels sont les récepteurs cellulaires et la zone hypothalamique impliqués ?
    \item \textbf{Application thérapeutique -- Induction de l'ovulation :} Une patiente présente une anovulation par insuffisance de la sécrétion de GnRH (hypogonadisme hypogonadotrope type Kallmann). On souhaite induire l'ovulation pour une FIVETE.
    \begin{itemize}
        \item[a.] Pourquoi l'administration de \textbf{clomifène} (anti-œstrogène) est-elle inefficace dans ce cas précis ?
        \item[b.] Propose le protocole de stimulation adapté (pompe à GnRH pulsatile ou injections de gonadotrophines). Explique le principe physiologique.
    \end{itemize}
    \item \textbf{Pathologie : Syndrome des Ovaires Polykystiques (SOPK).} C'est la première cause d'infertilité par anovulation au Cameroun (\textasciitilde 15-20\% des consultations). Le profil hormonal type montre : LH élevée, FSH normale ou basse, Testostérone élevée, AMH très élevée.
    \begin{itemize}
        \item[a.] Explique le mécanisme du cercle vicieux \og LH $\rightarrow$ Androgènes $\rightarrow$ Arrêt folliculaire $\rightarrow$ Pas de rétrocontrôle négatif $\rightarrow$ LH \fg.
        \item[b.] Le traitement de première intention pour induire l'ovulation est le \textbf{Letrozole} (inhibiteur de l'aromatase) ou le \textbf{Clomifène}. Explique le mécanisme d'action du Clomifène au niveau hypothalamo-hypophysaire pour rétablir un pic de LH.
    \end{itemize}
\end{enumerate}
\end{exercice}

\newpage

\coursec{Corrigés des exercices}

\begin{corrige}[Exercice 1 — QCM : Identification des hormones sexuelles]
\begin{enumerate}
    \item \textbf{Réponse B — La LH.} Le pic pré-ovulatoire de LH (vers le 14\ieme\ jour) déclenche la rupture du follicule de De Graaf (ovulation) et provoque la transformation des cellules de la granulosa et de la thèque interne en cellules lutéales : c'est la \cle{lutéinisation}, qui donne naissance au corps jaune. La FSH (A) agit surtout sur la croissance folliculaire ; l'œstradiol (C) et la progestérone (D) sont des hormones ovariennes, non hypophysaires.
    \item \textbf{Réponse B — Les cellules de Leydig.} Situées dans le tissu interstitiel entre les tubes séminifères, les cellules de Leydig sécrètent la testostérone sous le contrôle de la LH. Les cellules de Sertoli (A), situées dans les tubes séminifères, sont la cible de la FSH et sécrètent l'inhibine.
    \item \textbf{Réponse B — L'inhibine.} Sécrétée par les cellules de Sertoli, l'inhibine exerce un rétrocontrôle négatif sélectif sur la sécrétion de FSH par l'adénohypophyse. La testostérone (A) inhibe principalement la LH (et la GnRH).
    \item \textbf{Réponse C — La LH.} Les gonadotrophines hypophysaires sont la FSH et la LH, sécrétées par l'adénohypophyse. La progestérone, l'œstradiol et la testostérone sont des stéroïdes gonadiques.
\end{enumerate}
\end{corrige}

\begin{corrige}[Exercice 2 — Vrai ou Faux justifié : Régulation hormonale]
\begin{enumerate}
    \item \textbf{FAUX.} La GnRH est sécrétée de façon \cle{pulsatile} (une décharge toutes les 60 à 90 minutes environ) et non continue. C'est d'ailleurs cette pulsatilité qui est indispensable : une administration continue de GnRH (ou d'un agoniste) provoque une désensibilisation de l'hypophyse et l'arrêt des sécrétions de FSH et LH (principe utilisé en PMA).
    \item \textbf{VRAI.} En fin de phase folliculaire, le follicule dominant sécrète des taux élevés d'œstradiol (supérieurs à un seuil critique d'environ 200 pg/mL maintenu pendant au moins 48 h). Ce taux élevé et prolongé exerce un \cle{rétrocontrôle positif} sur l'axe hypothalamo-hypophysaire, déclenchant le pic de LH (et le petit pic de FSH) ovulatoire.
    \item \textbf{VRAI.} Pendant la phase lutéale, la progestérone (associée à l'œstradiol) exerce un \cle{rétrocontrôle négatif} sur l'hypothalamus et l'hypophyse, maintenant les taux de FSH et de LH bas. C'est ce blocage qui empêche le recrutement d'un nouveau follicule dominant pendant cette phase.
    \item \textbf{VRAI.} Les cellules de Leydig sécrètent la testostérone. Leur destruction fait chuter la testostérone ; le rétrocontrôle négatif de la testostérone sur l'axe hypothalamo-hypophysaire disparaît, donc la sécrétion de LH (et de FSH, par levée partielle du frein) augmente : c'est un \cle{hypogonadisme hypergonadotrope} (gonadotrophines élevées, stéroïdes bas), signe d'une insuffisance gonadique primaire.
\end{enumerate}
\end{corrige}

\begin{corrige}[Exercice 3 — Définitions et vocabulaire précis]
\begin{enumerate}
    \item \cle{Azoospermie} : absence totale de spermatozoïdes dans l'éjaculat, constatée sur au moins deux spermogrammes après centrifugation. Elle peut être \emph{obstructive} (spermatogenèse normale mais obstacle sur les voies excrétrices) ou \emph{non obstructive} (défaut de production testiculaire).
    \item \cle{FIVETE} : Fécondation In Vitro et Transfert d'Embryon. Technique de procréation médicalement assistée consistant à mettre en contact, \emph{in vitro} (en laboratoire, hors du corps), des ovocytes prélevés par ponction folliculaire et des spermatozoïdes, puis à transférer dans l'utérus un ou plusieurs embryons obtenus après quelques jours de culture.
    \item \cle{Rétrocontrôle négatif} : mécanisme de régulation par lequel une hormone gonadique (testostérone, œstradiol, progestérone, inhibine) inhibe la sécrétion des hormones qui ont stimulé sa propre production (GnRH hypothalamique et/ou FSH et LH hypophysaires), maintenant ainsi les taux hormonaux dans des limites compatibles avec l'équilibre physiologique.
    \item \cle{Phase folliculaire} : première moitié du cycle ovarien (en moyenne du 1\ier\ au 14\ieme\ jour, durée variable) ; la structure sécrétrice principale est le follicule en croissance puis le follicule de De Graaf ; l'hormone dominante est l'\textbf{œstradiol}. \cle{Phase lutéale} : seconde moitié du cycle (du 14\ieme\ au 28\ieme\ jour, durée quasi constante d'environ 14 jours) ; la structure sécrétrice est le \textbf{corps jaune} ; l'hormone dominante est la \textbf{progestérone}.
\end{enumerate}
\end{corrige}

\begin{corrige}[Exercice 4 — Lecture directe : Dosages hormonaux au jour 21]
\begin{enumerate}
    \item Le jour 21 d'un cycle de 28 jours correspond à la \cle{phase lutéale} (milieu de phase lutéale, environ 7 jours après l'ovulation). Deux arguments hormonaux :
    \begin{itemize}
        \item la \textbf{progestérone est élevée} (45 nmol/L, valeur caractéristique du pic lutéal, alors qu'elle est inférieure à 3 nmol/L en phase folliculaire) ;
        \item la \textbf{FSH et la LH sont basses} (4,5 et 6,2 UI/L), témoins du rétrocontrôle négatif exercé par la progestérone et l'œstradiol sur l'hypophyse. L'œstradiol présente en outre son second pic (450 pmol/L), typique de la sécrétion lutéale.
    \end{itemize}
    \item La structure ovarienne fonctionnelle est le \cle{corps jaune}. Il dérive de la transformation (lutéinisation) des cellules de la \textbf{granulosa} et de la \textbf{thèque interne} du follicule rompu, sous l'action du pic de LH.
    \item Une administration exogène de progestérone à forte dose renforcerait le \cle{rétrocontrôle négatif} sur l'hypothalamus (ralentissement des pulses de GnRH) et sur l'adénohypophyse : la sécrétion de \textbf{LH chuterait encore davantage}. Conséquence : le corps jaune, privé de son soutien lutéotrophique (LH), régresserait prématurément, et aucun nouveau follicule ne pourrait être recruté. C'est d'ailleurs le principe des contraceptifs progestatifs.
\end{enumerate}
\end{corrige}

\begin{corrige}[Exercice 5 — Analyse de courbes hormonales : cycle ovarien de 28 jours]
\begin{enumerate}
    \item Identification des courbes (critères distinctifs) :
    \begin{itemize}
        \item \textbf{Courbe présentant un pic unique, très élevé et bref au jour 14 : LH.} C'est le pic ovulatoire, le plus ample et le plus court du cycle.
        \item \textbf{Courbe à deux pics (un large pic pré-ovulatoire vers J13 et un second pic modéré en phase lutéale) : œstradiol.} Sécrété par le follicule puis par le corps jaune.
        \item \textbf{Courbe basse en phase folliculaire puis formant un large plateau en phase lutéale (J15--J25) : progestérone.} Hormone quasi absente avant l'ovulation, dominante après.
        \item \textbf{Courbe avec un petit pic en début de cycle et un pic modéré synchronisé avec celui de la LH : FSH.} Le pic initial recrute la cohorte folliculaire ; le pic de J14 accompagne celui de la LH.
    \end{itemize}
    \item L'ovulation se situe au \textbf{jour 14} (28 $-$ 14, la phase lutéale durant 14 jours). Elle est déclenchée par le \cle{pic de LH} (l'ovulation survient environ 24 à 36 h après le début du pic).
    \item Rôle du pic de LH :
    \begin{itemize}
        \item[(a)] \textbf{Rupture folliculaire :} la LH active des enzymes protéolytiques (collagénases, plasmine) qui digèrent la paroi du follicule de De Graaf au niveau du stigma, libérant l'ovocyte II bloqué en métaphase II.
        \item[(b)] \textbf{Lutéinisation :} la LH transforme les cellules de la granulosa et de la thèque interne en cellules lutéales, vascularisées et sécrétrices de progestérone : c'est la formation du corps jaune.
    \end{itemize}
    \item Pendant la phase lutéale, la \textbf{progestérone} (et l'œstradiol) sécrétées par le corps jaune exercent un \cle{rétrocontrôle négatif} puissant sur l'hypothalamus et l'hypophyse, maintenant la FSH (et la LH) à un taux bas malgré l'absence de follicule dominant fonctionnel. Ce n'est qu'à la chute du corps jaune (fin de cycle, en l'absence de fécondation) que ce frein est levé et que la FSH remonte pour recruter une nouvelle cohorte.
\end{enumerate}
\end{corrige}

\begin{corrige}[Exercice 6 — Expérience classique : Rôle de l'hypophyse chez le rat mâle]
\begin{enumerate}
    \item \textbf{Lot 2 vs Lot 1 :} après hypophysectomie, les vésicules séminales s'atrophient fortement et la spermatogenèse s'arrête au stade spermatogonie/spermatocyte I. L'hypophyse est donc \textbf{indispensable} au maintien de la fonction testiculaire : elle contrôle à la fois la sécrétion de testostérone (dont dépend la croissance des glandes annexes comme les vésicules séminales) et le déroulement complet de la spermatogenèse, via ses gonadotrophines FSH et LH.
    \item \textbf{Comparaison Lots 3 et 4 :} le Lot 3 (extraits hypophysaires = FSH + LH) restaure à la fois la masse des vésicules séminales et une spermatogenèse complète. Le Lot 4 (testostérone seule) restaure la masse des vésicules séminales (effet androgénique périphérique) mais \emph{pas} la spermatogenèse. La testostérone seule ne suffit donc pas : il manque un facteur hypophysaire agissant directement sur les tubes séminifères. De plus, la testostérone injectée exerce un rétrocontrôle négatif qui maintient les gonadotrophines endogènes à zéro, et la concentration \emph{intratesticulaire} de testostérone reste insuffisante sans stimulation des cellules de Leydig par la LH.
    \item \textbf{Conclusion :} la \textbf{LH} stimule les cellules de Leydig qui sécrètent la testostérone ; celle-ci assure le développement et le trophisme des organes annexes (vésicules séminales, prostate) et est nécessaire, à forte concentration locale, à la spermatogenèse. La \textbf{FSH} agit sur les cellules de Sertoli, qu'elle rend sensibles à la testostérone (synthèse de l'ABP, \emph{androgen binding protein}, qui concentre la testostérone dans le tube séminifère) et qui nourrissent et accompagnent les cellules germinales. La spermatogenèse complète exige donc l'action \textbf{conjointe} de la FSH et de la testostérone.
\end{enumerate}
\end{corrige}

\begin{corrige}[Exercice 7 — Cas clinique : Infertilité masculine et rétrocontrôle]
\begin{enumerate}
    \item Les taux de FSH (45 UI/L, norme 1,5--12,4) et de LH (32 UI/L, norme 1,7--8,6) sont \textbf{nettement augmentés} (environ 4 fois la norme), alors que la testostérone est \textbf{effondrée} (2,5 nmol/L, norme 10--35). L'axe hypothalamo-hypophysaire fonctionne et s'active fortement : c'est la levée du \cle{rétrocontrôle négatif} normalement exercé par la testostérone (et l'inhibine) sur l'hypophyse et l'hypothalamus. L'hypophyse, non freinée, hypersécrète FSH et LH.
    \item L'anomalie est \textbf{testiculaire} (insuffisance gonadique primaire, encore appelée \cle{hypogonadisme hypergonadotrope}). Raisonnement : si la lésion était hypophysaire ou hypothalamique, les gonadotrophines seraient \emph{basses} (hypogonadisme hypogonadotrope). Or ici FSH et LH sont élevées : l'hypophyse répond correctement à l'absence de frein, c'est donc le testicule qui ne répond pas à la stimulation (cellules de Leydig et lignée germinale défaillantes : azoospermie non obstructive).
    \item \textbf{Non.} Injecter de la testostérone corrigerait les signes d'hypoandrogénie mais : (a) elle ne restaure pas la spermatogenèse, qui exige la FSH et une concentration intratesticulaire élevée de testostérone ; (b) par rétrocontrôle négatif, elle ferait chuter FSH et LH, aggravant encore l'arrêt de la spermatogenèse. Ici la lésion testiculaire est de toute façon la cause première : le testicule ne peut pas répondre.
    \item Si la biopsie testiculaire retrouve des spermatozoïdes (même rares), la technique adaptée est la \cle{FIVETE avec ICSI} (injection intracytoplasmique d'un spermatozoïde) : prélèvement chirurgical de spermatozoïdes testiculaires ou épididymaires (TESE), puis injection d'un spermatozoïde directement dans le cytoplasme de l'ovocyte in vitro, suivie du transfert d'embryon.
\end{enumerate}
\end{corrige}

\begin{corrige}[Exercice 8 — Savoir-faire : Élaboration d'un outil de sensibilisation]
\begin{enumerate}
    \item \textbf{Structure du dépliant en 4 parties :}
    \begin{itemize}
        \item \textbf{Recto 1 — « Savoir » :} rappel simple : la fertilité dépend de la santé des testicules, des trompes et de l'équilibre hormonal ; elle se construit dès l'adolescence.
        \item \textbf{Recto 2 — « Risques évitables » :} les 5 causes majeures (voir question 2).
        \item \textbf{Verso 1 — « Signes d'alerte » :} douleurs pelviennes persistantes, écoulements génitaux, cycles irréguliers ou absents, douleur ou gonflement testiculaire, antécédent d'IST non traitée $\Rightarrow$ consulter sans délai.
        \item \textbf{Verso 2 — « Où consulter ? » :} adresses et gratuité/confidentialité (voir question 4).
    \end{itemize}
    \item \textbf{Cinq causes évitables et leur mécanisme lésif :}
    \begin{itemize}
        \item \emph{Infectieuse :} infection à \emph{Chlamydia trachomatis} non traitée $\rightarrow$ salpingite (infection des trompes) $\rightarrow$ cicatrisation et obstruction tubaire $\rightarrow$ impossibilité de rencontre des gamètes. Chez l'homme : orchi-épididymite $\rightarrow$ obstruction des voies excrétrices.
        \item \emph{Toxique :} tabac et alcool $\rightarrow$ stress oxydatif et toxicité directe sur les cellules germinales $\rightarrow$ baisse de la numération et de la mobilité des spermatozoïdes, altération de la réserve ovarienne.
        \item \emph{Thermique :} exposition répétée des testicules à la chaleur (saunas, slips serrés, ordinateur sur les genoux, varicocèle) $\rightarrow$ élévation de la température scrotale au-dessus de 34--35 \degre C $\rightarrow$ arrêt ou ralentissement de la spermatogenèse (qui exige une température inférieure d'environ 2 \degre C à celle du corps).
        \item \emph{Comportementale :} rapports sexuels non protégés et multi-partenariat $\rightarrow$ IST (Chlamydia, gonococcie) $\rightarrow$ séquelles tubaires ou épididymaires ; avortements clandestins $\rightarrow$ synéchies utérines et infections.
        \item \emph{Nutritionnelle :} dénutrition sévère ou obésité $\rightarrow$ perturbation de la pulsatilité de la GnRH $\rightarrow$ anovulation, aménorrhée ; carences $\rightarrow$ altération de la qualité des gamètes.
    \end{itemize}
    \item \textbf{Slogan :} « Protège ta fertilité aujourd'hui, pour choisir ta famille demain ! » \textbf{Phrase d'accroche :} « Une infection non traitée à 17 ans peut empêcher une grossesse à 27 ans : informe-toi, dépiste-toi, protège-toi. »
    \item \textbf{Structures de référence :} les centres de santé scolaires et universitaires (consultations confidentielles) ; l'Association Camerounaise pour le Bien-Être Familial (CAMNAFAW), qui propose des services « jeunes » gratuits ou à coût réduit ; on peut ajouter les centres de santé intégrés et les services de gynécologie des hôpitaux de district.
\end{enumerate}
\end{corrige}

\begin{corrige}[Exercice 9 — Sujet type Bac : Régulation hormonale féminine et contraception]
\textbf{Barème indicatif (sur 20) :} Q1 : 6 pts ; Q2 : 2 pts ; Q3 : 5 pts ; Q4 : 4 pts ; Q5 : 3 pts.
\begin{enumerate}
    \item \textbf{Identification des courbes :}
    \begin{itemize}
        \item \textbf{Courbe 2 = LH :} pic unique, très ample (0,9) et très bref, centré juste avant le jour 14 — c'est le pic ovulatoire.
        \item \textbf{Courbe 1 = FSH :} pic modéré synchronisé avec celui de la LH vers J13, avec un niveau de base plus élevé en début de cycle (recrutement folliculaire).
        \item \textbf{Courbe 3 = Œstradiol :} deux pics — un pic pré-ovulatoire (vers J12--13) et un second pic plus large en phase lutéale (vers J21).
        \item \textbf{Courbe 4 = progestérone :} quasi nulle en phase folliculaire, large plateau maximal vers J21 (phase lutéale), chute en fin de cycle.
    \end{itemize}
    \item \textbf{Ovulation :} jour 14 (ligne pointillée), déclenchée par le \textbf{pic de LH} (courbe 2), lui-même provoqué par le pic pré-ovulatoire d'œstradiol (courbe 3).
    \item \textbf{Succession des rétrocontrôles :}
    \begin{itemize}
        \item \emph{Début de phase folliculaire (J1--J10) :} l'œstradiol, à taux modéré, exerce un \textbf{rétrocontrôle négatif} sur l'hypophyse : FSH et LH restent basses pendant que le follicule dominant croît.
        \item \emph{Fin de phase folliculaire (J12--J13) :} l'œstradiol dépasse un seuil critique pendant plus de 48 h $\rightarrow$ \textbf{rétrocontrôle positif} : déclenchement brutal du pic de LH (et de FSH) $\rightarrow$ ovulation.
        \item \emph{Phase lutéale (J15--J28) :} la progestérone (courbe 4), associée à l'œstradiol, exerce à nouveau un \textbf{rétrocontrôle négatif} : FSH et LH retombent et restent basses jusqu'à la régression du corps jaune.
    \end{itemize}
    \item \textbf{Mécanismes de la pilule microprogestative (3 attendus) :}
    \begin{itemize}
        \item blocage de l'\textbf{ovulation} par rétrocontrôle négatif continu du progestatif sur l'axe hypothalamo-hypophysaire (suppression du pic de LH) ;
        \item épaississement de la \textbf{glaire cervicale}, devenant imperméable au passage des spermatozoïdes ;
        \item atrophie de l'\textbf{endomètre}, le rendant impropre à la nidation ;
        \item (bonus) ralentissement de la \textbf{mobilité tubaire}, perturbant le transport des gamètes.
    \end{itemize}
    \item \textbf{Obstacles et mesures (exemples) :}
    \begin{itemize}
        \item \emph{Obstacle socio-culturel :} pression familiale/religieuse et idées reçues (« la pilule rend stérile ») $\rightarrow$ \emph{mesure :} campagnes d'éducation à la santé reproductive en milieu scolaire et communautaire, impliquant les chefs traditionnels et religieux.
        \item \emph{Obstacle systémique :} coût et ruptures de stock, éloignement des centres de santé, exigence de l'accord parental $\rightarrow$ \emph{mesure :} gratuité des contraceptifs pour les adolescentes dans les centres de santé publics et services « conviviaux pour jeunes » garantissant la confidentialité.
    \end{itemize}
\end{enumerate}
\textbf{Points de vigilance :} ne pas confondre les pics de FSH et de LH (celui de la LH est toujours le plus ample) ; citer explicitement le \emph{seuil} et la \emph{durée} d'exposition aux œstrogènes pour le rétrocontrôle positif ; pour la contraception, distinguer clairement les mécanismes centraux (blocage ovulatoire) et périphériques (glaire, endomètre).
\end{corrige}

\begin{corrige}[Exercice 10 — Sujet type Bac : Choix d'une technique de PMA]
\textbf{Barème indicatif (sur 20) :} Q1 : 3 pts ; Q2 : 4 pts ; Q3 : 6 pts ; Q4 : 4 pts ; Q5 : 3 pts.
\begin{enumerate}
    \item Selon l'OMS, l'\cle{infertilité} est l'absence de grossesse après \textbf{12 à 24 mois} de rapports sexuels réguliers non protégés. Elle est dite \textbf{primaire} car le couple n'a jamais obtenu aucune grossesse (par opposition à l'infertilité secondaire, survenant après une grossesse antérieure).
    \item \textbf{Cause :} l'obstruction bilatérale des trompes (hydrosalpinx) empêche la rencontre des gamètes. Or la fécondation naturelle a lieu dans le \textbf{tiers externe de la trompe} (ampoule tubaire) : l'ovocyte II, capté par les franges du pavillon, y rencontre les spermatozoïdes remontés depuis le vagin. L'\textbf{insémination artificielle} dépose les spermatozoïdes dans l'utérus, mais ceux-ci ne peuvent pas franchir l'obstacle tubaire pour atteindre l'ovocyte : l'IA exige au moins \textbf{une trompe perméable}. Elle est donc contre-indiquée ici.
    \item \textbf{Les 5 étapes de la FIVETE :}
    \begin{enumerate}
        \item \textbf{Stimulation ovarienne contrôlée} (in vivo, environ 10--14 jours) : injections de gonadotrophines pour obtenir plusieurs follicules matures ; contrôle par échographie et dosage d'œstradiol.
        \item \textbf{Ponction ovocytaire} (prélèvement in vivo, ovocytes cultivés in vitro ; quelques minutes sous anesthésie, environ 36 h après le déclenchement par hCG) : aspiration des ovocytes par voie vaginale sous échographie.
        \item \textbf{Fécondation in vitro} (in vitro, 16--20 h) : mise en contact des ovocytes et des spermatozoïdes préparés en milieu de culture ; contrôle de la fécondation (2 pronuclei) le lendemain.
        \item \textbf{Culture embryonnaire} (in vitro, 2 à 5 jours) : développement des embryons jusqu'au stade 4--8 cellules ou blastocyste ; contrôle morphologique quotidien.
        \item \textbf{Transfert embryonnaire} (in vivo, quelques minutes) : dépôt d'un ou deux embryons dans la cavité utérine à l'aide d'un cathéter ; contrôle par dosage de $\beta$-hCG 12--14 jours plus tard.
    \end{enumerate}
    \item \textbf{Intérêt de l'agoniste de la GnRH (protocole long) :} administré en continu, l'agoniste provoque d'abord une stimulation brutale de l'hypophyse (effet \emph{flare-up} : libération massive de FSH/LH), puis, par désensibilisation des récepteurs hypophysaires (la GnRH naturelle étant pulsatile, une stimulation continue épuise la réponse), un \textbf{blocage complet des sécrétions endogènes de FSH et LH}. On obtient ainsi une hypophyse « au repos » : le pic de LH spontané (qui provoquerait une ovulation prématurée avant la ponction) est empêché, et l'administration exogène de FSH recombinante permet de recruter une \textbf{cohorte folliculaire homogène et synchronisée}, dont la maturation est entièrement contrôlée par le médecin.
    \item \textbf{Éthique et législation :} deux principes fondamentaux : le \textbf{consentement libre et éclairé} du couple (information sur les techniques, les risques et les chances de succès, recueilli par écrit) et le \textbf{respect de l'embryon} (encadrement du nombre d'embryons transférés et de leur conservation). La PMA est réservée à un couple hétérosexuel vivant, en âge de procréer ; dans les centres publics camerounais, la prise en charge en FIVETE est généralement limitée à environ \textbf{42--45 ans} pour la femme (au-delà, les chances de succès deviennent très faibles).
\end{enumerate}
\textbf{Points de vigilance :} pour la Q2, la justification anatomique (lieu de la fécondation = trompe) est exigée, pas seulement « les trompes sont bouchées » ; pour la Q3, préciser systématiquement le lieu (in vivo / in vitro) ; ne pas écrire que l'agoniste « stimule » durablement l'ovulation — c'est sa désensibilisation qui est recherchée.
\end{corrige}

\begin{corrige}[Exercice 11 — Sujet type Bac : Synthèse -- Régulation mâle vs femelle]
\textbf{Barème indicatif (sur 20) :} Q1 : 5 pts ; Q2 : 5 pts ; Q3 : 4 pts ; Q4 : 6 pts.
\begin{enumerate}
    \item \textbf{Tableau comparatif :}
\begin{center}
\begin{tabularx}{\linewidth}{|l|X|X|}\hline
\rowcolor{popblueL} \textbf{Critère} & \textbf{Homme} & \textbf{Femme} \\ \hline
Hormone hypothalamique & GnRH (pulsatile, rythme régulier) & GnRH (pulsatile, fréquence variable selon la phase du cycle) \\ \hline
Gonadotrophines hypophysaires & FSH et LH (sécrétion continue, peu fluctuante) & FSH et LH (sécrétion cyclique, pic ovulatoire) \\ \hline
Cible gonadique & Cellules de Leydig (LH) ; cellules de Sertoli (FSH) & Cellules de la thèque (LH) ; cellules de la granulosa (FSH) \\ \hline
Hormones gonadiques principales & Testostérone ; inhibine & Œstradiol ; progestérone ; inhibine \\ \hline
Rétrocontrôle(s) sur l'hypophyse & Uniquement négatif (testostérone sur LH ; inhibine sur FSH) & Négatif (œstradiol à taux modéré, progestérone) ET positif (œstradiol à taux élevé prolongé) \\ \hline
Cyclicité & Pas de cycle : production continue de spermatozoïdes de la puberté à la vieillesse & Cycle ovarien d'environ 28 jours, de la puberté à la ménopause \\ \hline
\end{tabularx}
\end{center}
    \item \textbf{Mécanisme du rétrocontrôle positif — expérience classique :} chez une femelle (singe ou brebis) \textbf{ovariectomisée}, les taux de FSH et LH s'élèvent fortement (levée du rétrocontrôle négatif). L'implantation d'une capsule d'œstradiol délivrant un taux \emph{faible} fait retomber FSH et LH : rétrocontrôle négatif. Si l'on élève ensuite la dose d'œstradiol pour atteindre un \textbf{seuil critique} (environ 200 pg/mL) maintenu pendant une \textbf{durée minimale} (environ 48 h), on observe, après un délai, une décharge massive de LH reproduisant le pic ovulatoire : c'est le rétrocontrôle positif. Si la dose est sous le seuil ou l'exposition trop brève, aucun pic n'apparaît. L'œstradiol agit via des \textbf{récepteurs nucléaires ER$\alpha$} exprimés par les neurones à GnRH (zone \textbf{préoptique} de l'hypothalamus, impliquant les neurones à kisspeptine) et par les cellules gonadotropes de l'adénohypophyse, dont il augmente la sensibilité à la GnRH.
    \item \textbf{Induction de l'ovulation (hypogonadisme hypogonadotrope type Kallmann) :}
    \begin{itemize}
        \item[a.] Le \textbf{clomifène} est un anti-œstrogène qui agit en bloquant les récepteurs aux œstrogènes au niveau hypothalamo-hypophysaire : l'axe « croit » manquer d'œstrogènes et répond en \emph{augmentant la sécrétion de GnRH}, donc de FSH/LH. Or dans le syndrome de Kallmann, les neurones à GnRH sont absents ou non fonctionnels : lever le frein ne peut rien produire, le clomifène est donc \textbf{inefficace}.
        \item[b.] Protocole adapté : soit la \textbf{pompe à GnRH pulsatile} (administration sous-cutanée d'une dose toutes les 90 minutes environ), qui imite la pulsatilité physiologique et restaure une sécrétion normale de FSH et LH ; soit les \textbf{injections de gonadotrophines} (FSH recombinante puis hCG pour déclencher l'ovulation), qui court-circuitent l'étage hypothalamo-hypophysaire défaillant en stimulant directement l'ovaire.
    \end{itemize}
    \item \textbf{Syndrome des Ovaires Polykystiques :}
    \begin{itemize}
        \item[a.] \textbf{Cercle vicieux :} la LH élevée stimule excessivement les cellules de la \textbf{thèque}, qui produisent trop d'\textbf{androgènes} (testostérone). L'hyperandrogénie bloque la maturation folliculaire (\textbf{arrêt de croissance des follicules} au stade 5--8 mm : ovaires polykystiques, anovulation). En l'absence de follicule dominant puis de corps jaune, il n'y a \textbf{ni pic d'œstradiol prolongé ni progestérone} : le rétrocontrôle négatif ne s'exerce pas, la LH reste élevée, ce qui entretient l'hyperandrogénie — et le cercle se referme.
        \item[b.] \textbf{Mécanisme du clomifène :} le clomifène bloque les récepteurs aux œstrogènes au niveau de l'hypothalamus et de l'hypophyse. L'axe hypothalamo-hypophysaire, privé du signal de rétrocontrôle négatif, « perçoit » un faux déficit en œstrogènes et augmente la sécrétion pulsatile de GnRH, donc de \textbf{FSH} (et modérément de LH). Cette relance de FSH permet le recrutement et la maturation d'un follicule dominant ; l'œstradiol produit atteint alors le seuil du rétrocontrôle positif, déclenchant un \textbf{pic de LH} ovulatoire. (Le Letrozole, lui, bloque l'aromatase, diminuant la synthèse d'œstradiol, ce qui relance la FSH par le même type de levée du rétrocontrôle.)
    \end{itemize}
\end{enumerate}
\textbf{Points de vigilance :} dans le tableau, exiger la mention des deux types de rétrocontrôle chez la femme (c'est la différence fondamentale avec l'homme) ; pour le Kallmann, le raisonnement « le clomifène agit en amont d'un étage défaillant » doit être explicité ; dans le SOPK, ne pas écrire que la FSH est élevée — c'est le rapport LH/FSH qui est augmenté.
\end{corrige}

\newpage

\coursec{Fiche bilan}

\begin{popfigure}
\centering
\begin{tikzpicture}[
    mindmap,
    grow cyclic,
    every node/.style={concept, execute at begin node=\hskip0pt},
    root concept/.append style={concept color=popblue, font=\large\bfseries, text width=3.5cm, minimum size=3.5cm},
    level 1 concept/.append style={level distance=5cm, sibling angle=51.4, font=\normalsize\bfseries, text width=2.8cm, minimum size=2.8cm},
    level 2 concept/.append style={level distance=3.5cm, sibling angle=45, font=\small, text width=2.2cm, minimum size=2.2cm},
    concept color=popblue!70
]
\node[root concept, align=center]{Santé\\ Reproductive}
    child[concept color=poporange] { node[align=center]{Régulation\\ Homme}
        child { node {Axe HPG} }
        child { node {Testostérone} }
        child { node {Rétroaction -} }
    }
    child[concept color=popgreen] { node[align=center]{Régulation\\ Femme}
        child { node {Cycle ovarien} }
        child { node {Œstrogènes/Progest.} }
        child { node {Rétroaction +/-} }
    }
    child[concept color=poppurple] { node {Infertilité}
        child { node {Causes mâles} }
        child { node {Causes femelles} }
        child { node {Mixtes/Idiop.} }
    }
    child[concept color=poppink] { node {PMA}
        child { node[align=center]{Insémination\\ artificielle} }
        child { node {FIVETE} }
        child { node {Éthique/Loi} }
    }
    child[concept color=popgold] { node[align=center]{Hormones\\ Clés}
        child { node {GnRH, LH, FSH} }
        child { node {Testostérone} }
        child { node {Œstradiol, Progest.} }
    }
    child[concept color=popdark] { node[align=center]{Prévention\\ \& Dépistage}
        child { node {IST/VIH} }
        child { node {Hygiène vie} }
        child { node {Consultation} }
    };
\end{tikzpicture}
\legende{Carte mentale récapitulative de la leçon 7 : Santé reproductive.}
\end{popfigure}

\begin{bilan}
\textbf{Définitions clés}
\begin{itemize}
    \item \cle{Hormones sexuelles} : Substances messagères (stéroïdes ou peptidiques) sécrétées par les gonades (testicules, ovaires) et le placenta, sous contrôle de l'axe hypothalamo-hypophysaire. Rôles : gamétogenèse, caractères sexuels secondaires, métabolisme, comportement.
    \item \cle{Spermatogenèse} : Formation continue des spermatozoïdes dans les tubules séminifères (durée $\approx$ 74 jours chez l'homme), stimulée par la FSH et la testostérone.
    \item \cle{Oogenèse} : Formation discontinue des ovocytes (arrêt en prophase I à la naissance, reprise à la puberté, un ovocyte par cycle), sous contrôle de la FSH et de la LH.
    \item \cle{Cycle ovarien} : Série cyclique d'événements (phase folliculaire, ovulation, phase lutéale) d'une durée moyenne de 28 jours, pilotée par les variations d'œstradiol et de progestérone.
    \item \cle{Infertilité} : Absence de grossesse après 12 à 24 mois de rapports sexuels réguliers non protégés. Distincte de la stérilité (incapacité absolue).
    \item \cle{PMA (Procréation Médicalement Assistée)} : Ensemble de techniques cliniques et biologiques permettant la conception in vivo (insémination artificielle) ou in vitro (FIVETE).
\end{itemize}

\textbf{Mécanismes régulateurs à connaître par cœur}
\begin{center}
\begin{tabularx}{\linewidth}{|>{\bfseries}l|X|X|}\hline
\rowcolor{popblueL}
\textbf{Axe / Phase} & \textbf{Hormones \& Cibles} & \textbf{Rétroactions \& Effets} \\ \hline
\cle{Axe HPG Mâle} & GnRH (Hypothalamus) $\rightarrow$ LH/FSH (Adénohypophyse) $\rightarrow$ Cellules de Leydig (LH) $\rightarrow$ Testostérone ; Cellules de Sertoli (FSH) $\rightarrow$ ABP, Inhibine. & \textbf{Rétroaction négative} : Testostérone $\searrow$ GnRH/LH ; Inhibine $\searrow$ FSH. Maintien taux constants. \\ \hline
\cle{Phase Folliculaire (Femme)} & FSH $\rightarrow$ Croissance follicule $\rightarrow$ Sécrétion \textbf{Œstradiol} (croissant). & \textbf{Rétroaction négative} (début) : Œstradiol $\searrow$ FSH (sélection follicule dominant). \\ \hline
\cle{Ovulation (J14)} & Pic d'Œstradiol $\rightarrow$ \textbf{Rétroaction positive} sur LH (et FSH) $\rightarrow$ \textbf{Pic de LH} $\rightarrow$ Rupture follicule, ovulation, lutéinisation. & Seuil critique Œstradiol $>$ 200 pg/mL pendant $\approx$ 48h. \\ \hline
\cle{Phase Lutéale (Femme)} & Corps jaune $\rightarrow$ \textbf{Progestérone} (forte) + Œstradiol. & \textbf{Rétroaction négative forte} : Progestérone $\searrow$ GnRH/LH/FSH $\rightarrow$ Atrésie corps jaune (si pas grossesse) $\rightarrow$ Règles. \\ \hline
\cle{Grossesse} & Trophoblaste $\rightarrow$ \textbf{hCG} (maintient corps jaune) $\rightarrow$ Progestérone/Œstradiol $\rightarrow$ Maintien endomètre, blocage cycles. & hCG détectable urines/sang (test grossesse). Prend le relais vers 10-12 SA (placenta). \\ \hline
\end{tabularx}
\end{center}

\textbf{Méthodes express}
\begin{enumerate}
    \item \textbf{Analyser un profil hormonal (femme) :} Repérer le pic de LH $\rightarrow$ Ovulation J+1. Taux progestérone J21 $>$ 10 ng/mL $\rightarrow$ Ovulation confirmée. Taux FSH J3 $>$ 10 UI/L $\rightarrow$ Réserve ovarienne diminuée.
    \item \textbf{Identifier cause infertilité :} Spermo-gramme (numération, mobilité, morphologie) $\rightarrow$ Mâle. Courbes de température / Dosages hormonaux / Hystérosalpingographie $\rightarrow$ Femelle. Bilan caryotype / sérologies $\rightarrow$ Mixtes.
    \item \textbf{Choisir technique PMA :} Infertilité mâle modérée / stérilité cervicale / inexpliquée $\rightarrow$ Insémination artificielle (IAC/IAD). Trompes bouchées / infertilité mâle sévère / échecs IA $\rightarrow$ FIVETE (FIV + Transfert Embryonnaire). Azoospermie obstructive $\rightarrow$ TESE/ICSI.
\end{enumerate}
\end{bilan}

\begin{aretenir}[Checklist avant le Bac]
\begin{itemize}
    \item $\square$ Citer les 5 hormones clés de l'axe HPG (GnRH, LH, FSH, Testostérone, Œstradiol, Progestérone) et leur origine glandulaire.
    \item $\square$ Expliquer la rétroaction négative de la testostérone sur LH/FSH chez l'homme (schéma à flèches +/-).
    \item $\square$ Différencier rétroaction négative (phase folliculaire début, phase lutéale) et positive (pic LH) chez la femme.
    \item $\square$ Relier le pic de LH à l'ovulation (J+1) et à la formation du corps jaune.
    \item $\square$ Décrire les effets de la progestérone sur l'endomètre (phase sécrétoire) et sur l'axe HPG (blocage).
    \item $\square$ Lister 3 causes d'infertilité masculine (oligo/asthénotératozoospermie, azoospermie, varicocèle, IST séquellaires) et 3 causes féminines (anovulation, trompes bouchées, endométriose, âge).
    \item $\square$ Différencier Insémination Artificielle (fécondation in vivo) et FIVETE (fécondation in vitro + transfert).
    \item $\square$ Citer les indications de l'ICSI (injection intracytoplasmique de spermatozoïde) : infertilité mâle sévère.
    \item $\square$ Rappeler le cadre légal/éthique au Cameroun (Loi n°2001/016 : don de gamètes anonyme, interdiction clonage, gestation pour autrui encadrée/interdite).
    \item $\square$ Interpréter un spermo-gramme selon normes OMS 2021 (concentration $\ge$ 16 M/mL, mobilité progressive $\ge$ 30\%, formes normales $\ge$ 4\%).
    \item $\square$ Construire un outil de sensibilisation (affiche, dépliant) sur la prévention de l'infertilité (dépistage IST, hygiène de vie, âge procréer).
\end{itemize}
\end{aretenir}

\findecours

\end{document}
